%!TEX TS-program = pdflatex
%!TEX encoding = UTF-8 Unicode

\documentclass[10.5pt,twoside]{scrbook}


\setparsizes{10pt}{0pt}{0pt plus 1fil}

\usepackage{geometry}
\geometry{paperwidth = 5in,paperheight=7in, textwidth=4in,textheight=5.5in}

%\usepackage[urw-garamond]{mathdesign}
%\usepackage{kpfonts}
%\usepackage{libertine}
\usepackage[T1]{fontenc}
\usepackage{amsmath}
\usepackage{libertine}
\usepackage[libertine]{newtxmath}
%\usepackage[nomath]{libertinus}

%
\usepackage{cabin}

\renewcommand*\ttdefault{lcmtt}

\usepackage[protrusion=allmath,tracking=small caps]{microtype}


\usepackage{physics}


\setcounter{secnumdepth}{0}


\usepackage[usenames,svgnames]{xcolor}

%\usepackage{amsmath}

\usepackage{graphicx}
%\usepackage{quattrocento}
\usepackage[per-mode=symbol]{siunitx}

\usepackage{nicefrac}

\usepackage{soul}

\newcommand{\hlred}[1]{\textcolor{Red}{#1}}% prints in red
\newcommand{\hlblue}[1]{\textcolor{SteelBlue}{\textbf{#1}}}% prints in blue
\newcommand{\hlgray}[1]{\textcolor{black!70}{{#1}}}% prints in gray 

\usepackage{mathrsfs}

\usepackage{scrlayer-scrpage}

\newcommand{\grotsq}{%
   \fontfamily{ugq}\selectfont}
\pagestyle{scrheadings}   
\setkomafont{pageheadfoot}{\grotsq\small}
\setkomafont{pagenumber}{\grotsq\small}

\usepackage{tocloft}

\setlength\cftparskip{-2pt}
\setlength\cftbeforesecskip{3pt}
\setlength\cftaftertoctitleskip{3pt}
 

% length before after figure
\setlength{\textfloatsep}{6pt plus 0.75pt minus 1.5pt}
\setlength{\intextsep}{10pt plus 1.0pt minus 2.0pt}

\usepackage{titlesec}

\titleformat{\chapter}
  {\fontfamily{ugq}\selectfont\huge\color{black!80}}
  {\thechapter}{1em}{}
\titleformat{\section}
  {\fontfamily{ugq}\selectfont\Large\color{black!70}}
  {\thesection}{1em}{}
%\setlength{\parskip}{0.25\baselineskip}
%\setlength{\parfillskip}{0pt plus 1fil}
%\setparsizes{10pt}{0pt}{0pt plus 1fil}
%\setparsizes{1em}{0.1\baselineskip plus .1\baselineskip}{1em plus 1fil}
%\usepackage{pstricks}

\renewcommand{\baselinestretch}{.9}

% for no dots in toc
\renewcommand{\cftdot}{} % no dot leaders
\renewcommand{\cftchapfont}{\fontfamily{ugq}\selectfont\large}
\renewcommand{\cftsecfont}{\fontfamily{ugq}\selectfont}
\renewcommand{\cfttoctitlefont}{\fontfamily{ugq}\selectfont\huge}
\renewcommand{\cftchappagefont}{\fontfamily{ugq}\selectfont\large}
\renewcommand{\cftsecpagefont}{\fontfamily{ugq}\selectfont}


\usepackage{caption}
\usepackage[font={color=black!70,footnotesize},figurename={Figure}]{caption}

\usepackage{enumitem}
% \setlength{\parskip}{1ex plus 0.5ex minus 0.2ex}
\usepackage[ colorlinks=true]{hyperref}
\hypersetup{
  colorlinks,
  linkcolor=black!70,
  linktoc=all
}

\titlespacing*{\section}
{0pt}{2.5ex plus 1ex minus .2ex}{2.3ex plus .2ex}

\usepackage{emptypage}


\usepackage{hyphenat}

%\usepackage{etoolbox}


\newcommand*\figr[1]{\hyperref[#1]{Figure~\ref{#1}}}

\usepackage{dialogue}

\usepackage{lips}

\DeclareSIUnit\kgf{kgf}
\DeclareSIUnit\amus{amu}
\DeclareSIUnit\calorie{cal}
\DeclareSIUnit\mercury{Hg}
\DeclareSIUnit\atmos{atm}
\DeclareSIUnit\gf{gf}
\DeclareSIUnit\dyne{dyn}
\DeclareSIUnit\erg{erg}
\DeclareSIUnit\gauss{G}
\DeclareSIUnit\fahrenheit{^{\circ}F}

\usepackage{pgfornament}

\newcommand{\TODO}{\textcolor{Red}{\textbf TODO!}\xspace}
\newcommand{\ie}{\textit{i.\hairsp{}e.}\xspace}
\newcommand{\eg}{\textit{e.\hairsp{}g.}\xspace}

%\DeclareMathAlphabet{\pazocal}{OMS}{zplm}{m}{n}
%\newcommand{\Ea}{\mathcal{E}}

\usepackage{url}
\renewcommand{\UrlFont}{\color{RoyalBlue}\small\texttt}

\usepackage{booktabs}

%\usepackage[version=4, arrows=pgf{To[length=3pt]}{.15ex}]{mhchem}
\usepackage[version=4]{mhchem}
%https://tex.stackexchange.com/questions/557753/using-mhchem-with-gfsneohellenic-gives-a-weird-looking-arrow

\usepackage{chemfig}
%\setatomsep{2.5em}


\usepackage{lipsum}

\usepackage{pdfpages}

\usepackage{float}

\title{Physics for Everyone: Book 1 Physical Bodies}
\author{L. D. Landau, A. I. Kitaigordsky}
\date{2020}

\begin{document}

\hyphenation{one-tho-us-and-th}

\frontmatter

%
% !TEX root = pfe-book1.tex
\thispagestyle{empty}


\begin{tikzpicture}[remember picture,overlay]
  \fill[gray!10]
        (current page.north west) rectangle (current page.south east);
\end{tikzpicture}

\begin{tikzpicture}[remember picture,overlay] 
\fill [fill=Maroon,yshift=-3cm] (current page.north west) rectangle (\textwidth,4);
\end{tikzpicture}


\begin{tikzpicture}[remember picture,overlay] 
%\fill [fill=IndianRed,yshift=-3cm] (current page.north west) rectangle (\textwidth,1);


\node [scale=2,font=\sffamily,darkgray,yshift=3.cm]	 at (current page.center) {L. D. Landau \quad A. I. Kitaigorodsky};
\end{tikzpicture}

%
\begin{tikzpicture}[remember picture,overlay] 
%\node [yshift=-3.1cm]	 at (current page.center) {\includegraphics[width=0.9\textwidth]{figs/front-back/tarasov-qasp-fc3.pdf}};
\node [scale=1.5,yshift=-5cm,font=\sffamily,darkgray]	 at (current page.center) {MIR PUBLISHERS MOSCOW};
\end{tikzpicture}



\cleardoublepage
\includepdf{pfe-book1-fc.pdf}
%todo front and back cover in pure tikz
\cleardoublepage
% !TEX root = pfe-book4.tex
%!TEX TS-program = pdflatex
%!TEX encoding = UTF-8 Unicode


\cleardoublepage
\thispagestyle{empty}
%\maketitle
\vspace*{1cm}
\begin{center}
\fontfamily{ugq}\selectfont{\Huge Physics for Everyone} \\[10pt] 
%\textsf{\Huge Physics for Everyone} \\[10pt] 
%\textsf{\huge Book 4: Photons and Nuclei} \\[10pt]
\fontfamily{ugq}\selectfont{\LARGE Book 4: Photons and Nuclei} \\[10pt]
\hrule
\vspace{10pt}

{\Large A. I. Kitaigordsky}\\[15pt]
\begin{flushright}
Translated from the Russian \\by Martin Greendlinger,\\
D.Sc.(Math.)\\[15pt]
Typeset in \LaTeX \, by\\ Damitr Mazanov
\end{flushright}
\vspace{4cm}
{\Large Mir Publishers Moscow}
\end{center}
\cleardoublepage
%normalfont
%\chapter{Back Cover}
% Physics for Everyone
% Book 1L.DLondou
% A, I. Kitoigorodsky
% PHYSOiL
% BODES
% Translated from
% the Russian
% by Martin Greendlinger,
% D.Sc.(Math.)
% .
% Mir Publishers Moscow

%\clearpage
\vfill
\vspace*{1cm}
This electronic version released in 2020 by 

\url{http://mirtitles.org}

{\small The Mir Titles project is an attempt to conserve the knowledge in the form of various books that came out during the Soviet era for the future generations.}

The project files for this series are available at:

\url{https://gitlab.com/mirtitles/pfe}

This edition published in 2022 by 

\includegraphics[width=0.5\textwidth]{figures/tantralipi-logo.pdf}

\url{http://tantralipi.com}


\thispagestyle{empty}
\clearpage

% !TEX root = pfe-book4.tex
%!TEX TS-program = pdflatex
%!TEX encoding = UTF-8 Unicode


\chapter{Preface} 
%\pagestyle{mystyle}
%\headrulewidth{0pt}
%\lhead{\nouppercase{\rightmark}}
%\rhead{\nouppercase{\leftmark}}
%\renewcommand{\chaptermark}[1]{%
%\markright{#1}{}}
This is the fourth and concluding book in the \emph{Physics for Everyone} series and it deals with the fundamentals of physics.

``Fundamentals'' is of course a rather vague word but we will think of it as meaning the general laws on which the whole edifice of modern physics rests. There are not so many of them, and so we can make a list: the laws of motion of classical mechanics, the laws of thermody­namics, the laws that are embodied in the equations of Maxwell and that govern charges, currents and electro­magnetic fields, and then the laws of quantum physics and the theory of relativity.

The laws of physics, like those of natural science at large, are of an empirical nature. They are arrived at by means of observation and experiment. Experiments establish a multitude of primary facts such as the build­ing up of matter from atoms and molecules, the nuclear model of the atom, the wave-particle aspect of matter, and so on. Now, both the number of basic laws and also the number of fundamental facts and concepts necessary for their description is not so very great: At any rate, it is limited.

During the past several decades, physics has grown and expanded to such an extent that workers in different branches cease to understand one another as soon as the discussion goes beyond what holds them together in one family, that is, beyond the limits of the laws and concepts underlying all branches of physics. Portions of physics are closely interwoven with technology, with other areas of natural science, with medicine, and even with be humanitarian sciences. It is easy to see why they
have set themselves up as independent disciplines. 

Surely no one would argue that any discussion of the various divisions of applied physics must be preceded by an examination of the basic laws and facts. And it is just as true that different writers select and arrange the material needed for laying the foundation of physics each in his own way, depending on his individual tastes and his own special field of inquiry what I have to offer here is merely one of many possible expositions of the fundamentals of physics.

The type of reader envisaged by this \emph{Physics for Every­one} series has been mentioned in the prefaces to the ear­ lier books. I will repeat that this series is aimed at repre­sentatives of all professions who wish to recall the phys­ics they studied, get a picture of the modern state of the science, and evaluate the effect it has on scientific and technological progress and on forming a materialist world outlook. Many pages of these books will, I am sure, be of interest to teachers of physics and to students at school that have come to like physics. And finally there may be something of interest for those readers who are depressed by even a simple algebraic equation.

Quite naturally, this series is not intended to take the place of a textbook. \emph{Photons and Nuclei} is an attempt, on the part of the author, to demonstrate to the reader how the laws of the electromagnetic field and quantum physics operate when we consider the behaviour of electromagnetic waves of different wavelength. Before taking up atomic nuclei, the reader will get some idea of what wave mechanics and the special theory of relativity are about. This is followed by a discussion of the basic facts concerning the structure of atomic nuclei, and then the topic will be sources of energy on the earth -- a topic of burning interest to humanity at large. We conclude our brief talk with a story about the physics of the universe.

The limited scope of this book has forced us to give up many traditional topics. The old must always give way to the new.

\begin{flushright}
 \emph{%April 1978\\
A. I. Kitaigorodsky}
\end{flushright}
\newpage

%\pagestyle{mystyle}


%{\pagestyle{plain}
%\markboth{Contents}{Contents} 
\tableofcontents
%\pagestyle{fancy}
\cleardoublepage

\mainmatter

% !TEX root = pfe-book4.tex
%!TEX TS-program = pdflatex
%!TEX encoding = UTF-8 Unicode


\cleardoublepage
%\mainmatter
\chapter{Soft Electromagnetic Radiation}
\label{ch-01}

\section{Exchange of Energy by Radiation}
\emph{Soft electromagnetic radiation} is that with wavelengths lying roughly in the interval from 0.1 to 100 micrometres. Also, bear in mind that when we speak of soft radiation we mean electromagnetic waves not dealt with in radio engineering. This stipulation is important because purely radio-engineering methods permit one to dip into the region of soft radiation. The soft radiation also rather frequently goes by the simple term ``light''. When applying that term, one must bear in mind that visible light occupies only a very narrow section of wavelengths -- for the ``average'' human eye it lies between 380 and 780 nanometres (or from 0.38 to 0.78 micrometres).

In future, whenever we have the occasion to speak of ``light'', we will do so in the broad sense of the word because the laws that hold true for the visible portion of the spectrum remain true for all other representatives of soft radiation.

Also note that radiation with shorter wavelengths than visible light is called \emph{ultraviolet radiation}; the longer wavelengths are termed \emph{infrared radiation}.

We can now turn to the topic of our discussion. It will be recalled that there are three modes of heat transfer. They are called heat conduction, thermal convection,
and thermal radiation. In order to study the exchange of energy that occurs in thermal radiation, we will have to  examine the behaviour of bodies in a vacuum (where convection is impossible) separated by a certain distance (this is to exclude the conduction of heat).

Experiments have shown that if two or more bodies form a closed system (the reader will recall that this means the absence of any exchange of energy between objects not in the system), the temperatures of the bodies equal out. Each one of the bodies of the system is at the same time a radiator and an absorber. What occur are numberless transitions of atoms and molecules from a higher level to a lower level (such events involve the emission of photons) and from a lower level to a higher level (photons are absorbed). Photons of all energies (or, what is the same thing, electromagnetic waves of all wavelengths) participate in these exchanges.

Quite naturally, the body does not absorb all the energy falling on it. There may be bodies that scatter more energy or transmit certain wavelengths. But this is of no consequence because sooner or later a thermal equilibrium is established nevertheless.

The condition of thermal equilibrium requires that the ratio of the energy of absorption to the energy of emission be the same for all wavelengths. This theorem was rigorously demonstrated in 1860 by the German physicist Gustav Robert Kirchhoff (1824-1887). The ratio can change for different temperatures, but if the temperature is fixed, then it is the same for photons of all energies.

This is a clear enough theorem and hardly needs any demonstration of proof. The idea behind the law is that the number of absorbed photons of a given kind (that is, of a definite energy) is equal, in the case of thermal equilibrium, to the number of radiated photons of that particular kind. From this we get the following rule: if an object is a strong absorber of any kind of rays, then those same rays are just as strongly radiated.

This rule helps to predict the conditions under which thermal equilibrium sets in. Why is water in a bottle with silvered sides so slow to heat up under the action of the sun's rays, whereas water in a flask made of black glass heats up very quickly? The explanation is obvious: a black body absorbs rays intensively and their energy goes to increase the temperature, and thermal equilibrium sets in after intense heating. A silvered surface, on the contrary, is an excellent reflector. Only a small amount of the energy is absorbed, it takes a long time to heat the body, and equilibrium sets in at a low temperature.

Now let’s reverse the experiment. Pour some hot water into both flasks and put them into a refrigerator. Which one will cool off quickest? The one that heats up faster will cool off faster. If more energy is absorbed, more is released.

Some very effective experiments can be performed with coloured ceramics. If the object is green, the piece absorbs all colours except green. This is because the eye sees those rays that are reflected (or scattered) by the material. Now heat up the fragment. How will it appear? The answer is right at the tip of your tongue: violet because violet is the complementary colour of yellow-green. Complementary colours are those that produce white if they are mixed. Newton was the one who introduced the term ``complementary colour'' when he decomposed light rays into a spectrum with the aid of a glass prism.

\section{The Radiation of Incandescent Bodies}
It is well known that a piece of metal, when heated, first becomes red and then white. Most chemical substances cannot be thus heated. They either melt or decompose. Therefore, what follows refers mostly to metals.

The most remarkable thing is that the radiation spectrum of all heated bodies is not at all specific. The point is this. From the basic law about energy levels it is clear that the radiation spectrum and the absorption spectrum of a body must coincide. Metals are opaque throughout the region of the spectrum of soft radiation. From this it follows that they must also radiate photons of all energies.
\label{bbr-ref}
Let's put this differently: a continuous spectrum appears due to the fact that in a multi-atomic system the energy levels of the atoms merge into overlapping bands. In such a system, all energy transitions are possible, that is, we can find any energy difference between the $m^{\textrm{th}}$ and $n^{\textrm{th}}$ levels $E_{m} - E_{n}$ and, hence, any frequencies of radiation and absorption. \figr{fig-1.1} shows the spectrum of an incandescent body for several temperatures (we give here the theoretical curves that hold true for a so-called ideal black body).
\begin{figure}[!ht]
\centering
\includegraphics[width=0.9\textwidth]{figures/fig-01-01.pdf}
\caption{Theoretical spectrum of an incandescent (black) body for several temperatures.}
\label{fig-1.1}
\end{figure}


It is well to point out here that the derivation of the shape of this curve (this was done by Max Planck in 1900) was the first step in the development of quantum physics. In order to obtain agreement of theory with experiment, Planck had to assume that radiation and absorption of light take place in separate portions. Planck was not prepared to take the next step and state that we are justified in speaking of particles of light (photons). That step was taken by Albert Einstein in 1905.

It was only in 1913 that Niels Bohr introduced the concept of quantization of energy. And if we want a logically rigorous theory of thermal radiation, the year of birth can be put as 1926.

Let's first discuss the shapes of the curves and only then talk about theory. First of all, note that as the temperature rises, the area under the curve rapidly increases. What is the physical meaning of the area bounded by the radiation curve? When constructing curves similar to those depicted in the figure, we say that the intensity of radiation for a given wavelength is laid off on the axis of ordinates (horizontal axis). But what does a ``given wavelength'' mean? Do we mean 453 or 453.2 nanometres? Maybe \num{453.257859987654}? It is probably clear that when we speak of a ``given wavelength'', we mean a very small interval of wavelengths. We agree, say, that the interval is equal to 0.01 nanometre. From this it follows that it is not the ordinate that has physical meaning but a tiny column with a base of 0.01 nanometre. The area of this column is equal to the energy radiated by the waves having lengths in that interval (for example, from 453.25 to 453.26 nanometres). Now if we break up the whole area into such columns, we get the total intensity of the whole spectrum. That is precisely the operation mathematicians perform and it is called integration. To summarize: the area under our curve yields the total intensity of the radiation, and it turns out to be proportional to the fourth power of the temperature.


In the figure we are discussing it is clear that with increasing temperature there is not only a change in the area occupied by the curve but there is a shift of its maximum to the left, that is, into the region of ultraviolet radiation.

The relationship between the wavelength of light in micrometres that corresponds to the greatest intensity of radiation (or absorption) and the temperature in kelvins is given by the following formula:
\begin{equation*}%
\lambda_{\textrm{max}} = \frac{2886}{T}
\label{lambdamax}
\end{equation*}
At the lowest temperatures, the maximum lies in the infrared region. That is precisely why infrared radiation is also termed thermal radiation. And it is a marvelous thing that we have instruments capable of sensing the thermal radiation emitted by bodies at room temperature and even lower. There are instruments today that can see in total darkness. Some animals, by the way, have the same capability. There is nothing at all strange in this fact since infrared rays have, in principle, the same properties as visible rays.

Also, don't forget that every animal is a source of radiation. We sometimes hear of a person being able to ``feel'' the presence of another person in darkness. No mysticism is involved, merely the one who ``feels'' has a highly sensitive perception of thermal rays.


%\newpage
\begin{center}
\includegraphics[width=0.8\textwidth]{figures/planck.pdf}
\end{center}
{\small \textsf{{Max Planck [1858-1947]}} -- \textsf{\footnotesize outstanding German scientist who laid the foundations of quantum theory. In an attempt to find a mathematical expression for a proper description of the spectral distribution of the emission of an ideal black body Planck demonstrated that such a formula could be obtained by introducing a ``quantum of action''. Planck assumed that a body emits energy in parcels, equal to the product of a constant (which later was named after him) by the frequency of the light.}}



I can't resist telling the reader about an interesting episode that demonstrates the necessity of taking thermal rays into account even when, in the very ordinary sense of this word, the source of rays is not a heated body. A few years ago I was asked to investigate some experiments conducted by a person who considered himself a magician capable of stopping a motor using only his will power.

My task was to find a rational explanation (sorcerers of the twentieth like to deal in pseudoscientific terminology and so call these experiments telekinesis).

\begin{figure}[!ht]
\centering
\includegraphics[width=\textwidth]{figures/fig-01-02.pdf}
\caption{Deciphering the sorcerer's trick.}
\label{fig-1.2}
\end{figure}

A diagram of the experiment is shown in \figr{fig-1.2}. A wing was set in rotation by a small motor, and the wing really did stop whenever the magician sat down next to the box where the axle of the motor emerged from below. I soon found that anybody who sat down in that position would be able to stop the wing. It usually took about 10 to 15 minutes to perform the operation. And it wasn’t the motor that came to a halt, as the magician claimed, but the little wing. It was clear then that some kind of force connected with the human body was interfering with the force of adhesion between the axle of the motor and the wing.

I pointed out that the wing could be stopped almost instantly if an electric lamp were brought up close to the side of the box. It was obvious that the trick lay in the heat emitted by the human body. I sent a whiff of tobacco smoke into the box and demonstrated that the convection currents of air inside the box move so as to prevent the wing from rotating. Exact measurements showed that the temperature of the side of the box closest to the human body is about one degree higher than the opposite side. Infrared rays emitted by a body heated to 60-70 degrees Celsius can be felt if the hand is held at a short distance. Thermal convection of course must be eliminated. Heated air rises upwards, so bring your hand close from below. You will then be certain that you are perceiving thermal rays.

We conclude our talk about thermal rays with an explanation of why the modern electric light bulb with a tungsten filament is a great step beyond a bulb with a carbon filament. The whole point is that a carbon filament can be heated to \SI{2100}{\kelvin}, while a tungsten filament can be heated all the way up to \SI{2500}{\kelvin}. Why are these 400 degrees so important? Because the purpose of an incandescent lamp is to provide light and not heat, and so the aim is to have the maximum of the curve located in the visible portion of the radiation. It will
be seen from the graph that the ideal is to have a filament that can withstand the temperature of the Sun's surface, or \SI{6000}{\kelvin}. But even the step from 2100 degrees to 2500 degrees raises the portion of energy involving visible radiation from 0.5\% to 1.6\%.

\section{The Theory of Thermal Radiation}
If a system of radiating and absorbing bodies is closed, then the photon ``gas'' (with the aid of which the bodies exchange energy) must be in equilibrium with the atoms supplying the photons. The number of photons with energy $h\nu$ depends on how many atoms lie in the $E_{1}$ level and how many in the $E_{2}$ level. In the case of equilibrium, these numbers remain unchanged.

However, the equilibrium is of a dynamic nature since the processes of excitation and radiation occur at the same time. In some way (either by collision with another particle or due to absorption of a photon from without) the atom or the atomic system climbs to a high level. The system persists in this excited state for some (in definite) time (ordinarily for a fraction of a second) and then reverts to a low level. This process is termed \emph{spontaneous radiation}. The atom behaves like a little ball on the top of a sharp peak with an intricate configuration: the slightest breath of air is enough to disrupt the equilibrium. The ball rolls down into a valley, usually the lowest part, and then only a strong impact can bring it out again. We say that an atom that has dropped to the lowest level is in a stable state.

But here we must bear in mind that there are also intermediate states in between the peak and the lowest portion of the valley. The ball may be at rest in a slight depression from which it can be extricated by a waft of air, so to speak, or at least by a little push. This is a metastable state. Thus, besides the excited and stable states there is a third, metastable, type of energy level.

To summarize, then, the transitions will occur in both directions. First one atom and then another atom will move into a higher level. In the next instant, they will fall to a lower level and emit light. But at the very same time, other atoms will receive energy and will rise to upper levels.

The law of conservation of energy requires that the number of transitions upwards equal the number of transitions downwards. What does the number of transitions upwards depend on? Two factors: first, the number of atoms in the lowest floor, and, second, the number of impacts that raise them to a higher floor. And the number downwards? It is of course determined by the number of atoms lying in the upper floor, and it would seem to be in dependent of any other factors. That is precisely what theoretical physicists thought at first, and yet the pieces didn’t fit. The number of transitions upwards, which is dependent on two factors, increased with temperature much faster than the number of transitions downwards, which is dependent on only one factor. This model, which had appeared to be so obvious, turned out to be nonsense: Sooner or later all the atoms would be chased up to the highest level: the system of atoms would be in an unstable state with no radiation.

It was precisely this impossible conclusion that Einstein, in 1926, picked up from the reasoning of his predecessors. Apparently, there was some other influence affecting the transitions of atoms from the upper floor to the lower floor. One could only conclude that there is a forced transition in addition to the spontaneous transition to the lower level.

What is this \emph{stimulated emission}, as it is called? In short, it is this. A system is in the upper level. It is separated from the lower level by the difference $E_{2} - E_{1}= h\nu$. Now, if a photon with energy $h\nu$ is incident on the system, then it makes the system move down to a lower level. The incident photon is not absorbed in the process but continues onwards accompanied by a fresh photon of exactly the same kind generated by the first one.

Do not seek any logic in this reasoning. It was intuition, a guess and experiment was to prove it right or wrong. Using the assumption of stimulated emission we are able to derive a quantitative formula that yields the graph of emission as a function of the wavelength of a heated body. The theory proved to be in brilliant agreement with experiment and so justified the hypothesis.

It is an exciting thought that the practical conclusions from the fact of the existence of stimulated emission that led to the invention of lasers were drawn only many years later.

\section{Optical Spectra}

Generally speaking, any body is a source of soft electromagnetic radiation. Using a spectrograph (this is an instrument whose principal component is a prism or a diffraction grating), light can be decomposed into a spectrum. The spectrum may turn out to be continuous, banded or line. The spectra of incandescent solids are very similar. For that matter, only a few substances can be heated to incandescence. A real rarity is a glowing liquid. The emission spectra of gases are highly informative. Such are the spectra of rays coming to us from distant stars. The bulk of the information we have concerning the structure of the universe comes to earth in the form of light rays of stellar matter in a gaseous state.

Under terrestrial conditions, it is easy to obtain emission spectra of atoms. Atoms are made to glow either by passing a current through the gas or by heating. In this manner we can obtain only the spectra of atoms but not the spectra of molecules. Before the gas begins to glow, the molecules break up into atoms. That is why absorption spectra are studied if the investigator is interested in liquids or solids. In the final analysis, the picture is determined by the system of energy levels. Transitions up or down yield the same information. Simply, do what is easiest.

Spectra that consist of separate clear-cut lines can be obtained only from a gas or a diluted solution. In the second book of this series it was stated that the behaviour of dissolved molecules resembles in many respects the behaviour of a gas. This also holds true for optical spectroscopy. Unfortunately, the solvent affects the character of the spectrum, but if we compare the type of spectra of molecules dissolved in different substances, it is possible to take that effect into account and extract from the experiment the fingerprints of the dissolved molecule.
%\newpage
\begin{center}
\includegraphics[width=0.8\textwidth]{figures/bohr.pdf}
\end{center}
{\small \textsf{{Niels Bohr [1885-1962]}} -- \textsf{\footnotesize the famous Danish physicist who creat­ed the first quantum model of the atom and thus discovered the law of quantization of energy. He was an active participant in developing the principles of quantum mechanics. He demonstrated the fundamental inapplicability -- to the microworld -- of concepts suitable in describing the behaviour of macroscopic bodies. He made a very considerable contribution to the theory of the structure of the atomic nucleus.}}

Obtaining a characteristic spectrum does not mean establishing the system of energy levels of a molecule. But for many practical purposes this is not required. With an album of information about spectra (that is, the list of spectral lines and their intensities, or the curves of intensity versus frequency) of some family of chemical substances, we can, by taking the spectrum of an unknown substance and comparing the experimental pattern with material from the album, determine the substance in the very same way that a criminal is detected from the fingerprints he leaves.


Just lately, optical spectral analysis has come up against a competitor: radiospectroscopy. Radio spectroscopic methods are still inferior in sensitivity to optical methods (though the inferiority will most likely not last long) but are far superior to optical methods in the identification and quantitative analysis of mixtures of substances.


We don’t aim here to acquaint the reader with concrete spectra of substances. It will suffice to discuss the pattern of energy levels of atoms of hydrogen and the fundamental scheme of energy levels of a free molecule.
\begin{figure}[!ht]
\centering
\includegraphics[width=0.5\textwidth]{figures/fig-01-03.pdf}
\caption{System of energy levels of hydrogen.}
\label{fig-1.3}
\end{figure}

\figr{fig-1.3} depicts the system of energy levels of hydrogen. Note the characteristic thickening of levels as we move away from the zero line.

Incidentally, the zero in the diagram is not a ``real'' zero actually. An unexcited atom of hydrogen naturally possesses some energy. But since spectra exhibit energy differences, it is convenient to reckon energies from the lower line. Depending on the intensity of the ``shock'' obtained, the atom can rise to any one of the ``floors'', hold on for a moment in the nonequilibrium state and then, via one of two possible modes (spontaneous emission or stimulated emission), revert to the lower level.

The resulting spectrum may conveniently be split up into a number of series. Each series is subordinate to its lower level. In the visible portion of the spectrum we have the so-called \emph{Balmer series}. The explanation of this series was the first triumph of Niels Bohr's theory of atomic structure.


Not all energy transitions are of equal probability. The higher the probability of transition, the more intense the appropriate line. There are also forbidden transitions.

A great achievement of theoretical physics was the exhaustive interpretation of the spectrum of the hydrogen atom via the solution of the famous equation of quantum mechanics derived in 1926 by the Austrian physicist Erwin Schr\"odinger (1887-1961).

Atomic spectra are affected by external fields. The lines split into several components under the action of an electric field (the \emph{Stark effect}) and under the action of a magnetic field (the \emph{Zeeman effect}). We will not go into these exciting phenomena here, but it is necessary to point out that an understanding of them came only after Samuel Goudsmit and George Uhlenbeck made the assumption that the electron possesses spin. How spin reveals itself in experiment directly was discussed in the third book of this series.

And finally the last remark regarding the pattern of energy levels. We see that the limit which the levels come up to is marked by the number 13.53. What kind of number is this? This is the \emph{ionization potential}. If we multiply the electron charge by the magnitude of this potential in volts, we obtain the amount of work needed to tear the electron away from the nucleus; in other words, this is the work that must be done to destroy the hydrogen atom.


Atomic spectra arise as a result of electron transitions. As soon as we move from atoms to a molecule, we immediately have to take into account two more components of energy. A molecule can rotate and the atoms of a molecule can perform vibrations with respect to one another. All these types of energy can likewise be quantized, which means they can have only specific discrete values. Thus, the energy state of a molecule is described by the state of its electron cloud (\emph{electronic level}), the state of oscillatory motion (\emph{vibrational level}), and the state of rotation (\emph{rotational level}). We thus have to deal with three kinds of information: the number of the house, the floor, and the flat.

\begin{figure}[!ht]
\centering
\includegraphics[width=0.8\textwidth]{figures/fig-01-04.pdf}
\caption{Schematic of electronic transfers in an atom with the required information.}
\label{fig-1.4}
\end{figure}

But what is the role of the ``floor'' and the ``flat''? What energy levels have big separations and what levels have small separations? All these questions are answered in \figr{fig-1.4}. This diagram shows two electronic levels $e'$ and $e''$ (the house numbers). The floors are the vibrational levels marked $\nu$, and the numbers of the flats are the rotational levels marked $j$. True, this differs somewhat from the ordinary numbering of houses and flats, which is continuous; in dealing with molecular spectra we number the flats on every floor from zero. Thus we see that the gaps between rotational levels are the smallest and between electronic levels ($e'$ and $e''$) the largest.

Suppose a molecule has the following possible electronic levels: 100, 200, 300, units of energy, vibrational levels at 10, 20, 30, \ldots\, units of energy, and rotational levels at 1, 2, 3, \ldots\, units of energy; then a molecule on the second electronic level, the first vibrational level, and the third rotational level will have an energy of 213 units.

Thus, we can give the energy of a molecule as follows:
\begin{equation*}%
W = W_{\textrm{el}} + W_{\textrm{vib}} + W_{\textrm{rot}}
\end{equation*}
The frequency of the emitted or absorbed light will always correspond to the difference (symbol: $\Delta$) of two levels:
\begin{equation*}%
\nu = \frac{1}{h} \left( \Delta W_{\textrm{el}} + \Delta W_{\textrm{vib}} + \Delta W_{\textrm{rot}} \right)
\end{equation*}

I would like to touch on those transitions that involve a change in only one type of energy. In practice, this occurs only in the case of rotational transitions, and why this is so will soon be seen.

We begin our investigation with the absorption of electromagnetic waves of a group of molecules starting with the longest wavelengths, that is, with the smallest portions of energy $h\nu$. Until the magnitude of the energy quantum has become equal to the distance between the two closest-lying levels, the molecule will not begin to absorb. By gradually increasing the frequency, we will reach quanta capable of raising the molecule from one ``rotational'' step to the next. Experiment shows that this occurs in the region of microwaves (the edge of the radio spectrum) or, to put it differently, in the region of the far infrared spectrum. Wavelengths of the order of 0.1 to 1 millimetre will be absorbed by the molecules. What we then have is a purely band spectrum.

New things happen when we irradiate the substance with energy quanta sufficiently high to move the molecule from one vibrational level to another. However, we will never attain a purely vibrational spectrum, that is, a series of transitions under which the number of the rotational level is preserved. On the contrary, transitions from one vibrational level to another will involve a variety of rotational levels. Say, a transition from the zero (lowest) vibrational level to the first can consist in moving up from the third rotational level to the second, or from the second to the first, and so forth. We thus obtain a vibration-rotation spectrum. We observe it in infrared light (from 3 to 50 micrometres). All transitions from one vibrational level to another will differ slightly in energy and will yield a group of very close lines in the spectrum. In the case of small resolution, these lines merge into a single band. Each band corresponds to a definite vibrational transition.

We thus move into a new spectral region, the region of visible light where the energy of the quantum becomes sufficient to move the molecule from one electronic level to another. Of course, it is not possible here to obtain either purely electron transitions or electron-vibrational transitions. Complex transitions arise in which the energy transition is accompanied by a change in the ``house'', the ``floor'',and the ``flat''. Since a vibrational-rotational transition constitutes a band, the spectrum in the visible region will be practically continuous.

The characteristic spectra of atoms and molecules have for many years played the role of helpers in determining the chemical structure and composition of substances. And their aid continues today. Revolutionary events in the field of spectroscopy have only just recently occurred.

\section{Laser Radiation}

The first thirty years of this century saw fantastic advances in theoretical physics with the discovery of such important laws of nature as the laws of the mechanics of high velocities, the laws of the structure of the atomic nucleus, and the laws of quantum mechanics. And the forty years that followed exhibited just as phenomenal a development in the applications of theory to practice.

This was a time when humanity harnessed the energy of atomic nuclei and invented semiconductor transistors that revolutionized radio engineering and led to the development of electronic computers and laser technology. These three applications were actually what produced the modern revolution in science and engineering.

In this section we discuss lasers. First let us give some thought to the problem of why, operating via traditional methods, we are not able to generate an intense directed beam of light.

The strongest stream of light collected into a very narrow beam disperses and loses its energy over small distances. It was only in the science fiction of the Russian writer Aleksei Tolstoi that the hero devises a ``hyperboloid'' capable of generating bursts of intense light rays that can burn and cut materials and carry tremendous energy over great distances. Of course we know that it is possible to manufacture a concave mirror that can generate a parallel beam of light. This requires placing a point source in the focus of the mirror. But a point is a mathematical abstraction. All right, suppose we have a small source, bigger than a point. But even then, if we heat the ball to 6000 degrees (and no known material can stand more), we obtain a beam of light of miserably low intensity. And as soon as we increase the dimensions of the source, then instead of a parallel beam of rays we obtain a spread-out fan of light ``filaments'' and the intensity of the ray of the projector begins to rapidly diminish with distance.

Thus, the first obstacle to creating a strong beam of light is that atoms emit light in all directions. That’s the first, but not the last. Atoms and molecules emit light without agreeing on how they’ll do it. The result is that rays from different atoms set out at different times, totally unmatched in their efforts. The emissions of different atoms do not agree in phase, and that means that rays from different atoms will frequently annihilate one another. This occurs, as you will recall, when the hump of one wave comes with the valley of another one.

\emph{Laser emission} is what overcomes these obstacles. The word ``laser'' stands for \textbf{l}ight \textbf{a}mplification by \textbf{s}timulated \textbf{e}mission of \textbf{r}adiation.

The underlying idea is made up of several elements. First of all, recall that there is stimulated emission and spontaneous radiation. We have already mentioned that this type of emission occurs when a light photon encounters an excited atom. If the excitation energy of the atom is exactly equal to the energy of the photon, then the photon de-excites the atom, which moves to a lower level and emits a photon. The marvelous peculiarity of stimulated emission is that the new photon is the same as the one that generated it, and not only in energy but also in phase and direction of motion.

The second element behind this idea is this. If the system of emitting atoms is placed in a tube, the ends of which are at a certain distance from each other and can serve as mirrors for the photons that interest us, then we can build up a bunch of photons generated (as they move back and forth) by identically excited atoms.
The third element in this idea is to retain the atoms in the excited state as long as possible and then, after the pumping is completed, to force all atoms to de-excite at the same time. Putting this laser idea (that is, the production of millions upon millions of identical photons from a single photon) into hardware should permit generating a light beam of unimaginable intensity. Such a beam would exhibit only the slightest spread and would have tremendous energy over its cross section.

But the question is: How can this be attained? For decades no one knew. Back in the 1930s, important ideas in this connection were expressed by Soviet physicist V. A. Fabrikant. Later, persistent research of Soviet scientists A. M. Prokhorov and N. G. Basov and, independently, the American physicist Charles Hard Townes led to the invention of lasers. All three received the Nobel prize in physics.

\begin{figure}[!ht]
\centering
\includegraphics[width=0.4\textwidth]{figures/fig-01-05.pdf}
\caption{Two level system for laser processes.}
\label{fig-1.5}
\end{figure}

Suppose a system has two energy levels. Most atoms or molecules are in the lower level. Thermal shocks can transfer a molecule to the upper level for a short time. But not for long, because the molecule is de-excited. In the process most of atoms go to the lower level spontaneously. Of course, some of emitted photons will carry some of the excited atoms to the lower state and generate photons of stimulated emission. But these are rare processes because there are few excited particles (the most occupied are the lower levels), and also the probability of a spontaneous transition is substantially higher than the probability of stimulated emission.

Let us suppose it was possible to find a substance whose atoms have the three energy levels marked in \figr{fig-1.5} by the numerals 1, 2, and 3. The distance 1-3 corresponds to the frequency of emission of green light, the distance 1-2 corresponds to the frequency of red light. Now suppose the probability of a transition from level 3 to level 2 is thousands of times higher than the frequency of transitions from level 2 to level 1. Let us irradiate the substance with green light. The atoms will rise to the third floor, then via spontaneous transitions will go to level 2 and will stay at that level. This is termed \emph{nonradiative transition}. The energy released goes into the vibrational energy of the atoms. Using our imagination further, let us suppose that we have carried most of the atoms to level 2. We have thus reversed the occupancy density, it is no longer ``normal''. There are more in the upper levels 2 than in the lower levels 1, which is impossible when the process is controlled solely by thermal motion.

And still and all there does begin a transition from level 2 to the lower level 1. An appropriate photon will encounter other atoms in the excited level 2. The result will be not absorption but the creation of a new photon. The first, accidentally generated photon in 2-1 will be joined by the very same photons of stimulated emission.

Thus arises a stream of 2-1 photons. They will all be identical and will generate a beam of tremendous intensity.

That precisely was the process that the three Nobel prize winners were able to create. Historically, the first was a ruby laser. The diagram of levels shown in the figure is precisely the diagram of ruby with an admixture of chromium atoms.

To make a laser we need a source of excitation that does the pumping of the laser, that is, that carries the atoms to higher levels.

If the source of laser emission is a solid, then it is made in the form of a cylinder whose bases play the part of mirrors. In the case of a liquid or gaseous laser, a tube is constructed with mirrors at the ends of the column. By performing a micrometre-precise positioning of the mirrors and thus fixing the length of the column, we put in a privileged position only those photons whose integer number of wavelengths fit into the length of the column. Only then do all the waves combine.

Perhaps the main peculiarity of the laser is the possibility of creating a narrow stream of radiation. Practically speaking, a laser beam can have almost any cross section. Technically, this is achieved by the fact that the ray is made to travel along a narrow glass capillary tube of sufficient length. Photons moving at an angle to the capillary do not participate in the photon build-up. A resonance cavity (that is, the mirrors that reflect photons first in one direction and then in the other during the pumping period in the operation of the laser) reproduces photons of only one direction. In some cases, if an angular dispersion of the beam of the order of one degree is not satisfactory, a supplementary lens is placed in the path of the released ray.

A laser device is a complicated piece of engineering if one has to do with high power outputs. A primary pulse is first set up in the column; then it is fed to amplifiers that function in the same manner as the first column but pump independently of the first column. We will not go into these details because we are only interested in the physical principles of pumping and the generation of laser emission. They can differ greatly, as is evident from a glance at \figr{fig-1.6} to \figr{fig-1.8} with diagrams of the action of lasers which today yield beams of maximum power output.

\begin{figure}[!ht]
\centering
\includegraphics[width=0.4\textwidth]{figures/fig-01-06.pdf}
\caption{Schematic of a neodymium laser.}
\label{fig-1.6}
\end{figure}


\figr{fig-1.6} depicts a so-called neodymium laser. Actually, the body of the laser is not the metal neodymium but ordinary glass with an admixture of neodymium. Ions of neodymium atoms are haphazardly distributed among the atoms of silicon and oxygen. Pumping is performed by flash bulbs. The lamps emit in wavelengths between 0.5 and 0.9 micrometre, and a broad band of excited states is obtained (shown here in the form of five bars). The atoms perform nonradiative transitions to the upper laser level (labelled 2 in all three figures). Each transition yields a different energy, which is converted into the vibrational energy of the whole ``lattice'' of atoms.

Laser emission, that is, the transition to an empty lower level labelled 1, has a wavelength of 1.06 micro metres. 

The dashed line, which depicts the transition from level 1 to the lowest level ``does not work'' (in the sense that energy is released in the form of noncoherent radiation).

A neodymium laser permits obtaining fantastic power outputs of up to \SI{d12}{\watt}. The energy is generated in the form of pulses lasting 0.1 nanosecond.

\begin{figure}[!ht]
\centering
\includegraphics[width=0.4\textwidth]{figures/fig-01-07.pdf}
\caption{Schematic of an iodine laser.}
\label{fig-1.7}
\end{figure}

A new competitor in this field is a laser using transitions in excited atoms of iodine (\figr{fig-1.7}). The working substance here is the gas \ce{C3F7I}. Here, too, flash bulbs are used for pumping, but the physical processes are different. Ultraviolet light of wavelength 0.25 micrometre is used for pumping. Under the action of this radiation there occurs a dissociation of the molecules. The remarkable thing is that the iodine atoms are torn out of the molecule and are in an excited state! As the reader will see, this is quite a different method for inverting the occupancy density. The operating transition $2 \to 1$ leads to laser emission with a wavelength of 1.3 micrometres, after which the iodine atom joins up with the molecular residue.

The reader has also probably heard of the widespread use of helium-neon lasers. They are used to obtain an intense infrared ray of wavelength 1.13 micrometres. These lasers are not record holders as far as power goes, and so we give a diagram of levels for a different laser that operates on a mixture of nitrogen and carbon dioxide (\figr{fig-1.8}).

\begin{figure}[!ht]
\centering
\includegraphics[width=0.4\textwidth]{figures/fig-01-08.pdf}
\caption{Schematic of nitrogen and carbon dioxide based laser.}
\label{fig-1.8}
\end{figure}
Before describing this laser, a natural question comes to mind, and that is: Why are mixtures of gases needed? The point is that some atoms and molecules are more easily excited while others are more easily de-excited, so that in a laser operating on a mixture, particles of one type effect the pumping process, these then transfer the energy via collisions to other atoms or molecules, and these in turn generate the laser beam. There are systems now functioning that consist of more than two gases.
For instance, in the nitrogen-carbon dioxide laser, it is advisable to add a variety of components including helium.

Pumping in the \ce{CO2} laser is done differently from the two just described. The mixture of gases is placed in a gas-discharge tube and a sufficiently high voltage is applied so that the system becomes a plasma. Electrons moving at high speeds excite the vibrations of nitrogen molecules. The diagram shows a transition of this molecule to the upper floor. The voltage applied to the electrodes plays a delicate role. The optimal energy for exciting nitrogen molecules is about \SI{2}{\electronvolt}. The nitrogen molecule is only an intermediary. It does not by itself produce any emission but rather transfers the energy obtained from the electrons to the \ce{CO2} molecule and lifts it to an upper laser level.

The upper laser levels 2 are ``flats of the third floor'' of \ce{CO2} molecules. A molecule of gas has a lifetime in the upper laser level equal to about 0.001 second. This is not so little, and the molecule has a good chance of encountering a photon of suitable energy that will force it down to a lower level.

It should be pointed out that ``interflat'' transitions are a much more frequent occurrence than ``interfloor'' transitions. Lifetimes on the rotational level are of the order of ten millionths of a second. This favourable circumstance results in the flats of each floor being occupied in a rather stable fashion, and so, using an engineering technique that we have already mentioned (setting a suitable distance between the mirrors), it is possible to isolate some one transition; let us say, from the sixth flat of the third floor to the fifth flat of the second floor.

The designing engineer must have at his disposal complete information about the residence time of an atom on any given sublevel and about the probabilities of transition. Then he is able to choose the optimal radiation of the given gaseous mixture. Ordinarily, a laser operating on carbon dioxide is tuned to a wavelength of \num{10.5915} micrometres. For a laser to function normally, it is necessary that the molecules not pile up on the lowest laser level, the idea being for them to do their jobs and get out. Now, at a gas pressure of 1 millimetre of mercury, carbon dioxide molecules experience 100 collisions per second in vacating the level. The respective figures are \num{4000} and \num{100000} if helium and water are present. This is a tremendous difference.

By selecting the proper admixtures for carbon dioxide, we can boost the power output of the instrument substantially. It would appear that this is the gold-medal winner in the laser field.

A carbon dioxide (\ce{CO2}) laser produces a beam that can be focussed on an area of \SI{0.001}{\centi\meter\squared} with an intensity of \SI{1000}{\kilo\watt\per\centi\meter\squared} in continuous operation and one million \si{\kilo\watt\per\centi\meter\squared} in pulsed operation with the pulse time equal to one nanosecond (which, as you know, is \num{d-9}, or one thousand millionth, of a second).

The search for suitable materials for lasers is a sort of art. One needs intuition, ingenuity, and memory to create an effectively operating laser. The user can now order lasers with a great variety of wavelengths ranging from a tenth of a micrometre to hundreds of micrometres.

The exceptional intensity and coherence of laser emission have revolutionized many areas of engineering. During the past decade, the manufacture of lasers has become a whole industry. Lasers have found application as generators of radiation that trasmit not only energy but also information. Intense research is in progress for their use in initiating thermonuclear reactions. Lasers are used in place of the surgeon's knife in medicine, as an instrument for the most delicate surgical operations, as a device for the separation of isotopes. We will have occasion later on to come back to further discussions of the marvelous laser.

\section{Luminiscence}
Thermal radiation is a universal property of all bodies. Thermal rays are emitted by every body at any temperature from absolute zero upwards. The thermal spectrum is a continuous one and is depicted by a curve that we have already discussed. True, our curve was that of a black body, but the behaviour of coloured bodies is in principle but slightly different from the behaviour of black bodies. Merely, the curve for coloured bodies is distorted somewhat. But the general increase in the energy of emission (as the temperature rises) and the displacement of the maximum to the left (if wavelengths are laid off on the axis of abscissas) are the general law.

All radiation consists in a transition from a higher energy level to a lower level. But the reasons for excitation of atoms or molecules may differ. In the case of
thermal radiation, it is the collisions of particles of the substance due to thermal motion.

But that is not the only reason compelling a body to emit waves. \emph{Luminescence}, which we are about to discuss, is of a different nature. This term embraces processes of excitation of molecules that are not connected with any increase in the temperature of the body. Causes of particle excitation may be encounters with beams of photons or electrons, mechanical impact, friction, and so on.

Practically all substances are capable of luminescence. But only some (called luminophors or phosphors) glow brightly and are of practical importance.

Luminophors are used as materials for covering television and oscillograph screens, in which case the luminescence occurs under the impact of electrons. Certain substances luminesce brightly under the action of ultraviolet radiation. The energy of the incident photon must be at least greater than the energy of the emitted photon. That is why the incident quantum of energy can come from the invisible portion of the spectrum while the emitted radiation can lie in the visible portion.

Admixtures of luminescent material measured in minute fractions (billionths, or \num{d-9}) are sufficient to make a substance luminesce under irradiation by ultraviolet light. That is why fluorometric analysis is sometimes used as a tool in chemical analysis. It is capable of detecting minute quantities of impurities.

Luminophors are used to cover the walls of daylight lamps.

There are two types of luminescence: fluorescence and phosphorescence. \emph{Fluorescence} consists in the de-excitation of an atom or molecule that occurs without the molecule remaining in the excited level. Contrariwise, \emph{phosphorescence} persists after the excitation has ceased. This occurs if the system, when excited, passes to a metastable level, from which transitions downwards have a low probability. As a rule, the radiation occurs after the molecule first absorbs the energy and rises to a higher level, after which de-excitation takes place, the transition to the lower level occurring without any stop at an intermediate, metastable, level.

A few words about \emph{electroluminescence} that occurs in certain semiconductor diodes on the boundary of a $p\!-	\!n$ region. This interesting phenomenon is of great practical value because it underlies the manufacture of semiconductor lasers. The idea is this: an electron and hole of the semiconductor can recombine with the emission of a photon.

For transitions of this type to take place continuously we have to pass an electric current through the diode. The problem is to find a suitable material that satisfies several requirements. First of all, the current has to inject (if that is the word) electrons into the $p$-type semiconductor material, that is, a semiconductor containing more holes, or it must pump holes into an $n$-type crystal. This is a necessary condition, but other factors such as, for example, the rate of transition from an upper to a lower level can play a decisive role. Then there are cases where all factors favour a transition of an electron downwards and electroluminescence takes place.

A particularly good electroluminescent material is the semiconductor gallium arsenide. It yields a sufficient quantity of photons. The photons move along the $p\!-\!n$ boundary. Two sections of the diode perpendicular to the boundary are polished and that sets up a resonant cavity. The photons generated in recombinations of holes and electrons are in phase, and for sufficiently large currents the radiation becomes stimulated emission with all the resultant consequences of being narrow, highly directed, and polarized.

Semiconductor lasers operate in a band of wavelengths from the ultraviolet to the far infrared and are widely used for a variety of purposes.


%\begin{figure}[!ht]
%\centering
%\includegraphics[width=\textwidth]{figures/fig-01-01.pdf}
%\caption{Using springs for measuring weights.}
%\label{fig-1.1}
%\end{figure}





% !TEX root = pfe-book2.tex
%!TEX TS-program = pdflatex
%!TEX encoding = UTF-8 Unicode


\cleardoublepage
%\mainmatter
\chapter{Structure of Matter}
\label{ch-02}

\section{Intramolecular Bonds}
Molecules consist of atoms. Atoms are bound in mol­ecules by forces which are called chemical forces.

There exist molecules consisting of two, three, four atoms. The largest molecules, protein molecules, consist of tens and even hundreds of thousands of atoms.

The molecule kingdom is exceptionally varied. By now, millions of substances built up out of various molecules have already been isolated by chemists from natural materials and created in their laboratories.

Properties of molecules are determined not only by how many atoms of one or another sort participate in their construction but also by the order and configura­tion in which they are bound. A molecule is not a heap of bricks, but a complicated architectural structure, where each brick has its place and its completely de­termined neighbours. The atomic structure forming a molecule can be rigid to a greater or lesser degree. In any case, each of the atoms carries out an oscillation about its equilibrium position. In certain cases, some parts of a molecule can even revolve around other parts giving different and the most fantastic configurations to free molecule in the process of its thermal motion.

Let us analyze the interaction between atoms in greater detail. The potential energy curve of a diatomic mole­cule is depicted in \figr{fig-2.1}. It has characteristic form -- it first goes down, then turns up forming a ``well'', and afterwards rises more slowly towards the horizontal axis on which the distance between the atoms is marked.

\begin{figure}[!ht]
\centering
\includegraphics[width=0.4\textwidth]{figures/fig-02-01.pdf}
\caption{The characteristic form of potential energy curve, a potential ``well'' for a diatomic molecule.}
\label{fig-2.1}
\end{figure}

We know that the state in which the potential energy has the minimum value is stable. When an atom forms a part of molecule, it ``sits'' in potential well, carrying out small thermal oscillations about its equilibrium position.

The distance from the vertical axis to the bottom of the well can be called the equilibrium distance. The atoms would be located at this distance if the thermal motion were to cease.

The potential energy curve tells about all the details of the interaction between atoms. Whether particles attract or repel each other at one or another distance, whether the strength of the interaction increases or decreases when the particles separate or approach all this information can be obtained from the analysis of the potential energy curve. Points to the left of the ``bot­tom'' of the well correspond to repulsion. On the contrary, points to the right of the bottom of the well characterize attraction. The steepness of the curve also yields important information: the steeper the curve, the greater the force.

When atoms are at great distances from each other, they are attracted; this force decreases rather rapidly with an increase in the distance -- between them. As they approach each other, the force of attraction grows and readies its maximum value when the atoms come very close to each other. As they come even closer, the attraction weakens and, finally, at the equilibrium distance the force of the interaction vanishes. When the atoms become closer than the equilibrium distance, forces of repulsion arise which sharply increase and quickly make a further decrease in the distance between the atoms practically impossible.

Equilibrium distances (below we shall say distances for the sake of brevity) between atoms are different for various types of atoms.

For various pairs of atoms, not only are the distances between the vertical axis and the bottom of the well different but so are the depths of the wells.

The depth of well has a simple meaning: in order to roll out of the well, an energy just equal to the depth is needed. Therefore, the depth of a well can be called the binding energy of the particles.

The distances between the atoms of molecule are so small that it is necessary to choose appropriate units for their measurement; otherwise, their values would have to be expressed, for example, in the following form: \SI{0.000000012}{\centi\meter}. This figure is for an oxygen molecule.

Units especially convenient for describing the atomic world are called angstroms (true, the name of the Swedish scientist in whose honour these units were named is properly spelt \AA ngstrom; in order to remember this, a small circle is placed over the letter A);
\begin{equation*}%
\SI{1}{\angstrom} = \SI{d-8}{\centi\meter}
\end{equation*}
\SI{1}{\angstrom} is one-hundred-millionth of a centimetre.

The distances between the atoms of a molecule lie within the limits of 1 to \SIrange{1}{4}{\angstrom} The equilibrium distance for oxygen, which has been written out above, is equal to \SI{1.2}{\angstrom}.

Interatomic distances, as you see, are very small. If we gird the Earth with a string at the equator, then the length of the ``belt'' will be as many times greater than the width of your palm as the latter is greater than the distance between the atoms of a molecule.

Ordinary calories are used for measuring binding energies, but they are related not to one molecule, which would, of course, yield a negligible number, but to one mole, i.e. to the number of grams equal to the relative molecular mass.

It is clear that the binding energy per mole divided by Avogadro’s number, $N_{A} = \SI{6.023d23}{\per\mole}$ yields the binding energy of a single molecule.

The binding energy of the atoms in a molecule, just as interatomic distances, varies within narrow limits. For the same oxygen, the binding energy is equal to
\SI{116 000}{\calorie\per\mole}, for hydrogen \SI{103000}{\calorie\per\mole}, etc. We have already said that the atoms in a molecule are distributed in an entirely definite manner with respect to each other, forming in complicated cases rather intricate structures.

Let us present several simple examples. In a molecule
of \ce{CO2} (carbon dioxide), all three atoms are lined up in a row, with the carbon atom in the middle. A molecule of \ce{H2O} (water) has an angular form, with the oxygen atom at the vertex of the angle (it is equal to \ang{105}).

In a molecule of \ce{NH3} (ammonia), the nitrogen atom
is at the vertex of a three-faced pyramid, in a molecule of \ce{CH4} (methane), the carbon atom is located at the centre of a four-faced figure with equal sides, which is called a tetrahedron.

The carbon atoms of \ce{C6H6} (benzene) form a regular hexagon. The bonds of the carbon atoms with the hydro­gen atoms go from all the vertices of the hexagon. All the atoms are situated in one plane.

\begin{figure}[!ht]
\centering
\includegraphics[width=0.4\textwidth]{figures/fig-02-02.pdf}
\caption{Structural form and distribution of atoms in carbon dioxide \ce{CO2} and water \ce{H2O}.}
\label{fig-2.2}
\end{figure}

Diagrams of the distribution of the centres of the atoms in these molecules are shown in \figr{fig-2.2} and \ref{fig-2.3}. The lines symbolize the bonds.

\begin{figure}[!ht]
\centering
\includegraphics[width=0.4\textwidth,angle=-1]{figures/fig-02-03.pdf}
\caption{Structural form and distribution of atoms in methane \ce{CH4} and benzene \ce{C6H6}.}
\label{fig-2.3}
\end{figure}

A chemical reaction occurred; there were molecules of one type, and then other, were formed. Some bonds were broken, while others, were newly created. In order to break the bonds between atoms (recall \figr{fig-2.1}) it is necessary to perform work, just as in rolling a ball out of a well. On the contrary, energy is liberated when new bonds are formed, just as when a ball rolls into a well.

Which is greater, the work involved in breaking or in creating bonds? We come across reactions of both types in nature.

The excess energy is called the \emph{thermal effect}, or more concisely, the \emph{heat of transformation (reaction)}. The heat of reaction is usually a quantity of the order of tens of thousands of calories per mole. The heat of reac­tion is often included as a summand in the formula for a reaction.
For example, the reaction whereby carbon in the form of graphite burns, i.e. unites with oxygen, is written out as follows:
\begin{equation*}%
\ce{C + O2 -> CO2 + \SI{94250}{\calorie}}
\end{equation*}
This means that when carbon combines with oxygen an energy of \num{94250} calories is liberated.

The sum of the internal energies of a mole of carbon and a mole of oxygen is equal to the internal energy of a mole of carbon dioxide plus \num{94250} calories.

Thus, such formulas have the transparent meaning of algebraic equalities written in terms of the values of the internal energies.

With the aid of such equations, one can find the heats of reaction for which direct methods of measurement, as a result of one or another cause, are unsuitable. Here is an example: if carbon (graphite) were to combine with hydrogen, the acetylene would be formed:
\begin{equation*}%
\ce{2C + H2 -> C2H2}
\end{equation*}
The reaction does not proceed in this manner. Never­theless, it is possible to find its thermal effect. We write down three known reactions:

\begin{enumerate}
\item the oxidation of carbon
\begin{equation*}%
\ce{2C + 2O2 -> 2CO2 + \SI{188000}{\calorie}}
\end{equation*}
\item the oxidation of hydrogen 
\begin{equation*}%
\ce{H2 + 1/2O2 -> H2O + \SI{68000}{\calorie}}
\end{equation*}
\item the oxidation of acetylene 
\begin{equation*}%
\ce{C2H2 + 5/2O2 -> 2CO2 + H2O + \SI{312000}{\calorie}}
\end{equation*}
\end{enumerate}

All these equalities may be regarded as equations for the binding energies of molecules. If so, we may operate on them as on algebraic equalities. Subtracting the first two equalities from the third, we obtain:
\begin{equation*}%
\ce{2C + H2 -> C2H2 - \SI{56000}{\calorie}}
\end{equation*}

Therefore, the reaction we are interested in is accom­ panied by the consumption of \num{56000} calories per mole.

\section{Physical and Chemical Molecules}

Until investigators had formed a detailed concept of the structure of matter, no such distinction was made. A molecule was simply a molecule, i.e. the smallest representative of a substance. It seemed that nothing more could be said. This is not so, however.

The molecules we have just discussed are molecules in both senses of the word. Molecules of carbon dioxide, ammonia and benzene (mentioned above), and the mol­ecules of practically all organic substances (which were not discussed) consist of atoms strongly bonded to one another. These bonds are not ruptured by dissolution, melting or evaporation. The molecule continues to be­ have as a separate particle or small physical body upon any physical action or change in state.

But this is not always true. For most inorganic sub­stances, we can speak of the molecule only in the chem­ical sense. The finest particles of such well-known inorganic substances as common salt or calcite or soda do not even exist. We do not find separate particles of these substances in crystals (this will be discussed a few pages further on); when they are dissolved, the molecules break down into their component atoms.

Sugar is an organic substance. Therefore, the sugar dispersed in a cup of sweetened tea is in the form of molecules. Salt is a different matter. We find no molecules of common salt (sodium chloride) in salty water. These ``molecules'' (we have to use quotation marks) exist in water in the form of atoms (actually, ions -- electrically charged atom -- that will be discussed later).

The same is true of vapours; and in melts a part of, the molecules live their own independent lives. 

When we speak of the forces binding the atoms together in a physical molecule, we call them valence forces. Intermolecular forces are not of the valency kind. The general shape of the interaction curve, of the type il­lustrated in \figr{fig-2.1}, is the same for both kinds of forces. The difference lies in the depth of the potential well. For valence forces the well is hundreds of times deeper.


\section{Interaction of Molecules}
There can be no doubt of the fact that molecules at­tract each other. If they stopped doing so for an instant, all liquids and solids would decompose into molecules.

Molecules repel each other, and neither can this be doubted, because liquids and solids would otherwise contract with extraordinary ease.
Forces are exerted between molecules which resemble in many respects the forces between atoms spoken of above. The potential energy curve which we have just drawn for atoms gives a true picture of the basic features of molecular interaction. However, there are also es­sential differences between these interactions.

Let us compare, for example, the equilibrium distances between oxygen atoms forming a molecule and oxygen atoms of two neighbouring molecules attracted in solidified oxygen before the equilibrium position. The difference will be very noticeable: the oxygen atoms forming a molecule settle down at a distance of \SI{1.2}{\angstrom}, while the oxygen atoms of different molecules approach each other to within \SI{2.9}{\angstrom}.

Analogous results have also been obtained for other atoms. Atoms of different molecules settle down farther from each other than atoms of the same molecule. It is therefore easier to tear molecules apart from each other than atoms from a molecule; moreover, the difference in energy is much greater than that in distance. While the energy necessary for breaking the bonds between oxygen atoms forming a molecule is about \SI{100}{\kilo\calorie\per\mole}, the energy needed to pull oxygen molecules asunder is less than \SI{2}{\kilo\calorie\per\mole}.

Hence, on a potential energy curve for molecules, the potential well lies farther away from the vertical axis and, furthermore, the well is much shallower.

However, this does not exhaust, the, difference between the interaction between atoms forming a molecule and the interaction of molecules.

Chemists have shown, that atoms are bound in a mole­cule with a fully determined number of other atoms. If two hydrogen atoms have formed a molecule, no third atom will join them to this end. A oxygen atom in water is bound to two hydrogen atoms, and it is im­possible to bind another atom to them.

We do not find anything similar in intermolecular interaction. Having attracted one neighbour to itself, a molecule does not lose its ``attractive force'' to any degree. The approach of neighbours will continue, as long as there, is enough room.

What does ``there is enough room'' mean? Are molecules really something, like apples or eggs? Of course, in a certain sense such a comparison is justified: molecules are physical bodies possessing definite size and shape. The equilibrium distance between molecules is nothing but their size.


\section{What Thermal Motion Looks Like}
The interaction between molecules can have greate or smaller values during the ``lives'' of the molecules.

The three states of matter -- gaseous, liquid and solid -- differ from one another in the role which molecular interaction plays in them. The word ``gas'' was thought up by scientists (derived from the Greek \emph{chaos} meaning ``disorder'').

And as a matter of fact, the gaseous state of matter is an example of the existence in nature of complete, perfect disorder in the mutual distribution and motion of particles. There is no microscope which would permit one to see the motion of gaseous molecules, but in spite of this, physicists can describe the life of this invisible world in sufficient detail.

There is an enormous number of molecules, approximate­ly \num{2.5d19} molecules in a cubic centimetre of air under standard conditions (room temperature and atmo­spheric pressure). Each molecule’s share is a volume of \SI{4d-20}{\centi\meter\cubed}, that of a small cube whose sides are approximately $\SI{3.5d-7}{\centi\meter} = \SI{35}{\angstrom}$. However, the molecules are much smaller. For example, molecules of oxygen and nitrogen -- the basic components of air -- have an average size of about \SI{4}{\angstrom}.

Therefore, the average distance between molecules is ten times as great as the size of the molecules. And this in turn implies that the average volume of air per molecule is approximately a thousand times as great as the volume of the molecule itself.

Imagine a plane surface on which coins have been thrown in a random manner where there is an average of a hundred coins to each square metre. This means one or two coins on a page of the book which you are reading. This is roughly how sparsely gas molecules are distrib­uted.

Every molecule of a gas is in a state of continual thermal motion.

Let us follow a single molecule. Here it is swiftly moving somewhere to the right. If it met no obstacles in its path, the molecule would continue its motion with the same velocity along a straight line. But the path of the molecule is crossed by its innumerable neighbours. Collisions are inevitable, and the molecules fly apart like two colliding billiard balls. In which direction will our molecule gallop? Will it acquire or lose its speed? Anything is possible: for its collisions can be of the most various kinds. Blows are possible from the front or from behind, from the right or the left, which are strong or weak. It is clear that being subject to such irregular impacts during these random collisions, the molecule which we are observing will rush about through all parts of the vessel in which gas is confined.

How far are gas molecules able to go without a col­lision?

It depends on the size of the molecules and the density of the gas. The larger the molecules and the more mole­cules there are in a vessel, the more often will they collide. The average distance travelled by a molecule without any impact -- it is called the mean free path -- is equal to $\SI{11d-8}{\centi\meter} = \SI{1100}{\angstrom}$ for hydrogen mole­cules and $\SI{5d-6}{\centi\meter} = \SI{500}{\angstrom}$ for oxygen molecules under ordinary conditions. The distance of \SI{5d-6}{\centi\meter} (one-twenty-thousandth of a millimetre) is very small, but it is far from small in comparison with molecular sizes. A distance of \SI{10}{\meter} for a billiard ball corresponds in scale to a path of \SI{5d-6}{\centi\meter} for an oxygen molecule.

The structure of a liquid differs essentially from that of a gas whose molecules are far from each other and only rarely collide. In a liquid, a molecule is constantly found in the immediate vicinity of others. The molecules of a liquid are distributed like potatoes in a sack. True, with one distinction: the molecules of a liquid are in a state of continual and chaotic thermal motion. Because they are so crowded, they cannot move around as freely as the molecules of a gas. Each of them is always ``marking time'' in practically one and the same place surrounded by the same neighbours, and only gradually moves through the volume occupied by the liquid. The more viscous the liquid, the slower this displacement. But even in such a ``mobile'' liquid as water, a molecule moves \SI{3}{\angstrom} during the time required by a gas molecule to cover \SI{700}{\angstrom}.

The forces of interaction between molecules deal very resolutely with their thermal motion in solids. In solid matter, the molecules are almost always in a fixed posi­tion. The only effect of the thermal motion is that the mole­cules are continually vibrating about their equilibrium positions. The lack of systematic displacements by the molecules is precisely the cause of what we call solidity. In fact, if molecules do not change neighbours, all the more will the separate parts of the body remain in a fixed bond with one-another.

\section{Compressibility of Bodies}

As raindrops drum on a roof, so do gas molecules beat against the walls of a vessel. The number of these blows is immense, and it is their united action that creates the pressure which can move the piston of an engine, explode a shell or blow up a balloon. A hail of molecular blows -- this is atmospheric pressure, this is the pressure that makes the lid of a boiling tea-kettle jump, this is the force driving a bullet out of a rifle.

What is gas pressure related to? It is clear that the
stronger the blow inflicted by a single molecule, the greater the pressure. It is no less obvious that the pressure will depend on the number of blows inflicted in a second. The more molecules in a vessel, the more frequent the blows and the greater the pressure. Hence, the pressure $p$ of a given gas is proportional, first of all, to its density.

If the mass of gas is constant, then decreasing its volume, we increase its density by the corresponding factor. Therefore, the pressure of a gas in a closed vessel will be inversely proportional to its volume. Or, in other words, the product of the pressure by the volume must be constant:
\begin{equation*}%
pV = \textrm{const}
\end{equation*}
This simple law was discovered by the English physi­cist Robert Boyle (1627-1691) and the French scientist Edme Mariotte (c.~1620-1684). \emph{Boyle's law} (also known as \emph{Mariotte's law}) is one of the first quantitative laws in the history of physical science. Of course, it holds when the temperature is constant.

As a gas is compressed, the Boyle equation is satisfied worse, and worse. The molecules approach each other and, the interactions between them begin to influence the, behaviour of the gas.

Boyle’s law is valid in those cases when the inter­ference of the forces of interaction between the gas mole­cules is completely insignificant. One therefore speaks of Boyle’s, law as a law of ideal gases.

The adjective ``ideal'' sounds rather funny when modi­fying the word ``gas''. Ideal means perfect, so that it is impossible to be better. 

The simpler a model or diagram, the more ideal it is for the physicist. Computations are simplified, explana­tions of physical phenomena become easy and clear. The term ``ideal gas'' pertains to the simplest model of a gas. The behaviour of sufficiently rarefied gases is practically indistinguishable from that of ideal gases.

Liquids are much less compressible than gases. In a liquid, the molecules are already in ``contact''. Compres­sion consists only in improving the ``packing'' of the molecules, and for very high pressures, in pressing the molecules themselves. The degree to which the forces of repulsion hinder the compression of a liquid can be seen from the following figures. A rise in pressure from one to two atmospheres entails a decrease in the volume of a gas by a factor of two, while the volume of water changes by 1/\num{20000} and that of mercury by total of 1/\num{250000}.

Even the enormous pressure in the depths of an ocean is incapable of compressing water to any noticeable extent. In fact, a pressure of one atmosphere is created by a ten-metre column of water. The pressure under a \SI{10}{\kilo\meter} layer of water is equal to \SI{1000}{\atmos}. The volume of water decreases by 1000/\num{20000}, i.e. by one-twen­tieth.

The compressibility of solids differs little from that of liquids. This is understandable since in both cases the molecules are already in contact, and so compression can only be achieved at the expense of a further drawing together of molecules which are already strongly repelling each other. By means of ultrahigh pressures of 50-100 thou­sand atmospheres, we are able to compress steel by one-thousandth, and lead by one-seventh, of its volume.

It is clear from these examples that, under terrestrial conditions, we cannot succeed in compressing solid mat­ter to any significant extent.

But in the Universe, there are bodies where matter is compressed with incomparably greater strength. Astron­omers discovered the existence of stars in which the density of matter reaches \SI{d8}{\gram\per\centi\meter\cubed}. Inside these stars, called white dwarfs (``white'' for the nature of their lumi­nosity, ``dwarfs'' because of their relatively small size), there should therefore be enormous pressures.


\section{Surface Tension}

Is it possible to emerge dry from water? Of course it is, if one smears oneself with a non-wettable substance. Rub your finger with paraffin and put it under water.

When you take it out, you will find that except for two or three drops there is no water on your finger. A slight motion -- and the drops are shaken off.

In this case we say that water does not wet paraffin. Mercury behaves in such a manner towards almost all solid bodies: it does not wet leather, glass or wood. 

Water is more capricious. It adheres closely to some bodies and tries not to touch others. Water does not wet oily surfaces, but thoroughly wets clean glass. Water wets wood, paper and wool.

If a drop of water is placed on a clean plate of glass, it will spread out and form a very shallow, small puddle. If such a drop is put on a piece of paraffin, it will just remain a drop, almost spherical in shape and slightly flattened by gravity.

Among the substances which ``stick'' to almost all bodies is kerosene. Striving to flow along glass or metal, kerosene is capable of creeping out of a loosely closed vessel. A puddle of spilled kerosene can spoil one’s existence for a long time: kerosene will seize a large sur­face, creep into cracks and penetrate one’s clothes. That is why it is so difficult to get rid of its not very pleasant odour.

The failure to wet bodies can lead to curious phenomena. Take a needle, grease it and carefully place it flat on water. The needle will not sink. Looking attentively, you can notice that the needle depresses the water and calmly lies in the small hollow so formed. However, a slight pressure is enough to make the needle go to the bottom. For this it is necessary that a considerable part of it turns out to be in the water.



 This interesting property is made use of by water strid­ers running swiftly along the surface of the water with­ out wetting their feet.
 
Wetting is used to dress ores by means of floatation. The word ``floatation'' means ``surfacing''. The essence of the phenomenon is as follows. Finely crushed ore is loaded into a vat containing water. A small amount of special oil is added; this oil must wet the particles of the mineral and not wet the particles of the gangue (this is what the valueless rock or aggregates of minerals in an ore are called). When mixed, the particles of the mineral are coated with an oily film.

Air is blown into the black mixture of ore, water and oil. There is formed a mass of little bubbles of air -- foam. The air bubbles come to the surface. The process of floatation is based on the fact that the particles coated with oil cling to the air bubble. A large bubble drags up a small particle like a balloon.

The mineral collects at the surface in the form of foam and the gangue remains at the bottom. The foam is moved and sent for further processing order to obtain a so-called ``concentrate'', which contains tens of times as little gangue.


Forces of cohesion between surfaces are capable of violating the levelling of a liquid in communicating vessels. It is very easy to verify the truth of this.

\begin{figure}[!ht]
\centering
\includegraphics[width=0.6\textwidth]{figures/fig-02-04.pdf}
\caption{The capillary action is the principle behind transport of sap in plants.}
\label{fig-2.4}
\end{figure}

If a thin glass tube (with a diameter of a fraction of a millimetre) is lowered into water, then in violation of the law of communicating vessels, the water in it will quickly begin rising, and its level will become consid­erably higher than in the large vessel (\figr{fig-2.4}).

But what took place? What forces are supporting the weight of the column of liquid that has risen up? The rise is accomplished by the forces of cohesion between the water and the glass.

Forces of cohesion between surfaces clearly manifest themselves only when a liquid rises in sufficiently thin tube. The narrower the tube, the higher the liquid and the more distinct the phenomenon. The name of these surface phenomena is related to the name of the tubes. The inside diameter of such tube is measured in fractions of a millimetre, such tube is called capillary (meaning ``thinner than a hair''. The phenomenon of the rise of liquids in thin tubes is called capillarity.\label{capillary-force}

But to what height are capillary tubes capable of raising a liquid? It turns out that water rises to a height of \SI{1.5}{\milli\meter} in a tube of \SI{1}{\milli\meter} diameter. For a diameter of \SI{0.01}{\milli\meter}, the height of the rise will increase as many
times as the diameter of the tube decreases, i.e. to \SI{15}{\centi\meter}.

Of course, the elevation of a liquid is possible only in the case of wetting. It is not hard to guess that mercury not rise in glass tubes. On the contrary, mercury falls in glass tubes. Mercury is so ``intolerant'' of contact with glass that it strives to reduce the total surface to the minimum allowed by gravity.

There exist many bodies which something like systems of very thin tubes. Capillary phenomena can always be observed in such bodies.


\begin{figure}[!ht]
\centering
\includegraphics[width=0.6\textwidth]{figures/fig-02-05.pdf}
\caption{A siphon works on the basis of capillary action.}
\label{fig-2.5}
\end{figure}

Plants and trees have an entire system of long ducts and pores. The diameters of these ducts are less than hundredths of a millimetre. Because of this, capillary forces raise soil moisture to a considerable height and distribute water throughout the plant.

Blotting paper is a very convenient thing. You spilled some ink on a page and want to turn it over. But you’re not going to wait until the blot dries up. You take a sheet of blotting paper, dip one of its edges in the drop, and the ink swiftly runs upwards against gravity.

A typical capillary phenomenon has occurred. If you look at the blotting paper through a microscope, you can see its structure. Such a paper consists of a sparse network of paper fibres forming thin and long ducts with each other. These ducts play the role of capillary tubes.

The same kind of system of long pores or ducts formed by fibres exists in wicks. Kerosene rises through the wick of a lamp. A siphon can also be created with the aid of a wick by placing one of its ends in a glass partially filled with water, in such a way that the other end hanging over the edge is lower than the first (\figr{fig-2.5}).

In the technology of the dyeing industry, frequent use is also made of the ability of a fabric to draw in a liquid through the thin pores formed by the threads of the fabric.

We have not said anything yet about the molecular mechanism of these interesting phenomena.

The differences in the surface tension of various substances can be excellently explained by intermolecular interactions.

A drop of mercury does not spread out over the surface of glass. The reason is that the energy of interaction between the mercury atoms is greater than the energy of cohesion between the glass and mercury atoms. For the same reason, mercury does not rise in narrow capillary tubes.

It is a different matter with water. It was found that the atoms of hydrogen of the water molecules readily cohere to the atoms of oxygen in the silica of which glass mainly consists. The water-glass intermolecular forces are greater than the water-water intermolecular forces. Therefore, water spreads out into a thin film over the surface of glass and rises in glass capillary tubes.

Surface tension or, more precisely, the energy of cohe­sion (depth of the potential well in \figr{fig-2.1}) for various pairs of substances can be both measured and calculated. A discussion on how this can be done would lead us too far away from our subject.


\section{Crystals and Their Shape}

Many people think that crystals are beautiful, rarely found stones. They occur in various colours, are usually transparent and, what is most remarkable, possess a beautiful regular shape. Crystals are most often polyhedra with ideally plane sides (faces) and strictly straight edges. They please the eye with a marvellous play of colours at the faces and an amazingly regular structure.

Among them are the unassuming crystals of rock salt natural sodium chloride, i.e. common salt. They are found in nature in the form of rectangular parallelepipeds or cubes. Calcite crystals also have a simple form -- trans­parent oblique-angled parallelepipeds. Quartz crystals are much more complicated. Each little crystal has a great many faces of different shapes, intersecting in edges of different lengths.

However, a crystal is anything but a museum-piece. Crystals surround us everywhere. The solid bodies with which we build homes and make machines, the substances which we use in our daily lives almost all of them are crystals. But why do we not see this? The reason is that in nature we rarely come across bodies in the form of single individual crystals (or, as is said, monocrystals). Substances are most often found in the form of firmly linked crystalline grains of very small size, less than a thousandth of a millimetre. Such a structure can only be seen through a microscope.

Bodies consisting of tiny crystalline grains are called polycrystaltine (derived from the Creek \emph{polys} meaning ``many'').

Of course, polycrystalline  bodies must also be included among the crystals. It will then turn out that almost all the solid bodies surrounding us are crystals. Sand and granite, copper and iron, the salol sold in a drugstore and paint all these are crystals.

There are exceptions; glass and plastics do not con­sist of small crystals. Such solid-bodies are called amorphous.

Thus, studying crystals means studying almost all the bodies surrounding us. It is obvious that this is important. 

Single crystals are recognized at once by the regularity of their shapes: Plane faces and straight edges are char­acteristic properties of a crystal; the regularity of shape is undoubtedly related to the regularity of the internal structure of a crystal. If a crystal has been especially stretched in a certain direction, it means that the struc­ture of the crystal in this direction is also special in some way.

But imagine that a ball has been made by machine out of a large crystal. Will we succeed in figuring out that we have a crystal in our hands, in distinguishing this ball from a glass ball? Since the different faces of a crystal are developed to different degrees, this suggests that the physical properties of a crystal also differ in different directions. This, of course, refers to the strength of a crystal, its electrical conductivity and to many other properties. This peculiarity of a crystal is called the anisotropy of its properties. Anisotropic means different properties in different directions.

Crystals are anisotropic. On the contrary, amorphous bodies, liquids and gases are isotropic, i.e. possess iden­tical (derived from the Greek \emph{isos} meaning ``equal'') properties in different directions (derived from the Greek \emph{tropos} meaning ``turning''). The anisotropy of the proper­ ties of a crystal is precisely what permits us to find out whether or not a transparent, formless piece of matter
is a crystal.

Let us visit a mineralogical museum and closely exam­ine various monocrystalline specimens of crystals of the same substance. It is quite possible that specimens of both regular and irregular shapes are on display. Some of the crystals resemble fragments, others have one or two abnormally developed faces.

From the collection of specimens we select the ones that seem to be of perfect shape and sketch them. The drawing we obtain is shown in \figr{fig-2.6}. Quartz has again been taken as an example. As the crystals of other substances, quartz can develop a different number of faces of each ``kind'', as well as a different number of ``kinds'' of faces. Even if their similarity in appearance is not self-evident, such small crystals do resemble one another as do close relatives or, sometimes, as twins. In what does their similarity consist?

\begin{figure}[!ht]
\centering
\includegraphics[width=0.6\textwidth]{figures/fig-02-06.pdf}
\caption{Different shapes of crystals.}
\label{fig-2.6}
\end{figure}

Look carefully at \figr{fig-2.6}, which illustrates a number of quartz crystals. All of them are close ``relatives'' They can be made truly similar by grinding off various faces by different amounts in such a way that they remain parallel to their initial positions. It is readily evident, for instance, that crystal $II$ can be made, in this way, entirely similar to crystal $I$. This is possible because the angles are equal between like faces, for example, be­tween faces $A$ and $B$, $B$ and $C$, etc.

It is this equality of angles that gives the crystals their ``family'' resemblance. As the faces are ground off parallel to their initial positions, the shape of the crystal is changed, but the angles between the faces are retained.

As a crystal grows, certain of its faces, due to various random factors, may be in more favourable, and others in less favourable conditions for the addition of new layers of atoms or molecules. Crystals grown under different conditions may not appreciably resemble one another, as far as their appearance is concerned, but the angles be­tween the corresponding faces of all crystals of the substance being considered are always the same. The shape (or \emph{habit}, as it is called) of a crystal is a matter of chance, but the angles between its faces (and you will understand why further on) depend upon its internal structure.

The flatness of its faces is not the only feature that distinguishes a crystal from shapeless bodies. A crystal also possesses symmetry. The meaning of this word, in the sense we are about to employ, can best be demonstrat­ed by examples.

\begin{figure}[!ht]
\centering
\includegraphics[width=\textwidth]{figures/fig-02-07.pdf}
\caption{Mirror symmetry.}
\label{fig-2.7}
\end{figure}


Illustrated in \figr{fig-2.7} is a statue (reminding us of the famous ones on Easter Island) standing in front of a large mirror. The reflection in the mirror is an exact copy (actually a mirror image) of the statue. The sculptor could have carved two statues and arranged them in the same manner as our statue and its reflection in the mirror.

This double sculpture is a symmetrical figure; it consists of two equal parts, one being a mirror image of the other and equal in the sense that one’s left hand is equivalent to one’s right.

Let us assume that a flat mirror is placed as in \figr{fig-2.7} Then the right-hand part of the sculpture exactly coincides with the reflection of its left-hand part. Such symmetrical arrangement would have a vertical plane of mirror symmetry passing through the centre between the two parts. The plane of symmetry is an imaginary one, but we sense it distinctly in examining a symmetrical body.

The bodies of animals have planes of symmetry; a ver­tical plane of external symmetry can be passed through a man as well. Symmetry is only approximate in the animal world, and, in general, there is no ideal symmetry in the world around us. An architector can design a house consisting of two ideally symmetrical halves. But when the house is built, no matter how well, you can always find some difference in the two symmetrical parts: one may have a crack at some place where the other has not.

The most exact symmetry is found in the world of crystals, though it is hardly ideal here either. Fissures invisible to the naked eye, scratches and other flaws always make equal faces differ slightly from one another.

The child’s toy called a pinwheel is illustrated in \figr{fig-2.8}. It is also symmetrical, but you cannot pass a plane of symmetry through it in any way. What then is symmetrical about this toy? First, let us consider its symmetrical parts. How many are there? Obviously, four. What makes their mutual arrangement regular? This, too, is simple enough. Let us turn the pinwheel counterclockwise through a right angle, i.e. one-fourth of a revolution. Now vane / has turned to the previous position of vane 2, 2 to that of vane 3, 3 to that of vane 4, and 4 to that of vane 1. The new position of the pinwheel cannot be distinguished from the previous one. We say that such a figure has an axis of symmetry or, more pre­cisely, an axis of four-fold symmetry because it coincides with itself after turning through one-fourth of a revolu­tion.

\begin{figure}[!ht]
\centering
\includegraphics[width=0.6\textwidth]{figures/fig-02-08.pdf}
\caption{A pinwheel has four-fold symmetry.}
\label{fig-2.8}
\end{figure}



Thus, an axis of rotation symmetry is an imaginary straight line about which a body, when turned through a unit fraction of a revolution, is in a position that can­ not be distinguished from its initial position. The order of rotation symmetry (four-fold in our case) indicates that such coincidence occurs after turning the body one-fourth of a revolution. Consequently, after four such turning motions we return to the initial position.

Do we find symmetry of all kinds in the crystal king­dom? Investigations have shown that we do not. In fact, the only axes of rotation symmetry found in crystals are the two-, three-, four- and six-fold ones. This is no mere chance. Crystallographers have shown that it is related with the internal structure of crystals. Therefore, the number of different kinds or, as we say, classes of crystal symmetry is relatively small -- only 32.

\section{Structure of Crystals}

Why is the regular shape of a crystal so beautiful? Its faces, so smooth and shining, look as though they were polished by a skilled lapidary. Different parts of the crystal repeat one another, forming a handsome sym­metrical figure. This exceptional regularity of crystals has been known since time immemorial. But the con­cepts of the ancient scholars on the nature of crystals differed only slightly from the tales and legends made up by poets whose imagination was fascinated by the beauty and elegance of crystals. It was believed that rock crystal is formed of ice, and diamonds of rock crystal. Many wonderous properties were attributed to crystals; they were thought to heal diseases, protect one against poisons, have an influence on one’s fate and many others.

\begin{figure}[!ht]
\centering
\includegraphics[width=\textwidth]{figures/fig-02-09.pdf}
\caption{The structure of crystals made up of tiny ``building blocks''.}
\label{fig-2.9}
\end{figure}



The first scientific views on the nature of crystals appeared only in the $17^{\textrm{th}}$ and $18^{\textrm{th}}$ centuries. An idea of these conceptions is given by \figr{fig-2.9}, which was reproduced from \emph{Traite de Mineralogie}, written by the French abbe, Rene Just Ha\"uy, near the end of the 18th century. 

According to its author, crystals are made up of tiny ``building blocks'', fitting tightly to one another. This conclusion is one that is naturally arrived at. Let us break up a crystal of Iceland spar (calcite, or calcium carbonate) with a sharp blow of a hammer. It flies apart into pieces of various sizes. Examining them carefully, we find that the pieces are of regular shape, quite similar to that of the large crystal, their ``parent'' Evidently, reasoned Ha\"uy, if we continue to break up the pieces into smaller and smaller ones, we will finally reach the smallest building block, invisible to the naked eye, which is a crystal of the given substance. These ultimate blocks are so small that steps formed by them, comprising the faces of the large crystal, seem to us to be immaculately smooth. Well, and what is this final building block like? The scientists of that day could not answer this question.

The ``building block'' theory of crystal structure was of great benefit to science. It explained the origin of the straight edges and flat faces of a crystal. As a crystal grows, new building blocks attach themselves to those of the crystal, and a face grows like the wall of a house built by the hands of bricklayers.

Hence, the question concerning the reason for the regularity and beauty of the shapes of crystals has been answered a long time ago. The reason for this phenom­enon is their internal regularity. This regularity consists in the endless repetition of one and the same ele­mentary parts.

Imagine a park fence made of rods of different lengths and distributed helter-skelter. An ugly scene. A good fence is constructed from identical rods distributed in a regular sequence at equal distances from each other. We find such a self-repeating pattern on wallpaper. Here an element of a drawing, say, a girl playing with a ball, is repeated not only in one direction, as in a park fence, but so that it fills a plane.

\begin{figure}[!ht]
\centering
\includegraphics[width=0.6\textwidth]{figures/fig-02-10.pdf}
\caption{Finding a ``unit cell'' in a wallpaper pattern.}
\label{fig-2.10}
\end{figure}

But what relationship do a park fence and wallpaper have to a crystal? A most direct one. A park fence consists of links repeating along a line, wallpaper of pat­terns repeating in a plane, and a crystal of groups of atoms repeating in space. One therefore says that the atoms of a crystal form a space (or crystal) lattice.

There are certain details concerning a three-dimensional, or space, lattice that we must now discuss. To sim­plify the work of the illustrator, we shall explain all that is required, using wallpaper as an example, rather than asking him to construct complicated drawings of three-dimensional figures.

Separated out in \figr{fig-2.10} is the smallest piece that can be transferred or shifted to obtain the whole pattern of the wallpaper. To separate out such a piece, we draw two straight lines from any point of the drawing, for instance, the centre of the beach ball, connecting it to the same points of the two adjacent balls. As can be seen in our drawing, we can use these two lines to construct a parallelogram. By shifting or transferring this parallelogram along the directions of the two initial basic lines over distances equal to its sides, we can obtain the whole wallpaper pattern. This smallest repeated portion, commonly called a unit cell, can be selected in various ways. It is evident from \figr{fig-2.10} that several different parallelograms could be selected, each containing one drawing. We emphasize that in the given case it makes no difference to us whether the picture inside the cell is a whole one or it is divided up by the lines bounding the cell.

It would be a mistake to suppose that after drawing the repeating picture, called the motif, for the wallpaper the artist can always consider his job to be finished. This would be true if the only way to make a pattern for wallpaper was to add to the given portion another identi­cal portion, shifted parallel to the first along the two initial basic lines.

\begin{figure}[!ht]
\centering
\includegraphics[width=\textwidth]{figures/fig-02-11.pdf}
\caption{Seventeen possible ways to fill the wallpaper design, that is filling the plane.}
\label{fig-2.11}
\end{figure}


In addition to this simplest method, however, there are sixteen more ways to fill in the wallpaper design with an orderly repeated drawing, or motif, i.e. seventeen types of mutual arrangement of drawings on a plane. They are illustrated in \figr{fig-2.11}. Here a simpler repeating motif has been selected, but, like the one in \figr{fig-2.10}, it has no symmetry in itself. The patterns composed of this motif are symmetrical and the difference in the patterns is due to the difference in the symmetrical arrangements of the motifs.

We see, for example, that in the first three cases the pattern has no plane of mirror symmetry: you cannot position vertical mirror so that one part of the pattern is a mirror reflection of another part. Cases 4 and 5, on the contrary, have planes of symmetry. Two mutually perpendicular mirrors can be set up in cases 8 and 9.

Case 10 has an axis of four-fold rotation symmetry, per­pendicular to the plane of the drawing. In case 11, the axis is of three-fold symmetry, and in cases 13 and 15, of six-fold rotation symmetry.

Planes and axes of symmetry are found on our drawings in parallel series or families, rather than singly. If you have found one point through which you can pass an axis (or plane) of symmetry, you can readily find the adjacent point and consecutive points, all spaced at the same distance from one another, through which the same kind of axes (or planes) of symmetry can be passed.


\begin{figure}[!ht]
\centering
\includegraphics[width=\textwidth]{figures/fig-02-12.pdf}
\caption{Two illustrations for Case 8 type of pattern.}
\label{fig-2.12}
\end{figure}

The seventeen types of symmetry of a plane pattern do not, of course, exhaust all of the great variety of patterns that can be composed from a single motif. The artist must specify one more condition: how the motif is to be positioned with respect to the boundary lines of the cell. \figr{fig-2.12} illustrates two wallpaper pat­terns with the same initial motif which is differently positioned in each pattern with respect to the lines rep­ resenting the vertical mirrors (in a plane pattern they are simply lines of symmetry). Both patterns come under case 8 in \figr{fig-2.11}.

Each body, including crystals, consists of atoms. Simple substances, or elements, consist of atoms of a single kind. Complex substances, or compounds, con­sist of atoms of two or several kinds. Assume that in some superpowerful microscope we could examine the surface of a crystal of common salt and see the centres of the atoms. \figr{fig-2.13} shows how we would see the atoms arranged along a face of the crystal like a wallpaper pattern. You can readily understand now how a crystal is constructed. We could say that a crystal resembles ``three-dimensional wallpaper''. A crystal is composed by joining three-dimensional unit cells tightly together instead of plane cells as for wallpaper patterns. These three-dimensional unit cells are the building blocks of which the crystals are built.

\begin{figure}[!ht]
\centering
\includegraphics[width=\textwidth]{figures/fig-02-13.pdf}
\caption{A crystal as a ``three-dimensional wallpaper''.}
\label{fig-2.13}
\end{figure}


How many different ways exist for building up three- dimensional wallpaper patterns from these elementary pieces? This complex mathematical problem was solved at the turn of the century by the founder of structural crystallography Yevgraf Stepanovich Fedorov (1853- 1919). This eminent Russian scientist established the fact that there can be only 230 different ways to build up a crystal or, as they say today, 230 Fedorov groups.

All that is known today about the internal structure of crystals was obtained by $X$-ray structure analysis, which we shall describe in some detail in the fourth book.

Simple crystals exist, made up of atoms of a single kind. Diamond, for example, is of pure carbon. Crystals of common salt consist of ions (electrically charged atoms) of two kinds: sodium and chlorine. More complex crystals may be made up of molecules which, in their turn, consist of atoms of many kinds.

We can always single out in a crystal the smallest repeating groups of atoms (or single atoms, in the simplest case). Such a group is called the unit cell. \label{unit-cell-def}

The dimensions of unit cells may vary in wide ranges. The shortest distances are found between the adjacent lattice points (corners of the unit cell) of the simplest crystals, consisting of atoms of a single kind. The maximum distances are found in complex protein crystals. These distances, called lattice constants, range from 2 or 3 angstroms to several hundred angstroms (hund­redths of one-millionth of a centimetre).

There are a great variety of crystal lattices. The prop­erties common to all crystals are due to their lattice structure. To begin with, we can readily understand that the ideally flat faces of a crystal are the planes passing through the lattice points at which atoms are located. But any number of lattice planes can be passed in the most diverse directions. Which of these planes bound the grown crystals, i.e. become its faces?

Let us turn our attention, first of all, to the following circumstance: various lattice planes and lines are not filled equally densely with lattice points. Experiments show that a crystal is faceted by the planes that are most densely studded with lattice points and that these planes intersect at edges which are also most densely occupied by lattice points.

\begin{figure}[!ht]
\centering
\includegraphics[width=0.6\textwidth]{figures/fig-02-14.pdf}
\caption{A view of a crystal perpendicular to a face.}
\label{fig-2.14}
\end{figure}

A view of a crystal lattice perpendicular to a face is shown in \figr{fig-2.14}. Shown also are the traces of lattice planes perpendicular to the drawing. It should be clear from what has been said that faces on the crystal parallel to planes $I$ and $III$ are apt to develop. No faces will be developed parallel to plane $II$ which is so sparsely dotted with lattice points (and atoms).

The structures of many hundreds of crystals are known today. We shall discuss the structures of the simplest crys­tals, first of all, those built up of atoms of a single kind.

Three types of lattices are most common. They are shown in \figr{fig-2.15}. The centres of the atoms are repre­sented by points; the lines joining the points do not have any actual meaning. They have been drawn merely to make the nature of the spatial distribution of the atoms clearer to the reader.

\figr{fig-2.15}~\textcolor{black!70}{(a)} and \textcolor{black!70}{(b)} depict cubic lattices. In order to visualize these lattices more clearly, imagine that you have arranged building blocks in the simplest manner -- edge to edge and face to face. If you now conceptually place points at the vertices and centres of the cubes, the cubic lattice depicted in \figr{fig-2.15}~\textcolor{black!70}{(a)} appears. Such a structure is called body-centred cubic. If points are placed at the vertices of the cubes and at the centres of their faces, the cubic lattice depicted in \figr{fig-2.15}~\textcolor{black!70}{(b)} appears. It is called face-centred cubic.


\begin{figure}[!ht]
\centering
\includegraphics[width=\textwidth]{figures/fig-02-15.pdf}
\caption{Three types of common lattices.}
\label{fig-2.15}
\end{figure}

The third lattice (\figr{fig-2.15}~\textcolor{black!70}{(c)}) is called close-packed hexagonal (i.e. having six angles). In order to understand the origin of this term and more clearly visualize the distribution of the atoms in this lattice, let us take some billiard balls and start packing them as closely as pos­sible. First of all, let us form a close layer -- it looks like billiard balls which have been gathered in by a rack before the beginning of a game (\figr{fig-2.16}). Note that the ball in the middle of the triangle is in contact with six neighbours, and these six neighbours form a hexagon.

We continue the packing by laying one layer upon ano­ther. If we place the balls of the second layer directly above the balls of the first, such a packing would not be close. Trying to distribute the greatest number of balls in a definite volume, we should place the balls of the second layer in the holes formed by the first, the balls of the third layer in the holes of the second, etc. In a close-packed hexagonal lattice, the balls of the third layer are placed in such a way that their centres lie directly above those of the balls of the first.

\begin{figure}[!ht]
\centering
\includegraphics[width=\textwidth]{figures/fig-02-16.pdf}
\caption{Two examples of hexagonal lattice: a rack of billiard balls, the structure of a beehive.}
\label{fig-2.16}
\end{figure}

The centres of the atoms in a close-packed hexagonal lattice are distributed just like those of the balls which
are closely packed in the manner described.

A great many elements crystallize in the lattices of the
three types described above:

\begin{center}
\begin{small}
\begin{tabular}{p{5.25cm}p{4cm}}
\toprule
Close-packed hexagonal lattice & \ce{Be}, \ce{Co}, \ce{Hf}, \ce{Ti}, \ce{Zn}, \ce{Zr} \\
Face-centred cubic & \ce{Al}, \ce{Cu}, \ce{Co}, \ce{Fe}, \ce{Au}, \ce{Ge}, \ce{Ni}, \ce{Ti} \\
Body-centred cubic & \ce{Cr}, \ce{Fe}, \ce{Li}, \ce{Mo}, \ce{Ta}, \ce{Ti}, \ce{U}, \ce{V}\\
\bottomrule
\end{tabular}
\end{small}
\end{center}

We shall mention only a few of the other structures. The structure of diamond is depicted in \figr{fig-2.17}. What is characteristic of this structure is that a carbon atom in diamond has four immediate neighbours. Let us compare this number with the corresponding numbers for the three most common structures just described. As is evident from the figures, each atom has 12 immediate neighbours in a close-packed hexagonal lattice, the atoms forming a face-centred cubic lattice have just as many neighbours, and each atom has eight neighbours in a body-centred lattice.

\begin{figure}[!ht]
\centering
\includegraphics[width=0.5\textwidth]{figures/fig-02-17.pdf}
\caption{Crystal structure of a diamond.}
\label{fig-2.17}
\end{figure}

We shall say a few words about graphite, whose struc­ture is shown in \figr{fig-2.18}. This structure has a striking peculiarity. Graphite consists of layers of atoms, where atoms of a single layer are more firmly bound to each other than atoms of neighbouring layers. This is related to the values of the interatomic distances: the distance between neighbours in a single layer is 2.5 times as small as the shortest distance between layers.

\begin{figure}[!ht]
\centering
\includegraphics[width=0.5\textwidth]{figures/fig-02-18.pdf}
\caption{Crystal structure of graphite.}
\label{fig-2.18}
\end{figure}

The presence of weakly bound atomic layers makes it easy for graphite crystals to break up along these layers. This is why solid graphite can serve as a lubricant in those cases when it is impossible to apply lubricating oil -- for example, at very low or very high temperatures. Graphite is a solid lubricant.

Friction between two bodies reduces, roughly speaking, to the fact that microscopic protuberances of one body fall in hollows of the other. The force required for breaking up a microscopic graphite crystal is much smaller than the frictional forces; therefore, a graphite lubrication considerably facilitates the sliding of one body along another.

There is an endless variety in the structures of crystals of chemical compounds. The structures of rock salt and carbon dioxide depicted in Figures 2.19 and 2.20 can serve as extreme examples.

\begin{figure}[!ht]
\centering
\includegraphics[width=0.5\textwidth]{figures/fig-02-19.pdf}
\caption{Crystals of rock salt.}
\label{fig-2.19}
\end{figure}

Crystals of rock salt (\figr{fig-2.19}) consist of atoms of sodium (small dark balls) and chlorine (large light balls) alternating along the axes of a cube. Each sodium atom has six equidistant neighbours of the other kind. The same is also true of chlorine. But where are the molecules of sodium chloride? There are none; not only are groups of single atoms of sodium and single atoms of chlorine absent from salt crystals, but in general, no group of atoms whatsoever can be distinguished from the others by their proximity.

The chemical formula of NaCl does not give us any
grounds for saying that ``this substance is built up out of molecules of \ce{NaCl}''. The chemical formula merely indicates that the substance is constructed from the same number of atoms of sodium and chlorine.

\begin{figure}[!ht]
\centering
\includegraphics[width=0.5\textwidth]{figures/fig-02-20.pdf}
\caption{Crystal structure of carbon dioxide \ce{CO2}.}
\label{fig-2.20}
\end{figure}

The question of the existence of molecules in a substance is decided by its structure. If no groups of close atoms can be distinguished in it, there are no molecules.


A crystal of carbon dioxide, \ce{CO2} (the dry ice which lies in cartons of ice cream), is an example of a molecular crystal \figr{fig-2.20}.

The centres of the oxygen and carbon atoms of a molecule of \ce{CO2} are situated along a straight line (see \figr{fig-2.2}­). The distance \ce{C-O} is equal to \SI{1.3}{\angstrom}, and the distance between oxygen atoms of neighbouring molecules is about \SI{3}{\angstrom}. It is obvious that under such conditions we immediately ``recognize'' a molecule in a crystal.

Molecular crystals are close packings of molecules. In order to see this, we must outline the contours of the molecules. This is precisely what has been done in \figr{fig-2.20}.


\begin{figure}[!ht]
\centering
\includegraphics[width=\textwidth]{figures/fig-02-21.pdf}
\caption{Crystal structure of an organic substance.}
\label{fig-2.21}
\end{figure}

All organic substances consist of molecular crystals. Organic molecules frequently consist of many tens and even hundreds of atoms (those made up of tens of thousands of atoms are to be discussed in a separate chapter). It is impossible to properly depict their packing arrangements graphically. For this reason, you may find drawings similar to \figr{fig-2.21} in books in this field. The mole­cules of this organic substance are built up of carbon atoms. The bars between the atoms represent the valence bonds. The molecules seem to be suspended in air. But do not believe what you see here. They have been drawn thus to give you an idea of how the molecules are arranged in a crystal. For the sake of simplicity, the illustrator did not show the hydrogen atoms joined to the outer atoms of carbon (as a matter of fact, chemists frequently make their similar omissions). He did not even consider it necessary to outline the molecule, imparting a definite shape to it. If he had, we would see that the principle of molecular packing, with the ``key fitting the lock'', is just as valid in this case as in other similar ones.



\section{Polycrystalline Substances}

We have already mentioned the fact that amorphous bodies are rare in the world of solids. Most of the objects surrounding us consist of tiny crystalline grains, about a thousandth of a millimetre in size.

Investigators discovered the granular structure of met­als in the $19^{\textrm{th}}$ century with the aid of ordinary optical microscopes. All that was required was an arrangement enabling the specimens to be observed in reflected instead of transmitted light. Up-to-date metallurgical microscopes operate on this same principle.

The image seen in such a microscope may be like that shown in \figr{fig-2.22}. The grain boundaries are usually quite distinct. As a rule, impurities accumulate at the boundaries.

\begin{figure}[!ht]
\centering
\includegraphics[width=0.6\textwidth]{figures/fig-02-22.pdf}
\caption{Crystal structure of a metal.}
\label{fig-2.22}
\end{figure}

The properties of materials depend to an exceptionally large degree on the size of the grains, their orientation and what is happening at their boundaries. For this reason, physicists have devoted much research to the study of polycrystalline substances. By means of $X$-ray structure analysis, about which we have already promised to tell our readers, they found that each grain is a small crystal. We repeat our promise.

All treatment of a metal affects its grains. Suppose we have a piece of cast metal. Its grains are in disorder and their size is quite large. Then we make wire of the piece of metal by drawing it through a die. What happens to the crystalline grains in this plastic working operation? Investigations have shown that the change in shape of a solid in drawing wire or in some other plastic working technique breaks up the crystalline grains. At the same time, a certain element of order is produced in the arrange­ment of these grains by the action of the mechanical forces. But what possible kind of order can there be here? The fragments of the grains are absolutely shapeless.

This is true: the fragments may be of any random shape, but a fragment of a crystal is still a crystal and the atoms in its lattice are packed as regularly as in any well-faceted crystal. Hence, we can indicate in each fragment how its unit cells are oriented. Prior to plastic working, the cells are strictly ordered only within each separate grain; there is, as a rule, no total order. After the metal is worked, the grains align themselves so that a certain general order becomes evident in their cells. This is called \emph{texture}. For example, the diagonals of the cells in all the grains become oriented approximately parallel to the direction in which the metal is worked.

\begin{figure}[!ht]
\centering
\includegraphics[width=0.6\textwidth]{figures/fig-02-23.pdf}
\caption{\emph{Texture} in a crystal.}
\label{fig-2.23}
\end{figure}

Texture is illustrated in \figr{fig-2.23} (at the right) by the example of certain definite, indicated planes in the grains. These planes are most densely filled with atoms and are shown by rows of dots. The illustration at the left shows the grains of the metal before being worked.

The various kinds of plastic working operations (roll­ing, forging, and wire-drawing) produce textures of different types. In some operations, the grains turn so that the diagonals of their unit cells are aligned along the direc­tion of working, while in others the edges of the cells are aligned, etc. The more advanced the rolling or wire­ drawing technique, the more perfect the texture of the crystalline grains of metal. Texture has a striking effect on the mechanical properties of metal articles. A study of the arrangement and size of the crystalline grains in metals led to an understanding of the fundamental prin­ciples involved in the plastic working of metals. This, in turn, led to the improvements of the techniques em­ployed.

Another important kind of metal treatment, anneal­ing, is also associated with the rearrangement of the metal grains. If rolled or drawn metal is heated to a suffi­ciently high temperature, new crystals are formed, which grow at the expense of the old ones. Annealing gradually destroys the texture; the new crystals are arranged dis­orderly. As the temperature is raised (or the metal is held longer at the annealing temperature), new grains grow and the old ones disappear. The grains can grow to a size visible to the naked eye. Annealing drastically alters the properties of a metal. It becomes more ductile and softer. This occurs because the grains become coarser and the texture disappears.


%\newpage
%\begin{center}
%\includegraphics[width=\textwidth]{figures/helmholtz.pdf}
%\end{center}
%{\small \textbf{Hermann Helmholtz (1821-1894)} -- a famous Herman scientist. Helm­holtz worked in hie fields of physics, mathematics and physiology with great success. He was the first (1847) to give a mathematical interpretation of the law of conservation of energy emphasizing the universal character of this law. Helmholtz obtained outstanding \textcolor{red}{\textbf{incomplete}}}
%
%







% !TEX root = pfe-book2.tex
%!TEX TS-program = pdflatex
%!TEX encoding = UTF-8 Unicode


\cleardoublepage
%\mainmatter
\chapter{Temperature}
\label{ch-03}

\section{Thermometer}
If two differently heated bodies are brought into contact, the warmer one will cool off and the colder one will warm up. It is said that two such bodies exchange heat.

As we have already said, heat exchange is a kind of energy transfer; the body which gives off energy is said to be hotter. We feel that a body is hot if it makes our hand warm, transfers energy to it. On the contrary, if a body is felt to be cold, this means that it is taking energy away from our body.

Concerning a body which is giving off heat (i.e. giving off energy by means of heat exchange), we say: its temper­ature is higher than that of the body which is taking in this heat.

Observing whether an object of interest to us is cooling off or warming up in the presence of one or another body, we find ``its place'' in a row of heated bodies. Temperature is a kind of mark indicating for which bodies the object of interest to us will be a giver, and for which a receiver, of heat.

Temperature is measured by thermometers.

The structure of thermometers can be based on the utilization of various properties of bodies sensitive to temperature. The most frequently utilized property is the expansion of bodies during a rise in temperature.

If the body of a thermometer changes its volume when in contact with various bodies, this implies that these bodies have different temperatures. When the volume of the body of a thermometer is greater, the temperature is higher, and when the volume is smaller, the temperature is lower.

There are various substances that can serve as thermom­eters: liquids (such as mercury or alcohol), solids (such as metals), and gases. But different substances expand differently, and so mercurial, alcoholic, gaseous and other ``degrees'' will not coincide. Of course, two basic points, the melting point of ice and the boiling point of water, can always be marked on all thermometers. Therefore, all thermometers will indicate 0 and 100 degrees centi­grade identically. But bodies will not expand identically between 0 and 100 degrees. One body expands rapidly between 0 and 50 degrees on a mercury thermometer, and slowly in the second part of this interval, but another, vice versa.

Having made thermometers with differently expand­ing bodies, we will discover noticeable discrepancies in their readings, in spite of the fact that their readings will coincide for the basic points. Moreover, a water thermometer would lead us to the following discovery: if a body cooled to zero is placed on an electric stove, its ``water temperature'' would first fall and then rise. This happens because water at first decreases its volume when heated, and only later behaves ``normally'', i.e. increases its volume when heated.

We see that a rash choice of material for a thermometer can bring us to an impasse.

But then what should we be guided by in choosing a ``true'' thermometer? Which body would be ideal for this purpose?

We have already spoken about such ideal bodies. They are ideal gases. There are no interactions between the particles an ideal gas and studying the expansion of an ideal gas, we study how the motion of its molecules changes. This is precisely the reason why an ideal gas is an ideal body for a thermometer.

And it really is a striking fact that while water does not expand like alcohol (nor alcohol like glass, nor glass like iron), then hydrogen, oxygen, nitrogen and any other gas in a sufficiently rarefied state to deserve being called ideal expand in exactly the same fashion when heated.

Therefore, the changes in volume undergone by a def­inite amount of ideal gas serve as the basis for defining temperature in physics. Of course, in view of the fact that gases are highly compressible, one must be especially careful in seeing to it that the gas is under constant pres­sure.

In order to graduate a gas thermometer, we should accurately measure the volume of the gas we have chosen at \SI{0}{\celsius} and at \SI{100}{\celsius}. We shall divide the difference be­tween the volumes $V_{100}$ and $V_{0}$ into 100 equal parts. In other words, the change in the volume of the gas by $(V_{100} - V_{0})/100$ corresponds to one degree centigrade (\SI{1}{\celsius}).

Let us now suppose that our thermometer shows a volume $V$ What temperature $t\si{\celsius}$  corresponds to this volume? It is not difficult to comprehend that
\begin{equation*}%
t \si{\celsius} = \frac{V - V_{0}}{V_{100}-V_{0}} 100 \,\, \text{i.e.} \,\, \frac{t\si{\celsius}}{100} = \frac{V - V_{0}}{V_{100}-V_{0}}
\end{equation*}
By means of this equality, we assign each volume $V$ to a temperature $t$ and obtain the temperature scale\footnote{The centigrade scale at which \SI{0}{\celsius} is taken as the melting point of ice, and \SI{100}{\celsius} as the boiling point of water (both at the standard pressure of \SI{760}{\milli\meter\mercury}) is very convenient. In spite of this, the British and the Americans have so far been using a temperature scale which seems very strange to us. How, for example, will you react to the following sentence taken from an English novel: ``The summer wasn’t hot, the temperature was 60-70 degrees.'' A misprint? No, the Fahrenheit scale (\si{\fahrenheit}). The temperature in England rarely falls below \SI{-20}{\celsius}. Fahrenheit selected a mixture of ice and salt having approximately such a temperature and took this temperature for his zero. In the words of the inventor, the standard temperature of a human body was taken for 100 on this scale. However, in order to determine this point, Fahrenheit probably made use of the services of a slightly feverish person. On the \emph{Fahrenheit scale}, the average standard temperature of a human body is \SI{98}{\fahrenheit}. On this scale, water freezes at \SI{+32}{\fahrenheit} and boils at \SI{212}{\fahrenheit}. The conversion formula will be: 
$t \si{\celsius}= \dfrac{5}{9} (t - 32) \si{\fahrenheit}$.}
which physicists use.

With a rise in temperature, the volume of a gas increases without bound -- there is no theoretical limit to the growth in temperature. On the contrary, low (negative on the centigrade scale) temperatures have a limit.

What will happen when the temperature is lowered? A real gas will eventually turn into a liquid, and with an even greater fall in temperature, will solidify. The gas molecules will gather in a small volume. But what will this volume be equal to for a thermometer filled with an ideal gas? Its molecules do not interact with each other and do not have any volume of their own. Hence, a decrease in temperature brings an ideal gas to a zero vol­ume. It is quite possible to come as close as we wish in practice to a behaviour that is characteristic of an ideal gas, for example, to a zero volume. For this it is necessary to fill up the gas thermometer with more and more rare­fied gas. Therefore, we won’t go wrong by assuming the minimum volume of the gas equal to zero.

According to our formula, the lowest possible tempera­ture corresponds to a zero volume. This temperature is called the \emph{absolute zero} of temperature.

In order to determine the position of absolute zero on the centigrade scale, we must substitute zero for the vol­ume $(V = 0)$ in the temperature formula just derived. Consequently, the temperature of absolute zero is equal to $-100V_{0}/(V_{100}- V_{0})$.

It turns out that this remarkable point corresponds to a temperature of about $\SI{-273}{\celsius}$ (more precisely, $\SI{-273.15}{\celsius}$). 

Thus, there are no temperatures below absolute zero; for they would correspond to negative volumes of a gas. It doesn’t make sense to speak of lower temperatures. It is just as impossible to obtain temperatures below abso­lute zero as to make a wire with a diameter less than zero. It is impossible to cool a body at absolute zero, i.e. one cannot take energy away from it. In other words, bodies and the particles they are made of have the least possible energy at absolute zero. This implies that the kinetic energy equals zero and the potential energy as­sumes its least possible value at absolute zero.

Since absolute zero is the lowest temperature, it is only natural that the \emph{absolute scale} in which readings begin at absolute zero be used in physics, especially in those of its branches where low temperatures play an important role. It is clear that $T_{\text{abs}}  = (t + 273)\,\, \si{celsius}$. Room temperature will be about \SI{300}{\celsius} on the absolute scale. The absolute scale is also called the \emph{Kelvin scale}, in honour of the well-known 19th century English scien­tist, and the notation $T$ \si{kelvin} is employed in place of $T_{\text{abs}}$. A formula for a gas thermometer determining the abso­lute temperature $T$ can be written down in the form 
\begin{equation*}%
T = 100 \left(  \frac{V - V_{0}}{V_{100}-V_{0}} \right) + 273
\end{equation*}
Using the equality $100V_{0}/(V_{100}-V_{0}) = 273$, we arrive at the following simple result:
\begin{equation*}%
\frac{T}{273} = \frac{V}{V_{0}}
\end{equation*}
Therefore, the absolute temperature is simply proportion­al to the volume of an ideal gas.

Exact measurements of temperature require all kinds of contrivances on the part of the physicist. Mercury, al­cohol (for Arctic regions) and other thermometers are graduated by comparison with a gas thermometer over a rather wide temperature interval. However, it too is also unsuitable for temperatures very close to absolute zero (below \SI{0.7}{\kelvin}), when all gases liquefy, and also for temperatures above \SI{600}{\celsius}, when gases penetrate glass. Other principles of temperature measurement are used for high and very low temperatures.

As for practical methods of measuring temperature, they are manifold. Instruments based on electrical phenomena are of great significance. It is now important to remember only one thing -- during all measurements of temperature, we should be convinced that the reading obtained completely coincides with what a measurement of the expansion of a rarefied gas would give.

High temperatures arise in ovens, furnaces and burners. Temperatures of \SIrange{220}{280}{\celsius} are attained in baking ovens. Higher temperatures are applied in metallurgy -- hardening furnaces yield \SIrange{900}{1000}{\celsius}, forges yield \SIrange{1400}{1500}{\celsius}. Temperatures of \SI{2000}{\celsius} are attained in steel furnaces.

The records for high temperature in a furnace are ob­tained with the aid of electric arcs (about \SI{5000}{\celsius}). The arc flame makes it possible to deal with the most refractory metals.

And what is the temperature of the flame of a gas burn­er? The temperature in the inner bluish cone of flame is only \SI{300}{\celsius}. The temperature in the outer cone attains \SI{1800}{\celsius}.

Incomparably higher temperatures arise during the explosion of an atomic bomb. Judging by indirect esti­mates, the temperature at the centre of the explosion attains several million degrees.

Attempts have recently been made to obtain such ultrahigh temperatures in special laboratory installations built in the Soviet Union and in other countries. It has proved possible to attain temperatures up to several million degrees for very brief moments.

Ultrahigh temperatures also exist in nature, not on the Earth, but on other bodies in the Universe. In the centres of stars, the Sun in particular, the temperature attains tens of millions of degrees. But the surfaces of stars have a considerably lower temperature, not exceeding \SI{20000}{\celsius}. The surface of the Sun gets heated up to \SI{6000}{\celsius}.

\section{Ideal Gas Theory}
The properties of an ideal gas, giving us the definition of temperature, are very simple. The Boyle law is valid for constant temperatures: during changes in volume or pres­sure, the product $pV$ remains constant. For a constant pressure, the quotient $V/T$ is conserved, no matter how the volume or temperature changes. It is easy to unite these two laws. It is clear that the expression $pV/T$ re­mains the same as for a constant temperature but with changing $V$ and $p$; the same is true for a constant pressure but with changing $V$ and $T$. The expression $pV/T$ re­mains constant during a change not only in any pair, but also simultaneously in all three of the quantities $p$, $V$ and $T$. The law $pV/T = \text{const}$ is known as the \emph{equation of state of an ideal gas}.

An ideal gas is chosen for a thermometer because its properties depend only on the motion (but not on the interaction) of its molecules.

What is the nature of the relationship between the notion of molecules and temperature? To answer this ques­tion, it is necessary to find the relationship between the pressure of a gas and the motion of the molecules in it.

\begin{figure}[!ht]
\centering
\includegraphics[width=0.4\textwidth]{figures/fig-03-01.pdf}
\caption{Finding pressure in a spherical vessel.}
\label{fig-3.1}
\end{figure}

In a spherical vessel of radius $R$, $N$ molecules of a gas are contained (\figr{fig-3.1}). Let us follow an arbitrary molecule, for example, one which is moving at a given moment from left to right along a chord of length $l$. We shall not pay attention to molecular collisions: such impacts do not affect the pressure. Having flown to the boundary of the vessel, the molecule will strike against the wall and fly off in some other direction with the same speed (the collision is elastic). Ideally, such a journey through the vessel might continue eternally. If $v$ is the molecular velocity, each impact will occur after $l/v$ seconds, i.e. each molecule will strike the wall $v/l$ times a second. The continual hail of impacts by the $N$ mole­cules unites into a single force of pressure.

According to Newton’s law, the force is equal to the change in momentum during a unit of time. Let us de­note the change in momentum at each impact by $\Delta$. This change occurs $v/l$ times a second. Consequently, the con­tribution to the force on the part of a single molecule will be $\Delta v/l$.

The momentum vectors before and after an impact, and also the momentum transfer $\Delta$ have been constructed in \figr{fig-3.1}. It follows from the similarity of the triangles arising in the construction that $\Delta/l = mv/R$. The con­tribution to the force on the part of a single molecule will take the following form:
\begin{equation*}%
\frac{mv^{2}}{R}
\end{equation*}
Since the length of the chord does not occur in the for­mula, it is clear that molecules moving along arbitrary chords make an identical contribution to the force. Of course, the change in momentum will be smaller for an oblique impact, but then the impacts in this case will be more frequent. Calculations show that these two effects exactly compensate for each other.

Since there are $N$ molecules in the sphere, the resultant force will be equal to
\begin{equation*}%
\frac{Nmv^{2}_{\textrm{av}}}{R}
\end{equation*}
where $v_{\textrm{av}}$ is the average molecular velocity.

The pressure $p$ of the gas equal to the force divided by
the surface area of the sphere, $4\pi R^{2}$, is
\begin{align*}%
p &= \frac{Nmv^{2}_{\textrm{av}}}{4 \pi R^{2}R} = \frac{\frac{1}{3}Nmv^{2}_{\textrm{av}}}{\frac{4}{3}\pi R^{2}} \\
& = \frac{Nmv^{2}_{\textrm{av}}}{3V}
\end{align*}
where $V$ is the volume of the sphere. 

Therefore,
\begin{equation*}%
pV = \frac{1}{3} Nmv^{2}_{\textrm{av}}
\end{equation*}
This equation was first derived by Daniel Bernoulli in 1738.\footnote{Of Swiss origin, D. Bernoulli worked and lived in Russia; he was a member of the St. Petersburg Academy of Sciences. No less well known is the activity of Johann (Jean) Bernoulli and Jakob (Jacques) Bernoulli. All three were brothers, not namesakes.}

It follows from the equation of state of an ideal gas that $pV = \textrm{const} \cdot T$; we see from the equation just derived that $pV$ is proportional to $v^{2}_{\textrm{av}}$. Hence,
\begin{equation*}%
T \propto v^{2}_{\textrm{av}} \,\, \text{or} \,\,  v_{\textrm{av}} \propto \sqrt{T}
\end{equation*}
i.e. the average velocity of an ideal gas molecule is proportional to the square root of the absolute temperature.

\section{Avogadro's Law}

Assume that a substance is a mixture of different mole­cules. Isn’t there a physical quantity characterizing a motion which would be identical for all these molecules, say, for hydrogen and oxygen, provided their tempera­tures are identical?

Mechanics yields an answer to this question. It can be proved that the average kinetic energy $mv^{2}_{\textrm{av}}/2$ of the trans­lator motion will be identical for all molecules.

This implies that for a given temperature, the average square of the molecular velocities is inversely proportion­al to the mass of the particles:
\begin{equation*}%
v^{2}_{\textrm{av}} \propto \frac{1}{m} \,\, \text{or} \,\, v_{\textrm{av}} \propto \frac{1}{\sqrt{m}} 
\end{equation*}
Let us return to the equation $pV = (1/3) Nmv^{2}_{\textrm{av}}$. Since the quantities $mv^{2}_{\textrm{av}}$ are identical for all gases at a given temperature, the number $N$ of molecules contained in a given volume $V$ at a definite pressure $p$ and temperature $T$ is identical for all gases. This remarkable law was first formulated by Amedeo Avogadro (1776-1856).

But how many molecules are there in \SI{1}{\centi\meter\cubed}? It turns out that there are \num{2.7d19} molecules in \SI{1}{\centi\meter\cubed} at \SI{0}{\celsius} and \SI{760}{\milli\meter\mercury}. This is an enormous number. So that you can feel just how great it is, let us give an example. Suppose that gas is flowing out of a \SI{1}{\centi\meter\cubed} vessel with such a speed that a million molecules leave each second. It isn’t hard to calculate that it will take the vessel a million years to get rid of the gas!

Avogadro’s law shows that under a definite pressure and temperature, the ratio of the number of molecules to the volume in which they are contained, $N/V$, is a quan­tity that is identical for all gases.

Since the density of a gas $p = Nm/V$, the ratio of the densities of gases is equal to that of their molecular masses:
\begin{equation*}%
\frac{\rho_{1}}{\rho_{2}} = \frac{m_{1}}{m_{2}}
\end{equation*}
The relative masses of molecules can therefore be determined by simply weighing gaseous substances. Such measurements once played a great role in the development of chemistry. It also follows from Avogadro’s law that for a mole of any substance in the state of an ideal gas, $pV = kN_{A}T$, where $k$ is a universal constant (named after the famous German physicist Ludwig Boltzmann) equal to \SI{1.38d-16}{\erg\kelvin}. The product $R = kN_{A}$ is called the universal gas constant.

The ideal gas law is often written as
\begin{equation*}%
pV = \mu RT
\end{equation*}
where $\mu$ is the amount of substance expressed in moles.

This equation is frequently applied in practice.

\section{Molecular Velocities}

Theory shows that for a constant temperature, the aver­age kinetic energy of molecules, $mv^{2}_{\textrm{av}}/2$, is identical. Ac­cording to our definition of temperature, this average kinetic energy of the translatory motion of the molecules of a gas is proportional to the absolute temperature.

Combining the ideal gas equation with Bernoulli’s equation we obtain 
\begin{equation*}%
\left( \frac{mv^{2}}{2} \right)_{\textrm{av}} = \frac{2}{3} kT
\end{equation*}
Temperature measurements with a thermometer filled with an ideal gas add a meaning of rare simplicity to this measure. The temperature is proportional to the average value of the energy of translatory motion of the mole­cules. Since we live in three-dimensional space, we can say that a point moving at random has three degrees of freedom. Consequently, there is $kT/2$ energy per degree of freedom of a moving particle.

Let us determine the average speed of oxygen mole­cules at room temperature, which we take to be $\SI{27}{\celsius} = \SI{300}{\kelvin}$ in round numbers. The molecular mass of oxygen is 32, so the mass of one molecule equals $32/(\num{6d23})\si{\gram}$. A simple computation yields $v_{\textrm{av}} = \SI{4.8d4}{\centi\meter\per\second}$, i.e. about \SI{500}{\meter\per\second}. Molecules of hydrogen move considerably faster. Their masses are 16 times as small, and their speeds are $\sqrt{16} = 4$ times as great, i.e. are about 2 km/s at room temperature. Let us estimate the thermal speed of a small particle which is visible through a microscope. An ordinary microscope permits us to see a dust particle of \SI{1}{\micro\meter} (\SI{d-4}{\centi\meter}) in diameter. The mass of such a particle with a density close to unity will be in the neighbourhood of \SI{5d-13}{\gram}. We obtain about  \SI{0.5}{\centi\meter\per\second} for its speed. It is not surprising that such motion is quite noticeable.

The speed of the Brownian movement of a particle with a mass of \SI{0.1}{\gram} will be only \SI{d-6}{\centi\meter\per\second} in all. It is no wonder that we do not see the Brownian movement of such particles.

We have spoken of the average speed of a molecule. But not all molecules move with the same speed; a certain fraction of the molecules move faster, but others move slower. It turns out that this can all be calculated. We shall only present the results.

At a temperature of about \SI{15}{\celsius}, for example, the average speed of nitrogen molecules is equal to \SI{500}{\meter\per\second}; 59\% of the molecules move with speeds between 300 and \SI{700}{\meter\per\second}. Only 0.6\% of the molecules move with small speeds -- from 0 to \SI{100}{\meter\per\second}. There are only 5.4\% of fast molecules with speeds greater than \SI{1000}{\meter\per\second} (\figr{fig-3.2}).

\begin{figure}[!ht]
\centering
\includegraphics[width=\textwidth]{figures/fig-03-02.pdf}
\caption{A graph showing the distribution of velocities in a gas.}
\label{fig-3.2}
\end{figure}

Each column is constructed with its base covering the velocity range it refers to and the area of each column is proportional to the percentage of molecules whose veloc­ity lies within this range.

It is also possible to calculate the distribution of molecules over the energy of their translatory motion.

The number of molecules whose energy is more than double the average is less than 10\%. The fraction of still more “energetic” molecules falls off faster and faster as the energy increases. Thus, the number of molecules whose energy is at least four times as large as the average is only 0.7\%, eight times as large as the average $0.06 \times
 \num{d-4}$\%, 16 times as large as the average $2 \times \num{d-8}$\%. 
 
The energy of an oxygen molecule moving with a speed of \SI{11}{\kilo\meter\per\second} is equal to \SI{32d-12}{\erg}. The average energy of a molecule at room temperature is equal to only \SI{6d-14}{\erg}. Therefore, the energy of an ``eleven-kilo­ metre molecule'' is at least 500 times as great as the ener­gy of a molecule with the average speed. It is not surpris­ing that the fraction of the molecules with speeds higher than \SI{11}{\kilo\meter\per\second} is equal to an unimaginably small number -- of the order of \num{d-300}.

But why are we intrigued with the speed of \SI{11}{\kilo\meter\per\second}? In the first book we spoke of the fact that only bodies having this speed can escape from the Earth. Hence, molecules which have risen to a great height can lose their ties to the Earth and take off in a distant interplane­tary trip, but for this it is necessary to have a speed of \SI{11}{\kilo\meter\per\second}. The fraction of such fast molecules, as we have seen, is so negligible that there is no danger of the Earth’s losing its atmosphere even in the course of a thousand million years.

The rate of leaving the atmosphere depends to an ex­traordinarily great degree on the gravitational energy $GMm/r$. If the average kinetic energy of a molecule is many times less than the gravitational energy, the escape of the molecules from the Earth is practically impos­sible. The gravitational energy on the surface of the Moon is 20 times as small, which gives an oxygen molecule a ``runaway'' energy of \SI{1.5d-12}{\erg}. This value exceeds that of the average kinetic energy of a molecule by a factor of only 20-25. The fraction of the molecules capable of breaking away from the Moon is equal to \num{d-17}. This is already entirely different than \num{d-300}, and computations show that the air would leave the Moon quickly enough for interplanetary space. It is not surprising that there is no atmosphere on the Moon.


\section{Thermal Expansion}

If a body is heated, the motion of its atoms (molecules) will be more intensive. They will start pushing each other away and will occupy more space. The following well-known fact is explained by this: solids, liquids and gases expand when heated.

We don’t have to say much about the thermal expansion of gases: in fact, the proportionality of gas temperature to volume was made the basis of our temperature scale.
We see from the formula $V = V_{0}T/273$ that the volume of a gas under a constant pressure grows by 1/273 (i.e. by 0.0037) of its value at \SI{0}{\celsius} for each \SI{1}{\celsius} increase in tem­perature (this situation is sometimes called Gay-Lussac’s law).

Under ordinary conditions, i.e. at room temperature and standard atmospheric pressure, most liquids expand one-third to one-half as much as gases.

We have already spoken more than once of the anoma­lous expansion of water. The volume of water decreases as it is heated up from 0 to \SI{4}{\celsius}. This anomaly in the expansion of water plays a colossal role in organic life on the Earth. In autumn, the upper layers of water become denser and sink to the bottom as they cool off. Warmer water, rising from below, takes their place. But such a mixing takes place only until the temperature of the water falls to \SI{4}{\celsius}. With a further fall in temperature, the upper layers will no longer contract, will therefore not become heavier and will not sink to the bottom. Starting with this temperature, the upper layer, gradually cooling off, reaches zero degrees and freezes.

It is only this peculiarity of water that prevents rivers from freezing down to their beds. If water were to sud­denly lose its remarkable anomaly, the disastrous con­sequences of this can be easily pictured even by a person without a rich imagination.


The thermal expansion of solids is considerably less than that of liquids. It is hundreds and thousands of times less than the expansion of gases.

Thermal expansion is an annoying hindrance in many cases. Thus, a change in the sizes of the moving parts of a clockwork with a change in temperature would lead to a change in the speed of the clock, if a special alloy, in­ var (invariant means unchanging, whence the name ``invar''), were not used for these delicate components. Invar, steel with a large nickel content, is widely used in the instrument manufacture. An invar rod is lengthened by only one-millionth when its temperature increases by \SI{1}{\celsius}.

An apparently negligible thermal expansion of a solid body can lead to serious consequences. The reason for this is that the low compressibility of solids makes it hard to hinder their thermal expansion.

When heated by \SI{1}{\celsius}, a steel rod will increase in length by only one-hundred-thousandth, i.e. by an amount unnoticeable to the unaided eye. However, in order to prevent the expansion and compress the rod one-hundred-thousandth, a force of \SI{20}{\kgf} on \SI{1}{\centi\meter\squared} is needed. And this is merely for cancelling the effect of a rise in temperature by only \SI{1}{\celsius}.

The forces arising from thermal expansion can lead to breakages and catastrophes if they are not reckoned with. Thus, in order to avoid the action of these forces, the rails of a railroad-bed are laid with clearances. One has to remember these forces when handling glassware which is easily cracked by non-uniform heating. It is therefore the practice in laboratories to use vessels made of quartz glass (fused quartz, silicon dioxide, exists in an amorphous state), which lack this drawback. For one and the same rise in temperature, a copper bar will be lengthened by a millimetre, while the same sized bar of quartz glass will change its length by the unnoticeable amount of 30-40 \si{\micro\meter}. The expansion of quartz is so insignificant that a quartz vessel can be heated by several hundred degrees and then thrown into water without any fear.

\section{Heat Capacity}

The internal energy of a body also depends, of course, on its temperature. The more a body must be heated, the greater is the energy required. In order to raise the tem­perature of a body from $T_{1}$ to $T_{2}$, it is required to supply an energy
\begin{equation*}%
Q=C(T_{2}-T_{1})
\end{equation*}
to it in the form of heat. Here $C$ is the proportionality factor, which is called the \emph{heat capacity} of the body. The definition of the concept of heat capacity follows from the formula: $C$ is the amount of heat necessary for raising the temperature by \SI{1}{\celsius}. The heat capacity also derends on the temperature: rises in temperature from 0 to \SI{1}{\celsius} and from 100 to \SI{101}{\celsius} require somewhat different amounts of heat.

The quantity $C$ is usually referred to a unit mass and called the \emph{specific heat}. It is then denoted by the small letter $c$.

The amount of heat which goes to heat up a body of mass m is given by the following formula:
\begin{equation*}%
Q = mc (T_{2} — T_{1})
\end{equation*}
In what follows we shall make use of the concept of specific heat capacity, but shall speak of the heat capac­ity of a body for the sake of conciseness. An additional guide will always be the dimension of the quantity.

The value of heat capacity varies within a rather wide range. Of course, the heat capacity of water in calorie per degree is equal to unity by definition.

Most bodies have a heat capacity less than that of water. Thus, most oils, alcohols and other liquids have heat capacities close to \SI{0.5}{\calorie\per\gram\per\kelvin}. Quartz, glass and sand have a heat capacity of the order of \SI{0.2}{\calorie\per\gram\per\kelvin}. The heat capacity of iron and copper is about \SI{0.1}{\calorie\per\gram\per\kelvin}. And here are examples of the heat capacities of some gases: hydrogen, \SI{3.4}{\calorie\per\gram\per\kelvin}; air, \SI{0.24}{\calorie\per\gram\per\kelvin}.

The heat capacities of all bodies decrease, as a rule, with a fall in temperature, and assume negligible values for most bodies at temperatures close to absolute zero. Thus, the heat capacity of copper is equal to only 0.0035 at \SI{20}{\kelvin}; this is twenty-four times less than at room tem­perature.

A knowledge of heat capacities may prove useful for solving various problems on the distribution of heat among bodies.

The difference between the heat capacities of water and soil is one of the causes determining the distinction between maritime and continental climates. Possessing approximately five times as great a heat capacity as soil, water warms up slowly and cools off just as slowly.

In summer in maritime regions, the water, having warmed up more slowly than the land, cools the air, but in winter, the warm sea gradually cools off, yielding heat to the air and making the frost less severe. It is not diffi­cult to calculate that \SI{1}{\meter\cubed} of sea water, cooling off by \SI{1}{\celsius}, warms up \SI{3000}{\meter\cubed} of air by \SI{1}{\celsius}. Consequently, in maritime regions the variations in temperature and the difference between winter and summer temperatures are less substantial than in continental regions.

\section{Thermal Conductivity}
Each object can serve as a ``bridge'' along which heat passes from a warmer body to a cooler one. For example, a tea spoon placed in a glass of hot tea is such a bridge.

Metallic objects conduct heat very well. The top of the spoon placed in the glass will become warm in the course of a second.

If it is necessary to stir a hot mixture, the handle of the stirrer must be made of wood or plastic. These solids conduct heat a thousand times worse than metals. We say ``conduct heat'', but could just as well have said ``con­duct cold''. Of course, the properties of a body do not change as a result of the direction in which a heat flow is passing through it. In freezing weather we are careful not to touch metals with our bare hands outdoors, but grasp wooden handles without fear.

Among the poor heat conductors, they are also called heat insulators, are wood, brick, glass and plastic. The walls of houses, ovens and refrigerators are made of these materials.

Among the good conductors are all the metals. The best conductors are copper and silver -- they conduct heat twice as well as iron.

Of course, not only solids can serve as ``bridges'' for the transfer of heat. Liquids also conduct heat, but much worse than metals. The thermal conductivity of metals is hundreds of times greater than that of solid and liquid non-metallic bodies.

In order to demonstrate the poor thermal conductiv­ity of water, the following experiment is performed. A piece of ice is fastened to the bottom of a test-tube filled with water, while the top of the test-tube is heated over a gas burner; the water begins boiling, but the ice is still not melting. If the test-tube were without water and made of metal, then the piece of ice would begin melting almost immediately. Water conducts heat about two hundred times worse than copper.

Gases conduct heat tens of times worse than condensed non-metallic bodies. The thermal conductivity of air is twenty thousand times smaller than that of copper.

The poor thermal conductivity of gases permits us to hold in our hands a piece of dry ice whose temperature is \SI{-78}{\celsius}, and to even hold on our palms a drop of liquid nitrogen having a temperature of \SI{-196}{\celsius}. If we do not squeeze these cold objects with our fingers, there will be no ``bum''. The reason consists in the fact that when the drop of liquid or the piece of solid is boiling very energet­ically, it is covered by a ``vapour jacket'', and the layer of gas so formed serves as a heat insulator.

The spheroidal state of a liquid, this is what one calls the state in which drops are covered by vapour, is formed whenever water gets into a very hot frying-pan.\footnote{This is the so called Leidenfrost effect. -- Damitr} Drops of boiling water, having fallen on one’s palm, severely burn one’s hand, although the difference in temperature between boiling water and a human body is less than that between a hand and liquid air. Since one’s hand is colder than the drops of boiling water, heat leaves the drops, the boiling ceases and no vapour jacket is formed.

It isn’t hard to understand that the best heat insulator is a vacuum -- emptiness. There are no carriers of heat in a vacuum, and so the thermal conductivity will be at a minimum.

Therefore, if we want to create a thermal shield, hide something warm from something cold or vice versa, the best thing to do is to erect a casing with double walls and pump the air out of the space between them. When doing this, we come across the following curious phenomenon. If we keep track of the change in thermal conductivity of the gas as it is being rarefied, we shall observe that right until the moment when the pressure reaches sev­eral millimetres of mercury column the thermal conduc­tivity remains practically constant. Our expectations are only justified when, with the passage to a higher vacuum, the thermal conductivity sharply falls off.

But what is the cause of this?

In order to understand this phenomenon, we must try to visualize the process of heat transfer in a gas.

Heat transfer from a warm place to a cold one takes place by means of the transmission of energy from one molecule to a neighbouring one. It is clear that collisions between fast and slow molecules usually lead to an ac­celeration of the slow molecules and a deceleration of the fast ones. And this means that the hot place will become colder, while the cold place will warm up.

But how will a decrease in pressure affect heat transfer? Since a decrease in pressure lowers the density, the num­ber of collisions between fast and slow molecules, during
which a transmission of energy occurs, will also decrease. This would decrease the thermal conductivity. However, a decrease in pressure leads, on the other hand, to an increase in the mean free path of the molecules, which therefore transfer heat by greater distances, and this tends to increase the thermal conductivity. Computations show that these effects compensate for each other, and so the ability to transfer heat does not change for some time as the air is being pumped out.

This will be the case until the vacuum becomes so con­siderable that the mean free path is comparable to the distance between the walls of the vessel. Now a further decrease in the pressure can no longer change the mean free path of the molecules, which are ``hopping around'' from wall to wall: the fall in density is not compensated for, and so the thermal conductivity rapidly falls in proportion to the pressure, reaching negligible values as a high vacuum is attained. The construction of a vacuum bottle is based on the use of the properties of a vacuum. Vacuum bottles are very widespread: they are applied not only for the conservation of hot and cold foods but also in science and technology. In such a case they are called Dewar flasks (vessels), in honour of their inventor. Liq­uid air, nitrogen and oxygen are transported in such vessels. Later we shall tell how these gases are obtained in a liquid state.\footnote{Everyone who has seen cylinders of vacuum bottles noticed that their walls are always silver-plated. But why? The fact is that thermal conductivity is not the only means of trans­ferring heat. There exists yet another way to transfer heat, which we shall speak of in another book, so-called radiation. Under ordinary conditions, it is much weaker than thermal conductivity, but is nevertheless quite noticeable. The walls of vacuum bottles are covered with a coating of silver precisely in order to weaken the radiation.}

\section{Convection}

But if water is such a poor heat conductor, how does it warm up in a tea-kettle? Air conducts heat even worse; then it isn’t clear why the same temperature is established in all parts of a room.

Water in a tea-kettle quickly boils because of gravity. The lower layers of water, having warmed up, expand, become lighter and rise to the top, with cold water tak­ing their place. A rapid heating occurs thanks only to \emph{convection} (derived from the Latin \emph{convectus} meaning ``bring together''). It wouldn’t be so easy to heat up water in a tea-kettle located in an interplanetary rocket.

Somewhat earlier explaining why rivers do not freeze down to the bottom, we spoke of another case of convec­tion currents of water without using this word.

Why are the radiators of a central heating placed near the floor? Why is the ventilation window made in the upper part of a window? It might be more convenient to open a ventilation window if it were lower down, and it might not be a bad idea to place radiators under the ceiling, so that they did not get in the way. If we were to follow such advice, we would soon discover that the room is not being warmed up by the radia­tor and not being aired when the ventilation window is open.

The same thing takes place with the air in a room as with the water in a tea-kettle. When the radiator is turned on, the air in the lower layers of the room begins warm­ ing up. It expands, becomes lighter and rises towards the ceiling. Heavier layers of cold air arrive in its place. And they, having warmed up, leave for the ceiling. A contin­uous air current thus arises in the room, with warm air moving up from below and cold air moving down from above. Opening a ventilation window in winter, we admit a stream of cold air into the room. It is heavier than the air in the room, and so goes down, forcing out the warm air, which rises towards the top of the room and leaves through the ventilation window.

A kerosene lamp flames up well only when it is covered with a tall piece of glass. One should not think that the glass is needed only in order to shield the flame from the wind. Even in the calmest weather, the brightness of the flame immediately increases when the glass is put on the lamp. The role of the glass consists in intensifying the stream of air approaching the flame -- in creating a draught. This occurs because the air inside the glass, deprived of the oxygen that was used for the burning, quick­ly warms up and rises, while pure cold air moves into its place through the holes made in the burner of the lamp.
The taller the glass, the better will the lamp burn. In fact, the speed with which cold air rushes into the burner of the lamp depends on the difference in weight between the heated column of air in the lamp and the cold air outside it. The higher the column of air, the greater this difference in weight, and so the faster the movement.

Factory chimneys are also made high for this reason. An especially rapid influx of air, a good draught, is needed for a factory furnace. It is achieved as a result of a high chimney.

The lack of convection in a rocket devoid of weight makes it impossible to use matches, lamps or gas burners: the products of combustion would smother the flame. Air is a poor conductor; we can conserve heat with its aid, but only under one condition: if we avoid convection, the mixing of warm and cold air, which will bring to naught the thermal-insulation properties of air.


The elimination of convection is achieved by applying various kinds of porous and fibrous bodies. It is difficult for air to move inside such bodies. All bodies of this kind are good heat insulators, thanks only to their ability to retain a layer of air. But the thermal conductivity of the substance itself of which the fibres or the walls of the
pores consist can be not very small.

A good fur coat is made of a dense fur containing as many fibres as possible; eiderdown can be used to make warm sleeping bags weighing less than half a kilogram, due to the exceptional thinness of its fibres. Half a kilo­gram of this down can ``detain'' as much air as tens of kilograms of sheet wadding.

Storm windows are made in order to reduce the con­vection. The air between the panes does not participate in the mixing of layers of air which takes place within the room.

On the contrary, every movement of the air intensifies the mixing and increases the transfer of heat. This is precisely why we fan ourselves or turn on the ventilator when we want the heat to go away faster. This is also what makes it colder in the wind. But if the air temperature is higher than one’s body temperature, the mixing has the opposite effect, and a wind feels like a hot breath of air.

The problem involved in a steam boiler consists in obtaining steam heated to the required temperature as quickly as possible. The natural convection in a gravita­tional field is quite insufficient for this. Therefore, the creation of an intensive circulation of water and steam, leading to the mixing of warm and cold layers, is one of the basic problems in the construction of steam boilers.


% !TEX root = pfe-book2.tex
%!TEX TS-program = pdflatex
%!TEX encoding = UTF-8 Unicode


\cleardoublepage
%\mainmatter
\chapter{States of Matter}
\label{ch-04}


\section{Iron Vapour and Solid Air}

A strange combination of words, isn’t it? However, this is by no means nonsense: iron vapour and solid air exist in nature, only not under ordinary conditions.

But what conditions are we talking about? The state of a substance is determined by two circumstances: tem­perature and pressure.

Our lives proceed under conditions which change rela­tively little. The air pressure varies by several percent about a value of one atmosphere; the temperature of the air, say, near Moscow, lies in the interval from $-30$ to \SI{+30}{\celsius}; in the absolute scale, in which the lowest pos­sible temperature ( \SI{-273}{\celsius}) is taken as zero, this interval will look less impressive: 240-300 \si{\kelvin}, which is also only $\pm$10\% of the average value.

It is quite natural that we have become accustomed to these ordinary conditions, and so when speaking simple truths such as ``iron is a solid, air is a gas'', we forget to add ``under standard conditions''.

If iron is heated, it will first melt and then vaporize. If air is cooled, it will first liquefy and then solidify.

Even if the reader has never come across iron vapour or solid air, he or she will probably believe without difficulty that by means of a change in temperature it is possible to obtain any substance in a solid, liquid and gaseous state or, as is also said, in a solid, liquid or gaseous phase.

It is easy to believe this because everyone has observed a substance without which life on the Earth would be impossible, and that substance is in the form of a gas, a liquid, or a solid. We are speaking, of course, of water.

But under what conditions does a transformation of a substance from one state to another occur?

\section{Boiling}

If we lower a thermometer into water which has been poured into a tea-kettle, turn on the electric stove and watch the mercury in the thermometer, we shall see the following: the level of the mercury will inch upwards almost immediately. Now it is already 90, 95 and finally \SI{100}{\celsius}. The water begins boiling and simultaneously the mercury stops rising. The water has already been boiling for many minutes, but the level of the mercury does not change. The temperature will not change until all the water has boiled away (\figr{fig-4.1}).

\begin{figure}[!ht]
\centering
\includegraphics[width=0.4\textwidth]{figures/fig-04-01.pdf}
\caption{Temperature stays constant at boiling.}
\label{fig-4.1}
\end{figure}

But what is the heat used for if the temperature of the water does not change? The answer is obvious. The process of transforming water into steam requires energy.

Let us compare the energy of a gram of water and a gram of the steam created out of it. The molecules of the steam are distributed farther from each other than the water molecules. It is obvious that because of this the potential energy of the water will differ from that of the steam.

The potential energy of attracting particles decreases as they approach. The energy of the steam is therefore greater than that of the water, and so the transformation of water into steam requires energy. This excess energy is imparted by the electric stove to the water boiling in the tea-kettle.

The energy needed for transforming water into steam is called its \emph{heat of vaporization}. In order to transform \SI{1}{\gram} of water into steam, 539 calories are required (this figure is for a temperature of \SI{100}{\celsius}). If 539 calories are used for \SI{1}{\gram}, then $18 \times 539 \approx 9700$ calories will be supplied to 1 mole of water. This amount of heat must be consumed in breaking the intermolecular bonds. One can compare this figure with the amount of work necessary for breaking the intramolecular bonds. In order to split one mole of steam into atoms, about 220 000 calories are required, i.e. 25 times as much energy. This directly proves the weakness of the forces binding molecules to each other, as compared to the forces finding atoms together in a mole­cule.

\section{Dependence of Boiling Point on Pressure}
The boiling point of water is equal to \SI{100}{\celsius}; one might think that this is an inherent property of water, that water will always boil at \SI{100}{\celsius}, no matter where and under what conditions it may be.

But this is not so, and people who live high up in the mountains are perfectly well aware of this.

There is a tourist cabin and a scientific station near the top of Mt. Elbrus. Novices are sometimes amazed at ``how hard it is to boil an egg in boiling water'' or wonder ``why boiling water doesn’t scald.'' In such cases, it is pointed out to them that water is already boiling at \SI{82}{\celsius} on the top of Mt. Elbrus.

But what causes this? What physical factor interferes with boiling? And does the height above sea level have any significance?

This physical factor is the pressure acting on the surface of the liquid. It isn’t necessary to climb to the top of a mountain in order to check the validity of what we have said.

If we place a bell glass over water that is being heated and pump air into or out of it, we can convince ourselves that the boiling point is raised by an increase in pressure and lowered by a decrease in pressure.

Water boils at \SI{100}{\celsius} only at a definite pressure -- \SI{760}{\milli\meter\mercury} (or \SI{1}{\atmos}).

The curve showing the dependence of the boiling point on the pressure is depicted in \figr{fig-4.2}. The pressure is equal to  \SI{0.5}{\atmos} on the top of Mt. Elbrus, and a boiling point of \SI{82}{\celsius} corresponds to this pressure.

\begin{figure}[!ht]
\centering
\includegraphics[width=\textwidth]{figures/fig-04-02.pdf}
\caption{Change in boiling point as a function of pressure.}
\label{fig-4.2}
\end{figure}

But it is even possible to refresh oneself in hot weather with water boiling at 10-15 \si{\milli\meter\mercury}. At such pressures, the boiling point will fall to 10-15 \si{\celsius}.

One can even obtain ``boiling water'' having a temper­ature of freezing water. One has to lower the pressure to \SI{4.6}{\milli\meter\mercury} for this.

It is possible to observe an interesting scene by placing an uncovered vessel with water under a bell glass and pumping the air out. This will make the water boil, but boiling requires heat. But since there is no air, the water has to give up its own energy. The temperature of the boiling water starts falling, but since the pumping con­tinues, the pressure also falls. Therefore, the boiling will not cease, the water will continue cooling off and will finally freeze.

Such a boiling of cold water occurs not only as a result of air being pumped out. For example, when a marine screw propeller rotates, the pressure in the layer of water moving rapidly about the metallic surface will fall sharp­ly, and so the water in this layer will boil, i.e. many bubbles filled with steam will appear in it. This phe­nomenon is called \emph{cavitation} (from the Latin \emph{cavus} meaning ``hollow'').

Decreasing the pressure, we lower the boiling point. And raising it? A graph analogous to ours answers this question. A pressure of \SI{15}{\atmos} can so delay boiling of water that it will begin only at \SI{200}{\celsius}, while a pressure of \SI{80}{\atmos} will make water boil only at \SI{300}{\celsius}.

Thus, a definite boiling point correspond to a definite external pressure. But we may also turn this assertion ``around'' by saying that ``a definite pressure corresponds to each boiling point of water''. This pressure is called the \emph{vapour pressure}.

The curve depicting the dependence of the boiling point on the pressure is simultaneously the curve of the vapour pressure as a function of the temperature.

The numbers plotted on the boiling point graph (or on the vapour pressure graph) show that the vapour pressure changes very sharply with a change in the temperature. At \SI{0}{\celsius} (i.e. \SI{273}{\kelvin}) the vapour pressure is equal to \SI{4.6}{\milli\meter\mercury}, at \SI{100}{\celsius} (\SI{373}{\kelvin}) it equals \SI{760}{\milli\meter\mercury}, i.e. has increased by a factor of 165. With a doubling of the temperature from \SI{0}{\celsius} (i.e. \SI{273}{\kelvin}) to \SI{273}{\celsius} (i.e. \SI{546}{\kelvin}), the vapour pressure grows from \SI{4.6}{\milli\meter\mercury}  to almost \SI{60}{\atmos}, i.e. by a factor of about \num{10000}.

Therefore, the boiling point, on the contrary, changes rather slowly with a change in the pressure. When the pressure is doubled, from 0.5 to \SI{1}{\atmos}, the boiling point grows from \SI{82}{\celsius} (i.e. \SI{355}{\kelvin}) to \SI{100}{\celsius} (i.e. \SI{373}{\kelvin}), and with a doubling from 1 to \SI{2}{\atmos}, from \SI{100}{\celsius} (i.e. \SI{373}{\kelvin}) to \SI{120}{\celsius} (i.e. \SI{393}{\kelvin}).

The same curve that we are now considering also con­trols the condensation of steam into water.

Steam can be transformed into water by either compress­ing or cooling it.

During the course of condensation, just as for boiling, the point will not move off the curve until the transfor­mation of steam into water or water into steam is com­pletely finished. This can also be formulated as follows: the coexistence of the liquid and the vapour phase is pos­sible under the conditions of our curve and only under these conditions. If, moreover, no heat is supplied or removed, the amount of vapour and liquid in a closed vessel will remain constant. We say that such a vapour and liquid are in equilibrium, and a vapour in equilibrium with its liquid is called saturated.

The boiling and condensation curve has, as we see, yet another meaning -- it is the curve of the equilibrium of liquid and vapour. The equilibrium curve divides the plane of the diagram into two parts. To the left and above the curve (towards higher temperatures and lower pres­sures) is the stable-vapour region. To the right and below the curve is the stable-liquid region.

The vapour-liquid equilibrium curve, i.e. the curve of the dependence of the boiling point on the pressure, or, what is the same thing, of the vapour pressure on the temperature, is approximately identical for all liquids. In some cases, the change may be somewhat sharper, in others somewhat slower, but the vapour pressure always grows rapidly with rise in temperature.

We have already used the words ``gas'' and ``vapour'' many times. These two words are more or less synonyms. One may say: water gas is the vapour of water, oxygen gas is the vapour of liquid oxygen. Nevertheless, a certain habit has been formed regarding the usage of these two words. Since we are accustomed to a definite, rather small range of temperatures, we usually apply the word ``gas'' to those substances whose vapour pressure is higher than atmospheric pressure at standard temperatures. On the contrary, we speak of a vapour when a substance is more stable in the form of a liquid at room temperature and atmospheric pressure.

\section{Evaporation}

Boiling is a rapid process, and not even a trace of boil­ing water remains after a short time -- it is transformed into steam.

But there is also another phenomenon whereby water or some other liquid is transformed into a vapour -- \emph{evap­oration}. Evaporation takes place at any temperature and regardless of the pressure, which is always close to \SI{760}{\milli\meter\mercury} under ordinary conditions. Evaporation, unlike boiling, is a very slow process. A bottle of eau-de-cologne which we forgot to close will turn out to be empty after several days; water will remain in a saucer for a longer time, but sooner or later it too will turn out to be dry.

Air plays a big role in the process of evaporation. It does not, by itself, prevent water from evaporating. As soon as we uncover the surface of a liquid, water mole­cules will begin moving into the nearest layer of air. The density of the vapour in this layer will quickly increase; after a short time, the pressure of the vapour will become equal to the vapour pressure at the temperature of the surroundings. Moreover, the vapour pressure will be exactly the same as in the absence of air.

The passage of vapour into the air does not, of course, mean an increase in pressure. The total pressure in the space on top of the water surface does not increase; it is only the fraction of this pressure which is borne by the vapour that increases, and the fraction of the air corre­spondingly decreases as it is displaced by the vapour.

There is vapour mixed with air over the water; higher up are layers of air without vapour. They will inevitably mix. Water vapour will continually move into higher layers, and its place in the lower layer will be taken by air which does not contain any water molecules. There­fore, room will always be made for new water molecules in the layer closest to the water. Water will continually evaporate, maintaining the pressure of the water vapour at the surface equal to the vapour pressure, and the pro­cess will continue until the water has completely evapo­rated.

We began with examples involving eau-de-cologne and water. It is well known that they evaporate with different speeds. Ether flies away with exceptional rapidity, alcohol is rather quick and water is much slower. We shall immediately understand why this is so if we find the values of the vapour pressure for these liquids in a handbook, say, at room temperature. Here are the figures; ether -- \SI{437}{\milli\meter\mercury}, alcohol -- \SI{44.5}{\milli\meter\mercury} and water -- \SI{17.5}{\milli\meter\mercury}.

The greater the vapour pressure, the more vapour there will be in the adjacent layer of air and the faster the liquid will evaporate. We know that vapour pressure increases with temperature. It is clear why the rate of evap­oration increases with heating.

It is also possible to influence the rate of evaporation by other means. If we want to aid the evaporation, we must take the vapour away from the liquid more rapidly, i.e. speed up the mixing with air. This is precisely why evaporation is greatly speeded up by blowing on the liq­uid. Water, although it has a relatively low vapour pres­sure, will disappear rather quickly if the saucer is placed in the wind.

It is therefore clear why a swimmer, having come out of the water, feels cold in the wind. The wind speeds up the mixing of air with vapour and so increases the rate of evaporation, but the swimmer’s body is forced to give up heat for the evaporation.

The way a person feels depends on how much water vapour there is in the air. Both dry and moist air are unpleasant. The humidity is regarded as standard when it is equal to 60\%. This means that the density of water vapour is 60\% of the density of saturated water vapour at the same temperature.

If moist air is cooled, the pressure of the water vapour in it will eventually equal the vapour pressure at this temperature. The vapour will become saturated and will begin condensing into water with a further fall in tem­perature. The morning dew moistening the grass and the leaves appears precisely as a result of this phenome­non.

At \SI{20}{\celsius} the density of saturated water vapour is about \num{2e-5} \si{\gram\per\centi\meter\cubed}. We shall feel fine if the amount of water vapour in the air is 60\% of this figure, only a bit more than one-hundred-thousandth of a gram per cubic centi­metre.

Although this is a small number, it leads to an impres­sive amount of water in a room. It is not difficult to calculate that in an average sized room of \SI{12}{\meter\squared} area and \SI{3}{\meter} height, ``there will be room'' for about a kilogram of water in the form of saturated vapour.

Consequently, if we place an open barrel of water in a room sealed up tight, then, regardless of the barrel’s volume, a litre of water will evaporate.

It is interesting to compare this result for water with the corresponding figures for mercury. At the same tem­perature of \SI{20}{\celsius}, the density of saturated mercury vapour is \SI{d-9}{\gram\per\centi\meter\cubed}. There will be room for at least \SI{1}{\gram} of mercury vapour in a room of the size we have just considered.

Incidentally, mercury vapour is very poisonous, and one gram of it can seriously injure any person’s health. When working with mercury, it is necessary to see to it that not even the smallest drop is spilt.

\section{Critical Temperature}

How can we turn a gas into a liquid? The boiling point graph answers this question. A gas can be turned into a liquid by either lowering the temperature or raising the pressure.

In the $19^{\textrm{th}}$ century, the problem of raising pressures seemed to be easier than that of lowering temperatures. At the beginning of that century, the great English phys­icist Michael Faraday (1791-1867) succeeded in compress­ing gases to the value of their vapour pressures and in this manner transforming many gases (chlorine, carbon dioxide, etc.) into liquids.

However, certain gases, such as hydrogen, nitrogen, oxygen, simply could not be liquefied. No matter how much the pressure was increased, they did not turn into liquids. One might have thought that oxygen and other gases cannot be liquid. They were regarded as true, or constant, gases.

But as a matter of fact, the failures were caused by a lack of understanding of one important circumstance. Let us consider a liquid and a vapour which are in equi­librium, and think of what happens to them with an increase in the boiling point and, of course, a corresponding increase in the pressure. In other words, let us imagine that a point on the boiling point graph is moving upwards along the curve. It is clear that as the temperature rises, the liquid expands and its density falls. But as for the vapour, an increase in the boiling point is, of course, con­ducive to its expansion, but, as we have already said, the pressure of the saturated vapour grows considerably faster than the boiling point. 
Therefore, the density of the vapour does not fall, but, on the contrary, rapidly rises with an increase in the boiling point.

Since the density of a liquid falls, and the density of a vapour rises, moving ``upwards'' along the boiling point curve, we shall inevitably arrive at the point for which the densities of the liquid and the vapour are equal (\figr{fig-4.3}).

At this remarkable point, called \emph{critical}, the boiling point curve breaks off. Since all distinctions between a gas and a liquid are related to a difference in density, the properties of the liquid and the gas become identical at the critical point. Each substance has its own critical temperature and critical pressure. Thus, the critical point for water corresponds to a temperature of \SI{374}{\celsius} and a pressure of \SI{218.5}{\atmos}.
\begin{figure}[!ht]
\centering
\includegraphics[width=0.5\textwidth]{figures/fig-04-03.pdf}
\caption{The critical point.}
\label{fig-4.3}
\end{figure}

If we compress a gas whose temperature is below criti­cal, the process of its compression can be depicted by an arrow intersecting the boiling point curve (\figr{fig-4.4}). This implies that at the moment of attaining a pressure equal to the vapour pressure (the point of intersection of the arrow with the boiling point curve), the gas starts condensing into a liquid. If our vessel were transparent, then at this moment we would see a layer of liquid forming at the bottom of the vessel. If the pressure does not change, the layer of liquid will grow until all of the gas has finally been transformed into the liquid. A further compression will now require an increase in pressure.
\begin{figure}[!ht]
\centering
\includegraphics[width=0.5\textwidth]{figures/fig-04-04.pdf}
\caption{Implication of critical point in liquefaction of a gas.}
\label{fig-4.4}
\end{figure}
The situation is entirely different for the compression of a gas whose temperature is above critical. The process of compression can again be still depicted in the form of an arrow going upwards. But now this arrow does not intersect the boiling point curve. Hence, during the com­pression the vapour will not condense, but will only con­tinually become denser.

The existence of a gas-liquid interface is impossible at a temperature above critical. When compressed to arbi­trary densities, a uniform substance will be found under the piston, and it is hard to say when one may call it a gas and when a liquid.

The presence of the critical point shows that there is no difference in principle between the liquid and gaseous states. At first sight, it might seem that there is no such difference in principle only in case we are dealing with temperatures above critical. This, however, is not so. The existence of the critical point indicates the possi­bility of transforming a liquid, a genuine liquid, which can be poured into a glass, into the gaseous state avoiding the boiling process.

The path of such a transformation is shown in \figr{fig-4.4}. An unquestionable liquid is marked by a cross. If we lower the pressure somewhat (arrow pointing downwards), it will boil; it will also boil in case we raise the tempera­ture somewhat (arrow pointing to the right). But let us choose a different way. Let us compress the liquid so hard that its pressure exceeds critical. The point repre­senting the state of the liquid will rise vertically. Then let us heat the liquid -- this process will be depicted by a horizontal line. Now, after we have found ourselves to the right of the critical temperature, let us lower the pressure to its initial value. If we now decrease the tem­perature, we can obtain a most genuine vapour which could have been obtained from this liquid along a simpler and shorter path.

Therefore, by changing the pressure and temperature so that the critical point is avoided, it is always possible to obtain a vapour by means of a continuous transition from a liquid and vice versa. Such a continuous transition does not require boiling or condensation.

Early attempts to liquefy such gases as oxygen, nitro­gen and hydrogen were unsuccessful because the existence of the critical temperature was unknown. These gases have very low critical temperatures: \SI{-119}{\celsius} for oxygen, \SI{-147}{\celsius} for nitrogen and \SI{-240}{\celsius} or \SI{33}{\kelvin} for hydrogen. The record holder is helium, whose critical temperature is equal to \SI{4.3}{\kelvin}. It is possible to transform these gases to the liquid phase in only one way: we must lower their temperatures below the indicated values.

\section{Obtaining Low Temperatures}

A significant decrease in temperature can be attained by various means. But the idea involved in all these meth­ods is one and the same: we must force the body we want to cool to expend its internal energy.

But how can this be done? One of the ways is to make the liquid boil, not bringing in any heat from without. For this, as we know, it is necessary to decrease the pres­sure -- to reduce it to the value of the vapour pressure. The heat expended on the boiling will be taken from the liquid, and hence the temperature of the liquid and va­pour will fall; therefore, so will the vapour pressure. Consequently, in order that the boiling not cease but proceed more rapidly, it is necessary to continually pump air out of the vessel with the liquid.

However, a limit is reached to the fall in temperature during this process: the vapour pressure will eventually become completely negligible, and so even the most pow­erful pumps will be unable to create the required pres­sure.

In order to continue the lowering of temperature, we can, by cooling a gas with the aid of the liquid already obtained, convert it into a liquid with a lower boiling point. It is now possible to repeat the pumping process with the second substance and in this way to obtain lower temperatures. In case of necessity, such a cascade method of obtaining low temperatures can be extended.

Precisely in such a manner was this problem dealt with at the end of the past century, the liquefaction of gases was carried out in stages: ethylene, oxygen, ni­trogen and hydrogen, substances with boiling points of \SI{-103}{\celsius}, \SI{-183}{\celsius}, \SI{-196}{\celsius} and \SI{-253}{\celsius}, were successively converted into liquids. Having liquid hydrogen available, one can also obtain the lowest boiling liquid-helium (\SI{-269}{\celsius}). The neighbour ``to the left'' helped obtain the neighbour ``to the right''.

The cascade method of cooling is about a hundred years old. Liquid air was obtained by this method in 1877. Liquid hydrogen was first obtained in 1884-1885. Finally, the last stronghold was taken after another twenty years: helium, the substance with the lowest critical tempera­ture, was converted into a liquid in 1908 by Heike Kamerlingh Onnes (1853-1926) in Leiden, Holland. The 70th anniversary of this important scientific achievement was widely celebrated.

For many years, the Leiden laboratory was the only ``low-tempe-rature'' laboratory. But now there exist tens of such laboratories in many countries, not to mention the factories producing liquid air, nitrogen, oxygen and he­lium for technical purposes.

The cascade method of obtaining low temperatures is now rarely applied. In technical installations for lowering temperatures, a different means of decreasing the internal energy of a gas is applied: the gas is forced to expand ra­pidly and perform work at the expense of its internal energy.

If, for example, air is compressed up to several atmo­spheres and let into an expander, then when it performs the work involved in displacing a piston or rotating a turbine, it will cool off so abruptly that it liquefies. If carbon dioxide is let out of a cylinder with great speed, it will cool off so abruptly that it is converted into ``ice'' in the air.

Liquid gases have found wide application in technology. Liquid oxygen is used for explosives and as a component of the fuel mixture in jet engines.

The liquefaction of air is used in technology for sepa­rating the gases constituting air.
The temperature of liquid air is widely used in various branches of technology. But this temperature isn’t low enough for many physical investigations. In fact, if we convert the relevant temperatures expressed on the cen­tigrade scale to their values on the Kelvin scale, we shall see that the temperature of liquid air is approximately one-third of room temperature. Much more interesting for physics are ``hydrogen'' temperatures, i.e. temperatures of the order of 14-20 \si{\kelvin}, and especially ``helium'' tem­peratures. The lowest temperature obtained by pumping out liquid helium is \SI{0.7}{\kelvin}.

Physicists have succeeded in coming much closer to absolute zero. At the present time, temperatures have been obtained which are only several thousandths of a degree above absolute zero. However, these extremely low temperatures are obtained by methods which do not resemble those described above.

In recent years, cryogenics, the physics of low tem­peratures, has created the need leading to the founding of a new branch of industry. This branch is engaged in the manufacture of equipment and apparatus enabling large volumes and long conductors to be held at temperatures close to absolute zero.

\section{Supercooled Vapours and Superheated Liquids}

In passing through the boiling points, a vapour ought to condense, be transformed into a liquid. However, it turns out that if a vapour is very pure and does not come in contact with a liquid, we are able to obtain it in the form of a supercooled or supersaturated vapour -- a vapour which should have already become a liquid a long time ago.
A supersaturated vapour is very unstable. Sometimes it is sufficient to shake the vessel containing the vapour or throw a couple of grains into it for the delayed con­densation to immediately begin.

Experience shows that the condensation of steam mole­cules is greatly eased by the introduction of small alien particles. The supersaturation of water vapour does not occur in dusty air. It is possible to bring about conden­sation with puffs of smoke, since smoke consists of tiny solid particles. When the particles enter the steam, they gather molecules around themselves and become centres of condensation.

Thus, even though unstable, a vapour can exist in the region of temperatures fit for liquid ``life''.

But can a liquid ``live'' in the region of a vapour under those same conditions? In other words, is it possible to superheat a liquid?

It turns out to be possible. For this it is necessary to prevent the molecules of the liquid from breaking away from its surface. A radical means of achieving this is liquidating the free surface, i.e. placing the liquid in a vessel where it would be compressed on all sides by solid walls. Liquids have been successfully superheated in this manner by several degrees, i.e. one is able to displace a point depicting the state of a liquid to the right of its boiling point curve (see \figr{fig-4.4}).

Superheating is the displacement of a liquid into the region of a vapour; therefore, the superheating of a liquid can be achieved by lowering the pressure, as well as by supplying heat.

The former method can be used to obtain a surprising result. Water or some other liquid thoroughly freed of dissolved gases, which is not easy to do, is placed in a vessel with a piston reaching the surface of the liquid. The vessel and the piston should be wet by the liquid.

If we now draw the piston towards ourselves, the water cohering to the bottom of the piston will move along with it. But the layer of water cohering to the piston will pull the next layer of water after it, this layer will pull the one lying below it. As a result, the liquid will stretch.

The column of water will finally break (it is precisely the column that will break, but the water will not break away from the piston), but this will occur when the force on a unit of area attains tens of kilograms. In other words, a negative pressure of tens of atmospheres is created in the liquid.

The vapour phase of a substance is stable even for small positive pressures. And a liquid can be made to have a negative pressure. You couldn’t think of a more striking example of ``superheating''.

\section{Melting}
There is no solid body which would withstand a continual rise in temperature. Sooner or later a solid piece is transformed into a liquid; true, in certain cases we will not succeed in reaching the melting point -- a chem­ical decomposition may take place.
Molecules move more intensively as the temperature increases. Finally, the moment arrives when the preservation of order among the wildly ``swinging'' mole­cules becomes impossible. The solid body melts. Tungsten has the highest melting point: \SI{3380}{\celsius}. Iron melts at \SI{1539}{\celsius}, gold at \SI{1063}{\celsius}. Incidentally, there are also easily melted metals. Mercury, as is well known, will even melt at a temperature of \SI{-39}{\celsius}. Organic substances do not have high melting points. Naphthalene melts at \SI{80}{\celsius}, toluene at \SI{-94.5}{\celsius}.
 
It is not at all difficult to measure the melting point of a body, especially if it melts within the interval of temperatures which can be measured by an ordinary ther­mometer. It is completely unnecessary to keep one’s eye on the melting body. It is sufficient to look at the mercury column of the thermometer. The temperature of the body increases until the melting begins (\figr{fig-4.5}). As soon as the melting begins, the rise in temperature
ceases, and the temperature remains constant until the process of melting is completely finished.
\begin{figure}[!ht]
\centering
\includegraphics[width=0.5\textwidth]{figures/fig-04-05.pdf}
\caption{Finding melting point of a solid.}
\label{fig-4.5}
\end{figure}
Just as in the conversion of a liquid into a vapour, the conversion of a solid into a liquid requires heat. The amount of heat needed for this is called the \emph{latent heat of fusion}. For example, the melting of one kilogram of ice requires \SI{80}{\kilo\calorie}.

Ice is one of the bodies possessing a large heat of fusion. The melting of ice requires, for example, more than ten times as much energy as the melting of the same mass of lead. Of course, we are talking about the melting proper; here we are not dealing with the fact that before lead begins to melt, it must be heated up to \SI{+327}{\celsius}. The thawing of snow is delayed as a result of the large heat of fusion of ice. Imagine that its heat of fusion were ten times smaller. Then the spring thaws would lead each year to inconceivable disasters.

Thus, the heat of fusion of ice is great, but it is also small if we compare it with the heat of vaporization -- \SI{540}{\kilo\calorie\per\kilo\gram} (about seven times as small). This difference, by the way, is completely natural. Converting a liquid into a vapour, we must tear its molecules away from each other, but in melting a solid, we only have to destroy the order in the distribution of the molecules, leaving them almost at the same distances from each other. It is clear that less work is required in the latter case.

The possession of a definite melting point is an important feature of crystalline substances. It is on the basis of precisely this property that they are easily distinguished from other solids, called amorphous bodies or glasses. Glasses are found among organic substances as well as among inorganic ones. Window glass is usually made out of silicates of sodium and calcium; an organic glass (it is also called Plexiglas) is often placed on a desk.

In contrast to crystals, amorphous substances do not have a definite melting point. Glass does not melt, but softens. When heated, a piece of hard glass becomes soft, so that it can be easily bent or stretched; it begins to change its form at a higher temperature under the influence of its own weight. As it is further heated, the thick viscous mass of glass assumes the form of the vessel in which it is lying. This mass is at first as thick as honey, then as sour cream and, finally, becomes almost like a liquid with a small viscosity, such as water. With the best will in the world, here we are unable to single out a definite temperature at which the solid entered the liquid phase. The reasons for this lie in the radical differ­ence between the structure of glass and that of crystalline bodies. As has been said above, the atoms in amor­phous bodies are distributed disorderly. The structure of glass resembles that of a liquid. The molecules in hard glass are already distributed disorderly. Hence, a rise in the temperature of glass merely increases the amplitude of the oscillations of its molecules, gradually giving them more and more freedom of movement. Therefore, glass softens gradually and does not display the sharp transition from ``solid'' to ``liquid'', which is char­acteristic of a transition from the distribution of mole­cules in a strict order to their disordered distribution.

When discussing a boiling point curve, we said that a liquid and a vapour can exist, although in an unstable state, in alien regions -- a vapour can be supercooled and moved to the left of the boiling point curve, while a liq­uid can be superheated and drawn off to the right of the boiling point curve.

Are the analogous phenomena with respect to a crystal and a liquid possible? It turns out that the analogy here is incomplete.

If a crystal is heated, it will begin melting at its melting point. We shall not succeed in superheating a crystal. On the contrary, in cooling a liquid, we can, if we take certain steps, ``slip past'' the melting point with compa­rative ease. We are able to achieve considerable supercool­ings of certain liquids. There are even such liquids which are easy to supercool, but hard to make them crystallize. As such a liquid is cooled, it becomes more and more vis­cous and finally hardens without having crystallized. Glass is like that.

It is also possible to supercool water. Drops of mist can fail to freeze even during the most severe frosts. If a crystal of a substance (priming) is thrown into a super­ cooled liquid, then crystallization will immediately begin.

Finally, in many cases a delayed crystallization can begin as a result of shaking or other random events. It is known, for example, that crystalline glycerine was first obtained while being transported by train. After lying around for a long time, glass can begin to crystallize (devitrify).

\section{How to Grow a Crystal}

Under definite known conditions, crystals can be grown from almost any substance. Crystals can be obtained from a solution or from a melt of the given substance, as well as from its vapour (black rhombic crystals of iodine, for instance, are readily deposited from iodine vapour at standard pressure without an intermediate conversion into the liquid state).

Start by dissolving common salt or sugar in water. At room temperature (\SI{20}{\celsius}), you can dissolve up to \SI{70}{\gram} of salt in a thick glass (holding about \SI{200}{\gram} of water). If you keep adding salt, it will not dissolve, but will settle to the bottom in the form of a residue. A solution that can dissolve no more solute (as it is called) is said to be saturated with the given substance. If we change the temperature, the solubility of the solute in the solvent is changed as well. Everyone knows that hot water dissolves most substances more readily than cold water.

Imagine now that you have prepared a saturated solu­tion, say, of sugar, at a temperature of \SI{30}{\celsius} and cool it to \SI{20}{\celsius}. At \SI{30}{\celsius}, you can dissolve \SI{223}{\gram} of sugar in \SI{100}{\gram} of water, and at \SI{20}{\celsius}, only \SI{205}{\gram}. Then, in cooling from 30 to \SI{20}{\celsius}, \SI{18}{\gram} of sugar turn out to be ``redundant'' and, as they say, are precipitated from the solution. Hence, one of the possible ways to obtain crystals is to cool a saturated solution.

We can accomplish the same in a different way. Pre­pare a saturated solution of salt and let it stand in an open glass. After some time has passed you will find crystals of salt at the bottom. Why did they form? Look at the glass again, more carefully, and you can see that another change occurred together with the appearance of crystals: there is less water. The water evaporated and left ``redundant'' matter in the solution. Thus, still an­ other way to form crystals is to evaporate the solution.

How are crystals formed from a solution?

We mentioned that the crystals precipitate from the solution. Does this mean that no crystal was observed for a whole week and then, in a flash, it appears as if by magic? No, this is not so; crystals grow. You cannot, of course, observe the initial instant of growth with the naked eye. First, a few of the randomly moving mole­cules or atoms of the solute assemble by chance in approx­imately the same order required to form the crystal lattice. Such a group of atoms or molecules is called a \emph{nucleus}.

Experiments show that nuclei are more frequently formed when the solution contains extremely fine par­ticles, dust specks, of some foreign substance. Crystalli­zation proceeds at the highest rate and most easily if a tiny seed crystal is put into the saturated solution. Then the precipitation of the solid substance from the solution consists in the growth of the seed crystal rather than in the formation of new small crystals, i.e. in nucleation.

The growth of a nucleus does not, of course, differ in any way from the growth of a seed crystal. The advantage of using a seed crystal is that it ``draws'' to itself the sub­ stance separating out of the solution, hindering, in this way, the formation of a large number of nuclei. If a great many nuclei are formed at once, they impede each other in growing and we cannot obtain large crystals.

How do new portions of atoms or molecules, separating out of the solution, arrange themselves on the surfaces of the nucleus?

It has been found that the growth of a nucleus or seed crystal consists in the seemingly outward motion of its faces in a direction perpendicular to each face so that they remain parallel to their initial positions (the crys­tal seems to expand). Naturally, the angles between the faces remain constant (we already know that the con­stancy of these angles is one of the vital features of crys­tals and is due to their lattice structure).
\begin{figure}[!ht]
\centering
\includegraphics[width=\textwidth]{figures/fig-04-06.pdf}
\caption{Modes of crystal growth.}
\label{fig-4.6}
\end{figure}
The successive outlines during the growth of three differ­ent crystals of the same substance are shown in \figr{fig-4.6}. Similar pictures can be observed in a microscope. For the crystal shown at the left the number of faces remains constant during the growth. The middle crystal is an example of how a new face appears during the growth (at the upper right) and subsequently disappears.

It is important to note that the rate of growth of the faces, i.e. the velocity at which the faces seem to move while remaining parallel to their initial positions, is not the same for various faces. We also find that the faces that disappear are those that grow fastest, for instance, the lower left face in the middle crystal. The slowest-growing faces, on the contrary, become the videst or, as we say, best-developed, ones. It is a principle that a grown crystal is always bounded by its slowest-growing faces.

This is especially clear from the illustration at the right in \figr{fig-4.6}. Here a formless fragment acquires the same shape as the other crystals precisely owing to the anisotropy of growth rates. Quite definite faces devel­op at the expense of others more and more resolutely, and impart the shape, or habit, to the crystal that is inherent in all specimens of this substance.

Very handsome intermediate shapes are observed if a sphere is taken as a seed crystal and the solution is alternately cooled slightly and then heated slightly. Upon being heated the solution becomes undersaturated and the seed crystal is partly dissolved. Cooling leads to oversaturation of the solution and growth of the seed crystal. But the molecules attach themselves differently than they were before being dissolved, seeming to give preference to certain locations. The substance is thus transferred from certain parts of the sphere to other parts. First small circular faces appear on the surface of the sphere. These circles gradually increase in size and, contacting one another, merge along straight edges. The sphere is converted into a polyhedron. Then certain faces overtake others in their growth, a part of the faces become smaller and smaller and disappear. 
\begin{figure}[!ht]
\centering
\includegraphics[width=\textwidth]{figures/fig-04-07.pdf}
\caption{Attaining the characteristic shape of a crystal by growing from a sphere as a seed crystal.}
\label{fig-4.7}
\end{figure}
Finally, the crystal reaches its characteristic shape, or habit (\figr{fig-4.7}). In observing the growth of a crystal, one of its features may astound us -- the seemingly parallel motion of its faces. This implies that the substance separating out of the solution is deposited on the faces in layers, and that a succeeding layer is not started until the preceding one
is completed.

Illustrated in \figr{fig-4.8} is an ``incompleted'' packing arrangement of atoms. In which of the positions indicated by letters is a new atom, attaching itself to the crystal, most likely to remain secured? Without doubt, in posi­tion $K$ because here it is subject to attraction by its neigh­bours from three sides. In position $L$ the new atom is attracted from two sides, and in position $M$ from only one. This is why a row is first completed, then a whole layer and finally a new layer is started.
\begin{figure}[!ht]
\centering
\includegraphics[width=\textwidth]{figures/fig-04-08.pdf}
\caption{An ``incompleted'' packing arrangement of atoms..}
\label{fig-4.8}
\end{figure}
In many cases, crystals are formed from a molten mass, or melt. This takes place on an immense scale in nature; basalts, granites and other kinds of igneous rocks were formed of fiery magma.

Let us heat some kind of crystalline substance, for example, rock salt. Up to a temperature of \SI{804}{\celsius}, the salt crystals change very slightly; they only expand negligibly, but the substance remains a solid. A thermom­eter placed in the vessel with the substance indicates a continuous increase in temperature during heating. At \SI{804}{\celsius}, we discover two new interrelated phenomena: the substance begins to melt and the temperature ceases to increase any further. Until all of the substance is con­verted into a liquid, the temperature remains unchanged. A further increase in temperature indicates that we are heating the liquid. All crystalline substances have a def­inite melting point. Ice melts at \SI{0}{\celsius}, iron at \SI{1527}{\celsius} and mercury at \SI{39}{\celsius} below zero, etc.

As we already know, in each small crystal the atoms or molecules of the substance form an ordered packing arrange­ment and have small vibrations about their mean posi­tions. As the body is heated, the velocity of the vibrating particles increases together with the amplitude. This increase in the velocity of motion of particles as the tem­perature is raised is one of the basic laws of nature and concerns matter in any state -- solid, liquid or gaseous.

When the temperature of the crystal has reached a definite, sufficiently high value, the vibration of its particles becomes so violent that an accurate arrangement of the particles is no longer possible; the crystal melts. As melting begins, the heat supplied no longer increases the velocity of motion of the particles, it is employed to break up the crystal lattice. This is why the temperature remains constant until all of the substance melts. Subsequent heating increases the velocity of the particles of liquid.

In our case of crystallization from a melt, called freez­ing, all the phenomena described above occur in the reverse order: as the liquid is cooled, the chaotic motion of its particles slows down. When a definite, sufficiently low temperature has been reached, the velocity of the particles is so low that some of them begin to align themselves with others, forming a crystal nucleus. Until the whole substance has frozen, or solidified, the temperature remains constant. As a rule, this temperature is the same as the melting point.

Unless special measures are taken, freezing begins in the melt simultaneously at many places. The small crystals begin to grow as regular polyhedrons character­istic of the given substance in exactly the same way as we described above. Such free growth soon ends, however, because the growing small crystals run into one another, with further growth ceasing at the points of contact. This imparts a granular structure to the solidified body. Each grain is a small crystal that was unable to assume its regular shape.

Depending on many factors, and primarily on the cooling rate, a solid can consist of coarser or finer grains. The lower the cooling rate, the coarser the grains. The grain size of crystalline bodies ranges from a millionth of a centimetre to several millimetres. In the majority of cases, the granular crystalline structure can be observed under a microscope. Most solids have such a fine- crystalline structure.

The freezing of metals is a process of vital interest to engineering. Physicists have investigated in exceptionally great detail the phenomena occurring when molten metal is poured into a mould in a foundry and solidifies.

Mostly small tree-like crystals, called \emph{dendrites}, grow in a molten metal when it is cooled. Sometimes the dendrites are oriented at random, and sometimes they are parallel to one another.
\begin{figure}[!ht]
\centering
\includegraphics[width=\textwidth]{figures/fig-04-09.pdf}
\caption{Stages of growth of single dendrite.}
\label{fig-4.9}
\end{figure}
The stages of growth of a single dendrite are illustrated in \figr{fig-4.9}. Under such circumstances, the dendrite can become overgrown, occupying all the space between its branches, before encountering another dendrite in the melt. Then we find no dendrites in the solidified casting. But events can develop differently: dendrites may meet and grow into one another (with the branches of one filling the spaces between the branches of another) while they are still in an early stage of growth.

Hence, we may obtain castings whose grains (shown in \figr{fig-2.22}) have very different structures. The properties of metals depend essentially upon this structure. We can control the behaviour of a metal in solidifying by varying the cooling rate and the system of heat dis­posal from the mould.

If we wish to grow a large single crystal, we must take steps to have the crystal grow in one place. But if several crystals have already started growing, then in any case we must take steps to make the conditions for growth favourable for only one of them.

Here, for example, is what one does in growing crystals of fusible metals. The metal is melted in a glass test tube with a drawn-out bottom. The test tube, suspended on a thread inside a vertical cylindrical oven, is slowly low­ered. The drawn-out bottom gradually leaves the oven and cools off. Crystallization begins. At first, several tiny crystals are formed, but those which grow sideways come up against the test tube wall and their growth is slowed down. Only the crystal which grows along the axis of the test tube, i.e. into the heart of the melt, proves to be in favourable conditions. As the test tube is lowered, new portions of the melt coming into a low-temperature region will ``feed'' this unique crystal. Therefore, of all the tiny crystals, it alone will survive; as the test tube is lowered, it continues to grow along its axis. At last, all of the melted metal hardens in the form of a single crystal.

The same idea underlies the growing of refractory crys­tals of ruby. Fine ruby powder is poured in a stream through a flame. As a result, the particles of powder melt; tiny droplets fall on a refractory support of very small area forming a mass of crystals. 

During the further fall of droplets on the support, all the tiny crystals will grow, but once again only the one which is in the most advantageous position for ``receiving'' the falling drops will de­velop.

But what are large crystals needed for?

Industry and science are often in need of large single crystals. Of great significance for technology are crystals of Seignette salt and quartz possessing the remarkable property of transforming mechanical action (for example, pressure) into voltage.

The optical industry needs large crystals of calcite, rock salt, fluorite, etc.

Crystals of ruby, sapphire and certain other precious stones are needed for the watchmaking. The reason for this is that individual movable parts of ordinary watches perform up to \num{20000} oscillations per hour. This makes unusually heavy demands on the quality of the tips of the axes and the bearings. The wear will be least when ruby or sapphire serves as the bearing for the tip of an axis of 0.07-0.15 \si{\milli\meter} in diameter. 

Artificial crystals of these substances are very durable and are worn out very little by steel. It is remarkable that artificial stones prove to be better for this purpose than the same natural stones.

Of greatest significance in industry, however, is the growing of semiconductor (silicon) monocrystals. The communications-electronics of today is inconceivable without these crystals.

\section{Influence of Pressure on Melting Point}

If the pressure is changed, the melting point will also change. We came across such a regularity when dealing with boiling. The greater the pressure, the higher the boiling point. As a rule, this is also true for melting. However, there are a small number of substances which behave anomalously: their melting points decrease with an increase in pressure.

The fact of the matter is that the vast majority of solids are denser than their liquids. The exceptions to this rule are precisely those substances whose melting points change somewhat unusually with a change in pressure, for example, water. Ice is lighter than water, and the melting point of ice is lowered by an increase in pressure.

Compression facilitates the formation of the denser state. If the solid is denser than the liquid, then compres­sion aids solidification and hinders melting. But if melting is obstructed by compression, this means that the substance remains solid, whereas previously it would have melted at this temperature, i.e. the melting point rises with an increase in pressure. In an anomalous case, the liquid is denser than the solid, and so the pressure helps to form the liquid, i.e. lowers the melting point.

The influence of the pressure on the melting point is much less than the analogous effect for boiling. An increase in pressure of more than \SI{100}{\kgf\per\centi\meter\squared} lowers the melting point of ice by \SI{1}{\celsius}.

Why is it that skates glide on ice but do not slide on a highly polished parquet floor that is just as smooth as ice? The only feasible explanation, evidently, is the formation of water which lubricates the runners of the skates. But the pressure exerted by the runner does not in any case exceed \SI{100}{\kgf\per\centi\meter\squared}, and this would not seem to lower the melt­ing point of the ice sufficiently so that it melts under the runners. All I can say to eliminate this contradiction is the following: skates with blunt runners slide very poorly on ice. The runners must be ground so that they cut the ice. Then only the sharp edges of the runners contact the ice, exerting a pressure of tens of thousands of atmospheres and the ice \emph{does} melt.

\section{Evaporation of Solids}

When we say that a substance is evaporating, it is usually implied that a liquid is evaporating. But solids can also evaporate. The evaporation of solids is sometimes caed \emph{sublimation}.

Naphthalene, for example, is an evaporating solid. Naphthalene melts at \SI{80}{\celsius}, but evaporates at room tem­perature. It is precisely this property of naphthalene which enables it to be used for the extermination of moths. A fur coat powdered with naphthalene will become saturated with naphthalene vapours, creating an atmo­sphere which moths cannot bear. Every solid with an odour sublimes to a significant degree. In fact, the odour is created by the molecules which have broken away from the substance and reached our nose. However, the cases where a substance sublimes to an insignificant degree, sometimes to a degree that cannot be detected by even the most careful investigation, are more frequent. In prin­ciple, any solid (yes, any, even iron or copper) evaporates. If we do not detect sublimation, this merely means that the saturated vapour density is very negligible.
\begin{figure}[!ht]
\centering
\includegraphics[width=\textwidth]{figures/fig-04-10.pdf}
\caption{Saturated vapour density as a function of temperature.}
\label{fig-4.10}
\end{figure}
The saturated vapour density in equilibrium with the solid grows rapidly with temperature (\figr{fig-4.10}). It is possible to convince oneself that a number of substances having sharp odours at room temperature lose them at low temperatures.

In most cases it is impossible to considerably increase the saturated vapour density of a solid for a simple rea­son -- the substance will melt beforehand.

Ice also evaporates. This is well known to housewives who hang up their wash to dry during a frost. At first the water will freeze, but then the ice evaporates and the wash turns out to be dry.


\section{Triple Point}

Thus, there are conditions under which a vapour, a liquid and a crystal can exist pairwise in equilibrium.

Can all three states be in equilibrium? Such a point exists on the pressure-temperature diagram, it is called the \emph{triple point}. Where is it located?

If ice floating on water is placed in a closed vessel at a temperature of zero degrees, water (and ``ice'') vapour will start entering the free space. At a pressure of \SI{4.6}{\milli\meter\mercury}, evaporation will cease and saturation will set in. Now the three phases, ice, water and vapour, will be in a state of equilibrium. This is precisely the triple point.

\begin{figure}[!ht]
\centering
\includegraphics[width=\textwidth]{figures/fig-04-11.pdf}
\caption{Relationship between various states (phases) of water showing the triple point.}
\label{fig-4.11}
\end{figure}
The relationships between the various states are graph­ically and clearly shown by the diagram for water depict­ed in \figr{fig-4.11}. Such a diagram can be constructed for any substance.

We are acquainted with the curves in the diagram -- they are the equilibrium curves for ice and water vapour, ice and water, and water and water vapour. As is customary, the pressure is plotted along the vertical axis, and the temperature along the horizontal.

The three curves intersect in the triple point and divide the diagram into three regions -- the living spaces for ice, water and water vapour.

Phase diagram is a concise handbook. Its aim is to answer questions as to what state of a substance is stable at a given pressure and a given temperature.

If water or water vapour is placed under the conditions of the ``left-hand region'', it will turn into ice. If water or ice is introduced into the ``lower region'', water vapour will be obtained. In the ``right-hand region'', water va­pour will condense and ice will melt.

The phase diagram permits one to immediately say what will happen to a substance when it is heated or compressed. Heating at a constant pressure will be depicted on the diagram by a horizontal line. The point depicting the phase of the substance will move from left to right along this line.

Two such lines are depicted in the figure, one of which is heating under standard pressure. This line lies above the triple point. Therefore, it will first intersect the melting point curve, and then, beyond the bounds of the diagram, the evaporation curve too. Under standard pressure, ice melts at a temperature of \SI{0}{\celsius}, while the water so formed will boil at \SI{100}{\celsius}.

The situation will be different for ice heated under a very low pressure, say, a bit less than 5 mm Hg. The heating process is depicted by a line passing beneath the triple point. The melting and boiling point curves do not intersect this line. Under such a negligible pressure, heating leads to a direct transition of ice to water vapour.

\begin{figure}[!ht]
\centering
\includegraphics[width=0.5\textwidth]{figures/fig-04-12.pdf}
\caption{Phase digram can be used for predicting what happens at particular pressure and temperature.}
\label{fig-4.12}
\end{figure}
In \figr{fig-4.12}, this same diagram shows what an inter­esting phenomenon will occur when water vapour in the state marked on the figure by a cross is compressed. The vapour will first be converted into ice, which will then melt. The diagram enables us to say at once at what pressure the growth of crystals will begin and when the melting takes place.

The phase diagrams for all substances resemble each other. Significant, from the everyday point of view, differ­ences arise because the location of the triple point on the diagram can be most diverse for various substances.

After all, we exist under conditions which are close to ``standard'', i.e. in the first place, under a pressure close to one atmosphere. The way in which the triple point of a substance is located with respect to the standard pressure line is very important for us.
If the pressure at the triple point is less than atmospher­ic, then for us, living under ``standard'' conditions, the substance will be one of those that melt. As the temper­ature rises, it will first be transformed into a liquid, which will then boil.

In the opposite case, when the pressure at the triple point is higher than atmospheric, we will not see any liquid when the substance is heated, since the solid will be converted directly into a vapour. This is how ``dry ice'' behaves, which is very convenient for ice cream vendors.

They can distribute pieces of ``dry ice'' among the portions of ice cream without worrying that the ice cream will become wet as a result. “Dry ice” is solid carbon dioxide, \ce{CO2}. Its triple point lies at \SI{73}{\atmos}. Therefore, when solid \ce{CO2} is heated, the point depicting its state will move along a horizontal line intersecting only the evaporation curve of the solid (just as for ordinary ice at a pressure of about \SI{5}{\milli\meter\mercury}).

We have already explained how one degree of temper­ature on the Kelvin scale, or kelvin (as we are now required to call it by the SI system), is determined. The point of the discussion there was the principle involved in determining temperatures. But not all metrological institutes have ideal gas thermometers at their disposal. Therefore, a temperature scale is usually established with the aid of points of equilibrium, fixed by nature, between various states of substances.

Of especial significance in this matter is the triple point of water. A kelvin is determined today as 1/273.16 of the thermodynamic temperature of the triple point of water\footnote{In the 2019 redefinition of SI units, the kelvin was changed completely. Now the kelvin is defined in terms of the energy equivalent as given by Boltzmann's equation, , the scale remains the same numerically, i.e. the triple point of water is still the same. -- Damitr}. The triple point of oxygen is taken equal to \SI{54.361}{\kelvin}. The freezing point of gold is equal to \SI{1337.58}{\kelvin}; Using these datum points we can readily graduate any ther­mometer.

\section{The Same Atoms but Different Crystals}

The dull, soft graphite with which we write and the bright, transparent, hard, cutting glass, diamond are made of one and the same atoms -- atoms of carbon. Why then are there such differences between the properties of these two substances identical in composition?
Recall the lattice of flaky graphite each of whose atoms has three nearest neighbours, and the lattice of diamond whose atoms have four nearest neighbours. It is clearly evident from this example how the properties of crystals are determined by the mutual distribution of their atoms. Fireproof crucibles, withstanding temperatures up to two or three thousand degrees, are made of graphite, but diamond burns at temperatures above \SI{700}{\celsius}; the spe­cific gravity of diamond is 3.5, of graphite 2.3; graphite conducts electricity, but diamond does not, etc.

This property of forming various crystals is possessed not only by carbon. Almost every chemical element, and not only element but also any chemical substance, can exist in several varieties. Six varieties of ice, nine varieties of sulphur and four varieties of iron are known.

In discussing a phase diagram, we did not speak of the various types of crystals, but drew a single region for the solid. But for a great many substances, this region is divided up into sections each of which corresponds to a definite ``sort'' of solid body or, as one says, a definite solid phase (a definite crystal modification).

Each crystal phase has its own region of stable state bounded by definite intervals of pressure and temperature. The laws of transformation of one crystal variety into another are the same as the laws of melting and evapora­tion.

Given any pressure, one can find a temperature at which both types of crystals will peacefully coexist. If we raise the temperature, a crystal of one form will be converted into a crystal of the second. If we lower the temperature, the reverse transformation will take place.

In order to convert red sulphur into yellow under stan­dard pressure, a temperature below \SI{110}{\celsius} is needed. Above this temperature, right up to the melting point, the order of the distribution of atoms characteristic of red sulphur is stable. When the temperature falls, the oscillations of the atoms decrease, and, beginning with \SI{110}{\celsius}, nature finds a more convenient order for the dis­tribution of the atoms. A transformation of one crystal into another occurs.

Nobody has thought of names for the six different ices. This is what one says: ice one, ice two, \ldots, ice seven. How come seven if there are only six varieties? The rea­son is that repeated experiments have failed to detect ice four.

If water is compressed at a temperature of about zero, ice five will be formed at a pressure of about \SI{2000}{\atmos}, and ice six at a pressure of about \SI{6000}{\atmos}.

Ice two and ice three are stable at temperatures below zero degrees centigrade.

Ice seven is a hot ice; it appears when hot water is sub­jected to a pressure of about \SI{20000}{\atmos}.

All ices, except the ordinary one, are heavier than water. Ice obtained under standard conditions behaves anoma­lously; on the contrary, ice obtained under conditions differing from the norm behaves normally.

We say that each crystal modification is characterized by a definite region of existence. But if so, how can graph­ite and diamond possibly exist under identical condi­tions?
Such ``lawlessness'' is found very often in the world of crystals. The ability to live under ``alien'' conditions is almost a rule for crystals. While one must turn to various tricks in order to transfer a vapour or a liquid to alien regions of existence, a crystal, on the contrary, can almost never be forced to stay within the frontiers marked off for it by nature.

The superheating and supercooling of crystals are explained by the difficulty in transforming one order into another under conditions of extreme overcrowding. Yel­low sulphur should be transformed into red at \SI{95.5}{\celsius}. During a more or less rapid heating, we ``slip past'' this transformation point and drive it up to \SI{113}{\celsius}.

The true transformation point is most easily detected when different crystals are in contact. If we put one up close against the other and maintain a temperature of \SI{96}{\celsius}, the yellow sulphur will be eaten up by the red, but the yellow will swallow the red at \SI{95}{\celsius}. Unlike a ``crystal-liquid'' transition, a ``crystal-crystal'' transfor­mation is ordinarily delayed during superheating, just as during supercooling.

In certain cases, we come across states of a substance which are supposed to exist at entirely different temper­atures.

White tin must turn into grey when the temperature falls to \SI{+13}{\celsius}. We ordinarily use things made of white tin and know that nothing will happen to them in winter. White tin withstands supercooling of 20-30 degrees per­fectly well. However, under conditions of a severe winter, white tin is transformed into grey. The lack of knowledge of this fact was one of the circumstances destroying Scott’s expedition to the South Pole (1912). The liquid fuel taken along by the expedition was kept in vessels soldered with tin. During severe frosts, the white tin was transformed into a grey powder, the vessels came unsoldered and the fuel was spilt. Not without reason is the appear­ance of grey spots on white tin called tin plague.

Just as in the case of sulphur, white tin can be converted into grey at a temperature a bit lower than \SI{13}{\celsius}, if a tiny grain of the grey variety falls on a tin object.

The existence of several varieties of one and the same substance and the delays in their mutual transformations have great significance for technology.

At room temperature, iron atoms form a body-centred cubic lattice, occupying the vertices and centre of each cube. Every atom has 8 neighbours. At a high tempera­ture, iron atoms form a denser ``packing'' -- each atom has 12 neighbours. Iron with 8 neighbours per atom is soft; iron with 12 neighbours per atom is hard. It turns out that it is possible to obtain iron of the latter type at room temperature. The method whereby this is done, hardening, is widely applied in metallurgy.

Hardening is accomplished quite simply -- the metallic object is made red-hot and then thrown into water or oil. Cooling occurs so rapidly that there is no time for the transformation of the structure stable at a high temper­ature to take place. Consequently, the high-temperature structure will exist indefinitely under conditions unnat­ural to it; recrystallization into a stable structure occurs so slowly that it is practically unnoticeable.

In speaking of the hardening of iron, we were not quite exact. Steel, i.e. iron containing a fraction of a percent of carbon, is hardened. The presence of very small admix­tures of carbon delays the transformation of hard iron into soft and permits the hardening to be carried out. As for completely pure iron, one cannot succeed in hardening it -- there is time for the conversion of its structure to take place even during the most rapid cooling.
Depending on the form of the phase diagram, one or another transformation is achieved by changing the pres­sure or the temperature.

Many transformations of a crystal into a crystal are observed during a change of only the pressure. Black phosphorus was obtained in this way.
\begin{figure}[!ht]
\centering
\includegraphics[width=0.5\textwidth]{figures/fig-04-13.pdf}
\caption{The phase diagram of carbon.}
\label{fig-4.13}
\end{figure}
Graphite was able to be converted into diamond only by simultaneously using both high temperature and pressure. The phase diagram for carbon is depicted in \figr{fig-4.13}. At pressures below ten thousand atmospheres and at tem­peratures less than \SI{4000}{\kelvin}, graphite is the stable modi­fication. Therefore, diamond exists under “alien” con­ditions, and so it can be transformed into graphite without particular difficulty. But the inverse problem is of prac­tical interest. We cannot succeed in carrying out a transformation of graphite into diamond by only raising the pressure. A phase transformation in the solid state apparently goes too slowly. The form of the phase diagram suggests the correct solution: increase the pressure and heat the graphite simultaneously. We then obtain upper
right-hand corner of the diagram) melted carbon. Cooling it at a high pressure, we should enter the region of dia­mond.

The practical possibility of such a process was proved in 1955, and at the present time the problem is considered to be solved from technological point of view.

\section{An Amazing Liquid}

If we lower the temperature of a body, sooner or later it will solidify and acquire a crystal structure. Moreover, it makes no difference at what pressure the cooling takes place. This circumstance seems perfectly natural and understandable from the point of view of the physical laws with which we have already become acquainted. In fact, by lowering the temperature, we decrease the intensity of the thermal motion. When the motion of molecules becomes so weak that it has already ceased to interfere with the forces of interaction between them, the molecules will line up in an accurate order -- will form crystals. A further cooling will take away from the mole­cules all the energy of their motion, and at absolute zero, a substance should exist in the form of molecules at rest distributed in a regular lattice.

Experiments show that all substances behave in this manner. All but one unique substance: this ``freak'' is helium. We have already informed the reader of certain facts concerning helium. It holds the record for the value of its critical temperature. Not a single substance has its crit­ical temperature lower than \SI{4.3}{\kelvin}. However, this record by itself does not imply anything amazing. Something else is startling: cooling helium below the critical tem­perature and practically reaching absolute zero, we will not obtain solid helium. Helium remains liquid even at absolute zero.

The behaviour of helium is completely unexplainable from the point of view of the laws of motion presented by us, and is one of the signs of the limited validity of the laws of nature which seemed universal.

If a substance is liquid, its atoms are in motion. But in cooling it down to absolute zero, we have taken all the energy of motion away from it. We have to admit that helium has an energy of motion which cannot be taken-away. This conclusion is incompatible with the mechan­ics which we have been studying so far. According to the mechanics we learned, the motion of a body can always be slowed down to a complete halt by taking away all its kinetic energy; the motion of molecules can also be stopped in exactly the same way by taking energy away from them during collisions with the walls of the vessel being cooled. Such a mechanics will obviously not do for helium.

The ``strange'' behaviour of helium is an indication of a fact of enormous importance. This is the first time that we have come up against the impossibility of applying in the world of atoms the basic laws of mechanics established by means of a direct investigation of the motion of visible bodies -- laws which seemed to constitute a firm founda­tion for physics.

The fact that helium ``refuses'' to crystallize at absolute zero is by no means possible to reconcile with the mechan­ics which we have been studying until now. The contradiction which we have come across for the first time, that atoms do not obey the laws of mechanics, is merely the first link in a chain of sharper and more glaring con­tradictions in physics.

These contradictions have led to the necessity of revis­ing the foundations of the mechanics of atomic motion. This revision is very profound and has led to a change in our entire understanding of nature.
\begin{figure}[!h]
\centering
\includegraphics[width=0.5\textwidth]{figures/fig-04-14.pdf}
\caption{Phase diagram for Helium.}
\label{fig-4.14}
\end{figure}


The necessity for a radical revision of the mechanics of atomic motion does not imply that we must give up the laws of mechanics we have studied as a bad job. It would have been unfair to make the reader study unnecessary things. The old mechanics is completely valid in the world of large bodies. This is already enough for us to regard the corresponding chapters of physics with complete respect. However, it is also important that a number of laws of the ``old'' mechanics pass over without change into the ``new'' mechanics. Among them, in particular, is the law of conservation of energy.

The presence of ``inalienable'' energy at absolute zero is not a special property of helium. It turns out that all substances have ``zero'' energy. Only for helium does this energy prove sufficient to prevent the atoms from forming a regular crystal lattice.
One must not think that helium cannot occur in a crys­talline state. It is merely necessary to raise the pressure to approximately \SI{25}{\atmos} in order to crystallize helium. Cooling carried out above this pressure leads to the for­mation of solid crystalline helium with perfectly ordinary properties. Helium forms a face-centred cubic lattice.

The phase diagram for helium is shown in \figr{fig-4.14}. It differs greatly from the diagrams of all other sub­stances in the absence of a triple point. The melting and boiling point curves do not intersect.


% !TEX root = pfe-book2.tex
%!TEX TS-program = pdflatex
%!TEX encoding = UTF-8 Unicode


\cleardoublepage
%\mainmatter
\chapter{Solutions}
\label{ch-05}

\section{What a Solution Is}

If broth is salted and stirred with a spoon, no traces of salt will remain. It should not be thought that the grains of salt are simply not visible to the unaided eye. One will not succeed in detecting the tiny salt crystals in any way, because they have dissolved. If pepper is added to the broth, no solution will be obtained. One can even stir the broth for days on end, but the tiny black grains will not disappear.

But what do we mean when we say that a substance has dissolved? For the atoms or molecules composing it cannot vanish without a trace, can they? Of course not, and they do not vanish. When dissolving, only the grains of a sub­stance, crystals, accumulations of molecules of a single sort disappear. A \emph{dissolution} consists in stirring the particles a mixture in such a manner that the molecules of one substance are distributed between the molecules of another. A \emph{solution} is the mixture of molecules or atoms
of different substances.

A solution can contain various amounts of a solute.

The composition of a solution is characterized by its concentration, for example, the ratio of the number of grams of a solute to the number of litres of a solution.
As we add a solute, the concentration of a solution grows but not without bound. Sooner or later the solution will become saturated and will cease ``taking in'' the solute. The concentration of a saturated solution, i.e. the ``limiting'' concentration of the solution, is called the \emph{solubility}.

It is possible to dissolve a surprising amount of sugar in hot water. At a temperature of \SI{80}{\celsius}, a full glass of water will accept \SI{720}{\gram} of sugar with no residue. This saturated solution will be thick and viscous, and is called a sugar syrup by cooks. The figure we just gave for sugar is for a cup which holds \SI{0.2}{\liter}. Hence, the concentration of sugar in water at \SI{80}{\celsius} is equal to \SI{3600}{\gram\per\liter} (which is read ``grams per litre'').

The solubility of certain substances is highly dependent on the temperature. At room temperature (\SI{20}{\celsius}), the solubility of sugar in water falls to \SI{2000}{\gram\per\liter}. On the con­trary, the solubility of salt changes quite insignificantly with a change in temperature.

Sugar and salt dissolve well in water. But naphthalene is practically insoluble in water. Different substances dissolve quite differently in different solvents.

Solutions are used for growing monocrystals. If a small crystal of a solute is suspended in a saturated solution, then as the solvent evaporates, the solute will settle out on the surface of this crystal. Moreover, the molecules will preserve a strict order, and, as a result, the small crystal will be transformed into a large one, remaining a monocrystal.

\section{Solutions of Liquids and Gases}

Is it possible to dissolve a liquid in a liquid? Of course, it is. For example, vodka is a solution of alcohol in water (or, if you wish, of water in alcohol -- depending on what there is more of). Vodka is a real solution; the molecules of water and alcohol are completely mixed up in it.

However, such a result is not always obtained when two liquids are mixed.

Try adding kerosene to water. No matter how you stir, you will not succeed in obtaining a uniform solution; this is just as hopeless as trying to dissolve pepper in soup.

As soon as you stop stirring, the liquids arrange themselves in layers: the heavier water at the bottom, the lighter kerosene on the top. Kerosene with water and alcohol with water are systems having opposite properties with respect to their solubility.

However, there are also intermediate cases. If we mix ether with water, we shall clearly see two layers in the vessel. At first glance, it might appear that on the top is ether and at the bottom is water. But as a matter of fact, the lower and upper layers are solutions: at the bottom is water in which part of the ether is dissolved (the concentration being 25 grams of ether to a litre of water), while on the top is ether in which there is a noticeable amount of water (\SI{60}{\gram\per\liter}).

Let us now turn our attention to solutions of gases. It is clear that an unlimited amount of an arbitrary gas will dissolve in any other gas. Two gases always intermingle in such a way that the molecules of one penetrate between the molecules of the other. For gas molecules interact but slightly with each other, and so each gas behaves in the presence of another gas by, in a certain sense, not paying any ``attention'' to its cohabitant.

Gases can also dissolve in liquids. However, not in arbitrary amounts but in limited ones, not differing in this respect from solids. Moreover, different gases dis­solve differently, and these differences can be very great. Immense amounts of ammonia can dissolve in water (about 100 grams to half a glass of cold water), and large amounts of hydrogen sulphide and carbon dioxide. Oxygen and nitrogen are soluble in water in only negligible amounts (0.07 and 0.03 grams per litre of cold water).

Therefore, there is a total of only about one-hundredth of a gram of air in a litre of cold water. However, even this small amount plays a large role in life on the Earth, for fish breathe the oxygen of the air dissolved in water.

The greater the pressure of a gas, the more of it will be dissolved in a liquid. If the amount of a dissolved gas is not very great, there is a direct proportionality between it and the pressure of the gas above the surface of the liquid.

Who has not enjoyed a glass of cold soda, so good for quenching one's thirst. It is possible to obtain soda because of the dependence of the amount of a dissolved gas on the pressure. Carbon dioxide is driven into water under pressure (from the cylinders which are in every store where soda is sold). When soda is poured into a glass, the pres­sure falls to atmospheric and the water gives off the ``super­fluous'' gas in the form of bubbles.
Taking such effects into account, deep-sea divers should not rise rapidly from under the water to the surface. Addi­tional amounts of air dissolve in a diver’s blood under the high pressure of deep water. When rising, the pressure falls and air begins separating out in the form of bubbles which can stop up the blood vessels.

\section{Solid Solutions}

The word ``solution'' is applied to liquids in everyday life. However, there also exist solid mixtures whose atoms and molecules are uniformly distributed. But how can one obtain solid solutions? You won’t get them with the aid of mortar and pestle. Therefore, we must first turn the substances to be mixed into liquids, that is melt them, then mix the liquids and allow the mixture to solidify. One can also act otherwise, dissolving the two substances which we want to mix in some liquid and then evaporat­ing the solvent. Solid solutions might be obtained by such means. They might be, but they cannot ordinarily be so obtained. Solid solutions are rarities. If a piece of sugar is thrown into salty water, it will dissolve excel­lently. Evaporate the water; the minutest crystals of salt and sugar will be found at the bottom of the cup. Salt and sugar do not yield solid solutions.

It is possible to melt cadmium and bismuth in a single crucible. After cooling, we can see a mixture of cadmium and bismuth crystals through a microscope. Bismuth and cadmium do not form solid solutions either.

A necessary although not a sufficient condition for the emergence of solid solutions is the affinity of the molecules or atoms of the substances being mixed in form and dimensions. In this case, crystals of one sort are formed when the mixture freezes. The lattice points of each crystal are usually occupied randomly by atoms (molecules) of differ­ent sorts.

Alloys of metals of great technological value are fre­quently solid solutions. The dissolution of a small amount of an admixture can radically change the properties of a metal. A striking illustration of this is the obtaining of one of the most widespread materials in technology -- steel, which is a solid solution of small amounts of carbon (of the order of 0.5 of weight per cent, that is one atom of carbon to 40 atoms of iron) in iron, where the atoms of carbon are randomly distributed between the atoms of iron.

Only a small number of carbon atoms dissolve in iron. However, some solid solutions are formed by mixing substances in arbitrary proportions. Alloys of gold and copper can serve as an example. Crystals of gold and copper have lattices of the same type -- face-centred cubic. An alloy of copper with gold has the same lattice. An idea of the structure of alloys with an increasing percentage of copper can be obtained by conceptually deleting atoms of gold from the lattice and replacing them by atoms of copper. Moreover, the replacement occurs disorderly, with the copper atoms generally being distributed random­ly among the lattice points.

Alloys of copper with gold may be called solutions of replacement, whereas steel is a solution of a different type, a solution of introduction.

In the vast majority of cases, solid solutions do not arise, and, as has been said above, after cooling we can see in the microscope that the substance consists of a mixture of tiny crystals of both substances.

\section{How Solutions Freeze}

If a solution of any salt in water is cooled, it will be discovered that the freezing point has been lowered. Zero mark is past, but solidification does not occur. Only at a temperature of several degrees below zero do tiny crystals appear in the liquid. They are tiny crystals of pure ice: salt does not dissolve in solid ice.

The freezing point depends on the concentration of a solution. Increasing the concentration of a solution, we shall lower its freezing point. A saturated solution has the lowest freezing point. The decrease in the freezing point of a solution is not at all small: thus, the saturated solu­tion of sodium chloride in water freezes at \SI{-21}{\celsius}. With the aid of other salts, we can achieve an even greater fall in temperature; calcium chloride, for example, enables us to bring the freezing point of a solution down to \SI{-5}{\celsius}.

Let us now consider how the process of freezing takes place. After the first tiny crystals of ice have fallen out of a solution, the concentration of the solution increases. Now the relative number of alien molecules grows, the obstacles to the process of the crystallization of water also increase and the freezing point falls. If we do not further lower the temperature, then crystallization will cease. With a further lowering of temperature, tiny crys­tals of water (the solvent) continue separating out. Finally, the solution becomes saturated. A further enrich­ment of the solution by the solute becomes impossible, and the solution solidifies at once; moreover, if we look at the frozen mixture through a microscope, we can see that it consists of tiny crystals of ice and tiny crystals of salt.

Consequently, the freezing of a solution is unlike that of a simple liquid. The process of freezing is spread out over a wide temperature interval.

What will happen if we pour salt on some frozen surface? The answer to this question is well known to yard men -- as soon as salt comes in contact with ice, the ice begins melting. In order that this phenomenon take place, it is necessary, of course, that the freezing point of a saturated salt solution be lower than the temperature of the air. If this condition is met, the mixture of ice and salt is in an alien phase region, namely, in the region of stable exi­stence. Therefore, the mixture of ice with salt will turn into a solution, i.e. the ice will melt and the salt will dissolve in the water so formed. Finally, either all the ice will melt or else a solution of such a concentration will be formed that its freezing point will be equal to the tem­perature of the surroundings.

Suppose that a yard whose area is \SI{100}{\meter\squared} is covered by a \SI{1}{\centi\meter} layer of ice -- this is quite a bit of ice, about one ton. Let us calculate how much salt is needed to melt this ice if the temperature is \SI{-3}{\celsius}. For such a freezing (melting) point the concentration of a salt solution must be \SI{45}{\gram\per\liter}. Approximately one litre of water corresponds to one kilogram of ice. Hence, in order to melt one ton of ice at \SI{-3}{\celsius}, 45 kg of salt are needed. A much smaller amount is used in practice since one does not seek the complete melting of all the ice.

Ice melts when mixed with salt, and salt dissolves in water. But heat is required for melting, and the ice takes it from its surroundings. Therefore, the addition of salt to ice leads to a decrease in temperature.

We are now in the habit of buying ice cream, but it was previously made at home. In this connection, a mixture of ice and salt played the role of a coolant.

\section{Boiling of Solutions}

Boiling of solutions has a lot in common with their freezing.

The presence of a solute hinders crystallization. For the very same reasons, a solute also hinders boiling. In both cases it is as though alien molecules were fighting to retain a solution which is as diluted as possible. In other words, the alien molecules stabilize the state of the basic substance (i.e. facilitate its existence), which can dissolve them.

Therefore, alien molecules impede the crystallization of a liquid, and so lower the freezing point. In exactly the same way, alien molecules impede the boiling of a liquid, and so raise its boiling point.

It is curious that within certain limits of concentration (for not very concentrated solutions), neither the lowering of the freezing point of a solution nor the raising of its boiling point depends in the least on the properties of the solute but is determined only by the number of its molecules. This interesting circumstance is used for the determination of the molecular mass of soluble substances. This is done with the aid of a remarkable formula (we cannot give it here), which relates the change in the freez­ing or boiling point to the number of molecules in a unit volume of a solution (and to the heat of fusion or boiling).

The boiling point of water is raised one-third as much as its freezing point is lowered. Thus, sea water containing approximately 3.5\% of salt has a boiling point of \SI{100.6}{\celsius}, while its freezing point is lowered by \SI{2}{\celsius}.

If one liquid boils at a higher temperature than another, then (at the same temperature) its vapour pressure is lower. Hence, the vapour pressure of a solution is less than that of the pure solvent. One can judge the differ­ence on the basis of the following values: the vapour pres­sure at \SI{20}{\celsius} equals \SI{17.5}{\milli\meter\mercury} and the vapour pressure of a saturated solution of sodium chloride at the same temperature is \SI{13.2}{\milli\meter\mercury}.

Vapour with a pressure of \SI{15}{\milli\meter\mercury}, unsaturated for water, will be supersaturated for a saturated salt solu­tion. In the presence of such a solution, the vapour starts condensing and uniting with the solution. Of course, not only a salt solution but also powdered salt will take wa­ter vapour out of the air. In fact, the very first drop of water falling on the salt will dissolve it and create a saturated solution.

The absorption of water vapour from the air by salt leads to the salt’s becoming moist. This is well known to hosts and affords them grief. But this decrease in vapour pressure over a solution can also be of benefit: it is used for drying air in laboratory practice. Air is passed through calcium chloride, which is the record holder in taking moisture out of the air. While a saturated solution of sodium chloride has a vapour pressure of \SI{13.2}{\milli\meter\mercury}, that of calcium chloride is \SI{5.6}{\milli\meter\mercury}. The vapour pres­sure will fall to this value when it is passed through a sufficient amount of calcium chloride (\SI{1}{\kilo\gram} of which “has room for” approximately \SI{1}{\kilo\gram} of water). This is a negligible humidity, and such air may be regarded as dry.

\section{How Liquids Are Freed of Admixtures}

Distillation is one of the most important means of freeing liquids of admixtures. The liquid is boiled and the vapour is sent into a refrigerator. When cooled, the vapour turns into a liquid once more, but this liquid will be purer than the initial one.
It is easy to get rid of solids dissolved in a liquid with the aid of distillation. Molecules of such substances are practically absent in water vapour. Distilled water is obtained in this way -- completely tasteless pure water devoid of mineral admixtures.

However, using evaporation, we can also get rid of liquid admixtures and separate a mixture consisting of two or more liquids. In this connection, we make use of the fact that two liquids forming a mixture do not boil with the same ``ease''.

Let us see how a mixture of two liquids, for example, of water and ethyl alcohol, taken in equal amounts (100 proof vodka), will behave when boiled.

Under standard pressure, water boils at \SI{100}{\celsius} and al­cohol at \SI{78}{\celsius}. The mixture we are dealing with will boil at the intermediate temperature of \SI{81.2}{\celsius}. Alcohol boils more easily, so its vapour pressure is greater, and for an initial fifty-per cent composition of the mixture, the first portion of vapour will contain 80\% alcohol.

We can draw off the portion of vapour so obtained into a refrigerator and obtain a liquid enriched by alcohol. This process can be further repeated. However, it is clear that such a method is not practical, for with each suc­cessive distillation less and less substance will be obtained. In order that there be no such loss, so-called fraction­ating (i.e. purifying) columns are applied for the pur­pose of cleaning a liquid.

The idea behind the structure of this interesting appa­ratus consists in the following. Imagine a vertical column whose lower part contains a liquid mixture. Heat is supplied below the column and cooling is carried out above it. The vapour formed by boiling rises to the top and condenses; the resulting liquid flows down. For a fixed heat supply from below and heat removal from above, the countercurrents of vapour, going up, and liquid, flowing down, are established in the closed column.

Let us fix our attention on some horizontal cross-section of the column. The liquid passes downwards, and the vapour upwards, through this cross-section; moreover, none of the substances composing the liquid mixture stays behind. If we are dealing with a column containing a mixture of alcohol and water, the amount of alcohol passing upwards and downwards, just as the amount of water passing upwards and downwards, will be equal. Since the liquid is going down and the vapour is coming up, this means that the compositions of the liquid and of the vapour are identical at any height in the column.

As has been just explained, the equilibrium of the liquid and the vapour of a mixture of two substances requires, on the contrary, a difference between the composi­tions of these phases. Therefore, the conversion of liquid into vapour, and of vapour into liquid, takes place at any height in the column. Moreover, the high-boiling compo­nent condenses and the low-boiling component passes from the liquid to the vapour.

Consequently, it is as though the rising vapour were taking away the low-boiling component and the down­ flowing liquid were continually being enriched by the high-boiling component. The composition of the mixture will become different at each height: the higher the mixture, the greater the percentage of the low-boiling com­ponent. In the ideal case, a pure layer of the low-boiling component will be on the top and a pure layer of the high- boiling component will be at the bottom.

The substances should now be drawn off, only as slowly as possible, in order that the ideal picture just outlined not be disturbed, the low-boiling one from the top, and the high-boiling one from the bottom.

In order to carry out the separation, or rectification, in practice, we must make it possible for the counter-cur­rents of liquid and vapour to thoroughly mix. To this end, the currents of liquid and vapour are delayed with the aid of plates distributed one above the other and connected by means of overflow pipes. The liquid can flow down to a lower level from an overfilled plate. The rising vapour (0.3-1\si{\meter\per\second}) forces its way through a thin layer of liquid. The diagram of a column is shown in \figr{fig-5.1}.

\begin{figure}[!ht]
\centering
\includegraphics[width=0.5\textwidth]{figures/fig-05-01.pdf}
\caption{Separating liquids using distillation.}
\label{fig-5.1}
\end{figure}

One does not always succeed in purifying a liquid com­pletely. Certain mixtures possess an ``unpleasant'' prop­erty: for a definite composition, the ratio of the compo­nents of the evaporating molecules is the same as that in the liquid mixture. In this case, of course, a further pu­rification by means of the method described becomes im­possible. Such is the mixture containing 96\% alcohol and 4\% water: it yields a vapour of the same composition. Consequently, 96\% alcohol is the best that can be ob­tained by the evaporation method.

Rectification (or distillation) of liquids is an important process in chemical technology. For example, gasoline is obtained from oil with the aid of rectification.

It is a curious fact that rectification is the cheapest method of obtaining oxygen. For this, of course, one must change air to the liquid phase beforehand, after which one can separate it into almost pure nitrogen and oxygen by means of rectification.

\section{Purification of Solids}

On a phial containing a chemical substance alongside of the chemical name one can see, as a rule, the following notation: ``pure'', ``pure for analysis'' or ``spectroscopically pure''. They are used to mark the degree of purity of a substance: ``pure'' means a rather slight degree of purity -- there may possibly be an admixture of the order of 1\% in the substance; ``pure for analysis'' means that admixtures constitute not more than several tenths of a percent of the substance; ``spectroscopically pure'' is not easy to obtain, since spectral analysis detects an admixture of several thousandths of a percent. The inscription ``spectroscopically pure'' permits us to be confident that the purity of the substance is characterized by at least ``four nines'', i.e. that the content of the basic substance is not less than 99.99\%.

The need for pure solids is quite great. Admixtures of several thousandths of a percent can change many phys­ical properties, and in one special problem of exception­al interest to modern technology, namely the problem of obtaining semiconductors, technologists require a purity of seven nines. This means that one unnecessary atom among ten million necessary ones prevents us from solv­ing engineering problems! We take recourse in special methods for obtaining such ultrapure materials.

Ultrapure germanium and silicon (these are the main representatives of semiconductors) can be obtained by slowly drawing out a growing crystal from a melt. A rod on whose tip a seed crystal is attached is brought to the surface of melted silicon (or germanium). We then begin slowly raising the rod; the crystal coming out of the melt is made up of the atoms of the basic substance, while the atoms of the admixture remain in the melt.

The method of so-called zone refining has been applied more widely. A rod of arbitrary length and several mil­limetres in diameter is made out of the material being purified. A small cylindrical oven is placed alongside the rod enveloping it. The temperature in the oven is high enough for melting, and the part of the metal which is inside the oven melts. Therefore, a small zone of melted metal moves along the rod.

The atoms of an admixture usually dissolve much more easily in a liquid than in a solid. Consequently, at the boundary of the melted zone, the atoms of admixture pass from solid parts to the melted zone and do not pass back. It is as though the moving melted zone were dragging off the atoms of the admixture together with the melt. The oven is turned off during the reverse motion, and the operation of pulling the melted zone past the rod of metal is repeated many times. After a sufficient number of cycles, it only remains to saw off the polluted end of the rod. Ultrapure materials are obtained in a vacuum or in an atmosphere of inert gas.

When there is a large fraction of alien atoms, purifi­cation is carried out by other methods, the zone melting and pulling a crystal out of a melt being applied only for the final purification of the material.

\section{Adsorption}

Gases rarely dissolve in solids, i.e. rarely penetrate crystals. But there exists a different kind of absorption of gases by solids.  Molecules of a gas accumulate on the surface of a solid body -- this peculiar adherence is called \emph{adsorption}.\footnote{Adsorption should not be confused with absorption, which simply means taking in.} Thus, adsorption occurs when molecules cannot penetrate within a body, but successfully stick to its surface.

To be adsorbed means to be absorbed by a surface. But can such a phenomenon really play a significant role? In fact, a layer one molecule thick, even on a very large object, will weigh a negligible fraction of a gram.

Let us perform the appropriate calculations. The area of a small molecule is about \SI{10}{\angstrom\squared}, i.e. \SI{d-15}{\centi\meter\squared}. Hence, \num{d15} molecules will fit into a square centimetre. That many molecules, say, of water, do not weigh much: \SI{3d-8}{\gram}. Even in a square metre, there will be room for only \SI{0.0003}{\gram} of water.

More noticeable amounts of a substance will form on surface areas of hundreds of square metres. There is as much as \SI{0.03}{\gram} of water (\num{d21} molecules) per \SI{100}{\meter\squared}.

But do we come across such sizable surface areas in our laboratory practice? However, it is not hard to grasp that sometimes very small bodies, fitting onto the tip of a teaspoon, have enormous areas of hundreds of square metres.

A cube whose edges are \SI{1}{\centi\meter} long has a surface area of \SI{6}{\centi\meter\squared}. Let us cut such a cube into 8 equal cubes with \SI{0.5}{\centi\meter} edges. The small cubes have faces with an area of \SI{0.25}{\centi\meter\squared}. There are $6 \times 8 = 48$ such faces in all. Their total area is equal to \SI{12}{\centi\meter\squared}. The surface area has been doubled.

Thus, every splitting of a body increases its surface. Let us now break up a cube with a 1-cm side into small chips one micrometre long: $\SI{1}{\micro\meter} = \SI{d-4}{\centi\meter}$, so the large cube will be divided into \num{d12} pieces. Each piece (let us assume for the sake of simplicity that it, too, is cubic) has a surface area of \SI{12}{\micro\meter\squared}, i.e. \SI{6d-8}{\centi\meter\squared}. The total surface area of the chips is equal to \SI{6d4}{\centi\meter\squared}, i.e. \SI{6}{\meter\squared}. And such a splitting is by no means the limit.

It is quite understandable that the specific surface area (i.e. the surface area of one gram of a substance) can be immense. It grows rapidly with the crushing of a sub­stance -- for the surface area of a grain decreases in proportion to the square of its linear dimension, and the number of grains in a unit of volume increases in proportion to the cube of this dimension. A gram of water poured onto the bottom of a glass has a surface area of several square centimetres. The same gram of water in the form of raindrops has a surface area of about tens of square centimetres. But one gram of droplets of fog has a surface area of several hundred square metres.

If you break up a piece of coal (the finer the better), it will be capable of adsorbing ammonia, carbon dioxide and many toxic gases. This last property has assured for coal an application in gas-masks. Coal breaks up partic­ularly well, and the linear dimensions of its particles can be reduced to ten angstroms. Therefore, one gram of special coal has a surface area of several hundred square metres. A gas-mask with coal is capable of absorbing tens of litres of gas.

Adsorption is widely employed in the chemical indus­try. Molecules of different gases adsorbed on a surface come in close contact with each other and participate in chemical reactions more easily. In order to speed up chemical processes, one frequently makes use of finely split-up metals (nickel, copper and others) as well as of coal. Substances increasing the rate of a chemical reac­tion are called catalysts.

\section{Osmosis}

Among the living tissues, there are peculiar membranes which have the ability of letting water molecules pass through them, remaining impermeable to molecules of substances dissolved in water. The properties of these membranes are the causes of physical phenomena bearing the name ``osmotic'' (or simply ``osmosis'').

Imagine that such a semipermeable partition divides a pipe made in the form of the letter $U$ into two parts. A solution is poured into one of the elbows of the pipe, and water or some other solvent into the other elbow. Having poured the same amount of liquid into both elbows, we shall discover with surprise that there is no equilibrium when the levels are equal. After a short time, the liquids settle down at different levels. Moreover, the level is raised in the elbow containing the solution. The water separated from the solution by the semipermeable par­tition tends to dilute the solution. It is just this phenom­enon that bears the name \emph{osmosis}, and the difference in height is called the \emph{osmotic pressure}.

\begin{figure}[!ht]
\centering
\includegraphics[width=0.5\textwidth]{figures/fig-05-02.pdf}
\caption{Osmotic pressure and beverages.}
\label{fig-5.2}
\end{figure}
But what is the cause giving rise to osmotic pressure? In the right-hand elbow of the vessel (\figr{fig-5.2}), pres­sure is exerted only by water. In the left-hand elbow, the total pressure is the sum of the pressure of the water and that of the solute. But the door is open only for the water, and equilibrium in the presence of the semipermeable partition is established not when the pressure from the left equals the total pressure from the right, but when the pressure of the pure water is equal to the “water” portion of the pressure of the solution. The arising difference between the total pressures us equal to the pressure of the solute.

This excess of pressure is precisely the osmotic pres­sure. As experiments and computations show, the osmotic pressure is equal to the pressure exerted by a gas composed of the solute and occupying the same volume. It is therefore not surprising that osmotic pressure is measured in impressive numbers. The osmotic pressure arising in one litre of water when it dissolves 20 grams of sugar would balance a column of water \SI{14}{\meter} in height. 

Running the risk of arousing unpleasant memories in the reader, we shall now examine how the laxative ac­tion of certain salt solutions is related to osmotic pres­sure. The walls of the intestines are semipermeable to a number of solutions. If a salt does not pass through the walls of the intestines (this is true of Glauber’s salt), then in the intestines there arises an osmotic pressure and this pressure ``sucks'' water through the tissues from the
organism into the intestines.

Why does very salty water fail to quench one’s thirst?

It turns out that osmotic pressure is guilty of this, too. The kidneys cannot eliminate urine with the osmotic pressure greater than the pressure in the tissues of the organism. Consequently, an organism which has acquired salty sea water does not only fail to give it to the tis­sue liquids, but, on the contrary, eliminates water taken away from the tissues along with urine.

% !TEX root = pfe-book3.tex
%!TEX TS-program = pdflatex
%!TEX encoding = UTF-8 Unicode


\cleardoublepage
%\mainmatter
\chapter{Radio}
\label{ch-06}

\section{Some History}

Just as Faraday had no idea that the discovery of electromagnetic induction would lead to the founding of electrical engineering, and Ernest Rutherford considered idle talk and rank ignorance the feasibility of extracting energy from the atomic nucleus, so Heinrich Hertz, after discovering electromagnetic waves and showing how they can he detected at distances of several metres, had no notion of radio communication and even denied its possibility. Three amusing facts, are they not? But we shall leave their discussion to the psychologists.


Therefore, restricting ourselves to a mere statement of this remarkable circumstance, we shall see how events developed after Hertz's early death in 1894.
Hertz's classical experiments, which we have described in such detail, attracted the attention of scientists all over the world. Professor N. G. Yegorov of the St. Petersburg University accurately repeated these experiments. The spark in the resonator was barely visible. It could be seen only in complete darkness and then only with a magnifying glass.

Alexander Stepanovich Popov (1859-1906), an unassuming lecturer in electrical engineering at the Kronstadt Military Academy, set to work in 1889, at the age of 30, to improve Hertz's experiments. The sparks that he obtained in his resonators were much fatter than those other investigators succeeded in producing.

In 1894, a fall issue of the English journal \emph{Electrician} published an article by the well-known English physicist Sir Oliver Joseph Lodge (1851-1940) in which he claimed that Hertz's resonator could be improved by using the Branley tube. The French scientist Edouard Branley (1844-1940) was engaged in research on the conductivity of metal filings. He found that such filings do not always offer the same resistance to electric current. Loosely packed metal filings in a tube have practically infinite resistance, but if the tube is placed in the vicinity of an operating Hertz resonator, the resistance drops drastically. The explanation is that the small filings cohere owing to the welding action of the tiny sparks produced between them by the electromagnetic wave. When lightly tapped or shaken, the resistance of the filings is restored.

This property of metal filings was made use of by Lodge. He wired a circuit consisting of the Branley tube (which came to be known as a \emph{coherer}), a battery and a sensitive galvanometer. The hand of the instrument was deflected at the instant electromagnetic waves passed by. Lodge succeeded in detecting radio waves at distances up to \SI{40}{\meter}.

This system was inconvenient, however, because the coherer immediately stopped operating. It was necessary to find a way to restore the cohering (welded) filings to their initial state, with the required device designed so that the shaking occurs automatically.

%\newpage

\begin{center}
\includegraphics[width=0.8\textwidth]{figures/popov.pdf}
\end{center}
{\small \textsf{{Alexander Stepanovich Popov [1859-1906]}} -- \textsf{\footnotesize great Russian physicist and electrical engineer; invented the radio. Popov's scientific work was highly appraised by his contemporaries. In 1900 he won a Gold Medal at the World's Fair in Paris for his invention.}}


This problem was solved by Popov. He tried many kinds of coherers and finally decided to use one designed as follows (\figr{fig-6.1}). ``Two strips of thin sheet platinum, $AB$ and $CD$, are glued on the inside wall of a glass tube and extend over almost the whole length of the tube. One strip emerges at one end of the tube and the other strip at the other end. The strips are \SI{8}{\milli\meter} wide and are arranged with a space of about \SI{2}{\milli\meter} between them. The inner ends, $B$ and $C$, of the strips do not reach the corks that plug the tube so that filings jammed between the cork and the tube cannot form conducting chains that cannot be broken up when the tube is shaken. This happened in some of the earlier models. The length of the whole tube is from 6 to \SI{8}{\centi\meter} and its diameter is about \SI{1}{\centi\meter}. In operation the tube is arranged horizontally so that the strips are at the bottom and are covered by metal filings. Best operation is achieved when the metal filings fill not over one-half of the space in the tube.''


\begin{figure}[!ht]
\centering
\includegraphics[width=0.4\textwidth]{figures/fig-06-01.pdf}
\caption{Popov's coherer.}
\label{fig-6.1}
\end{figure}

Popov's coherer, just described in his own words, is shown in \figr{fig-6.1}. He used iron or steel powder in it. The main problem, however, was not to improve the
coherer, but to invent a method of restoring its initial state after detecting an electromagnetic wave. In Popov's first receiver, whose wiring diagram is shown in \figr{fig-6.2}, this job was performed by an ordinary electric bell. The striker of the bell replaced the galvanometer hand and its hammer struck the glass tube when the striker returned to its initial position.

\begin{figure}[!ht]
\centering
\includegraphics[width=\textwidth]{figures/fig-06-02.pdf}
\caption{Wiring diagram for Popov's receiver.}
\label{fig-6.2}
\end{figure}

What a simple solution of a puzzling problem! And really simple. Do you realize, dear reader, that this elementary arrangement, which neither Hertz nor Lodge hit upon, was the first application of what engineers now call a relay circuit? The negligible energy of radio waves is not directly detected, but is used to control a current circuit.

In the spring of 1895, Popov set up his apparatus outside, in an orchard. He began to move the receiver farther and farther away from the oscillator. At \SI{50}{\meter} the bell responded to the spark of the oscillator; at \SI{60}{\meter} the apparatus still worked, but at \SI{80}{\meter} the bell refused to ring. Then Popov took a coil of copper wire, threw one end up into a tree and connected the other end to the coherer. The bell rang. This is how the first antenna was devised.

In the USSR, the 7th of May is celebrated each year as Radio Day. On this day in 1895 Popov read a paper with the unassuming name ``On the Relationship of Metal Powders to Electric Oscillations'' at the regular session of the Russian Physico-Chemical Society. Many of those present had watched a demonstration, several years previously, of Hertz's experiments in which the tiny sparks could be seen only through a magnifying glass. But, when they heard the brisk clangour of Popov's receiver, they realized that they were witnessing the birth of the wireless telegraph, that a new and more efficient method of transmitting signals over long distances had been invented.

On March, 12, 1896, Popov transmitted the world's first radiogram. By pressing a key for short and long intervals, the words ``Heinrich Hertz'' were transmitted over a distance of \SI{250}{\meter} from one building to another and were recorded on telegraph tape.

By 1899, the range of radio communication between the training ships of a mine-layer detachment had already reached \SI{11}{\kilo\meter}. The practical importance of the wireless telegraph was no longer doubted even by the most skeptical officials.

The Italian electrical engineer and inventor Guglielmo Marconi (1874-1937) began his experiments somewhat later than Popov. Marconi carefully followed all the advances in the fields of electrical engineering and the study of electromagnetic waves. He skilfully employed them to improve the quality of radio reception and transmission. His especial contribution was more in the line of organization and management rather than the strictly technical aspect. This is no small matter, however, and Marconi's fame is well deserved. We should not forget that the priority in the discovery of the radio, based on the paper read on May 7, 1895, belongs to the unpretentious Russian scientist Popov who always refused to put his knowledge and research at the disposal of any country except his native land.

Marconi did not mention Popov in his articles and lectures. But not everyone knows that in 1901 he offered Prof. A. S. Popov a position in the commercial company that Marconi had founded and was the president of the range of.

The range of radio reception grew at a rapid rate. In 1899, Marconi established radio communications between France and England, and in 1901 a radio signal was sent from Europe to America.

What technical innovations facilitated these successes and made radio broadcasting feasible?

Beginning with 1899, radio engineering no longer based reception on the coherer. Instead of detecting radio waves by the drop in the resistance of a circuit due to the effect of an electromagnetic wave, an entirely different technique can be made use of. A rectified pulsating electromagnetic wave can be detected by the clicks heard in an ordinary telephone receiver.

This began a search for rectifiers. The extensively used contact detector, applied right up to the twenties of our century, consisted of a crystal with one-way conduction. Such crystals had been known since 1874. They include metal sulphides, copper pyrites and hundreds of different minerals. People of my age remember such radio sets and the irritating procedure of searching for ``good contact'' with a whisker (contact spring). Such contact was achieved when the point of the whisker had found a ``suitable'' spot on the crystal (\figr{fig-6.3}). By that time many broadcasting stations were in operation so that the set had to be tuned to the required wave. This was done with a multicontact switch for a small number of stations, or by steplessly varying the capacitance of the capacitor, which is also employed in up-to-date apparatus.

\begin{figure}[!ht]
\centering
\includegraphics[width=0.4\textwidth]{figures/fig-06-03.pdf}
\caption{Finding the whisker contact on radio.}
\label{fig-6.3}
\end{figure}

It was extremely difficult, or even impossible in some cases, to operate at high power with the spark-gap broadcasting stations because the spark-gap device became overheated. Such stations were soon replaced by one operating on the principle of an oscillating electric arc or a high-frequency alternator. After this the power ratings reached hundreds of kilowatts.

The real revolution in radio communications, enabling speech and music to be transmitted instead of only telegraph code, came with the development of the vacuum tube.

In October 1904, the English electrical engineer Sir John Ambrose Fleming (1849-1945) showed that high-frequency current could be rectified by a vacuum tube consisting of a filament heated by the current and surrounded by a metal cylinder. Its diagram is illustrated in \figr{fig-6.4}. Fleming realized the value of the vacuum-tube diode for converting electric oscillations into sound (he called this device a ``valve'' because it opened and shut the gate to the flow of electricity), but he could not achieve wide application of his detector.

\begin{figure}[!ht]
\centering
\includegraphics[width=0.4\textwidth]{figures/fig-06-04.pdf}
\caption{A vacuum-tube structure and operational setup.}
\label{fig-6.4}
\end{figure}


The fame of inventing the electron lamp was won by the American scientist Lee de Forest (1873-1961). In 1906, he converted the two-electrode tube (diode) into a triode by adding a third element. The tube was called an \emph{audion} (from the Latin word \emph{audire}, meaning \emph{I hear}. De Forest's vacuum tube radio set received signals on the grid (third element) in the lamp, rectified them and enabled them to be heard as telegraph code in headphones.

The feasibility of using the electron tube as an amplifier was evident to the American scientist. But only after seven years had passed, in 1913, a triode was first applied in a generator circuit by the German radio engineer Alexander Meissner (1883-1958).

Attempts to transmit speech, i.e. the modulation of an electromagnetic wave, had been made before the electron tube was used as a generator. But the difficulties were very great: the band of frequency modulation could not be made wide enough. Speech could be transmitted with some little success, but not music. Only in the twenties did radio transmitters and receivers, operating with electron tubes, demonstrate the truly inexhaustible capabilities of radio communications for transmitting the whole range of audio frequencies.


The next revolutionary breakthrough occurred not long ago, when semiconductor elements superseded electron tubes in radio circuits. This established a new branch of applied physics dealing with the huge complex of problems concerning the input, transmission and storage of information.

\section{Vacuum-Tube Triode and Transistor}

Vacuum-tube triodes incited an upheaval in radio engineering. But technology ages faster than people do. Today, the electronic tube has become a pensioner or, as they say in the USA, a senior citizen. Not so many years ago you could hear impatient prospective customers in shops selling TV sets demanding transistorized models, i.e. sets in which semiconductors have replaced electronic tubes.

But age demands respect. Moreover, the principles underlying the two fundamental applications of tubes and transistors, namely, the amplification and generation of waves of a definite frequency, can be more simply described by using the electronic tube as an example.


\begin{figure}[!ht]
\centering
\includegraphics[width=0.5\textwidth]{figures/fig-06-05.pdf}
\caption{Schematic showing how the grid enables plate current to be controlled.}
\label{fig-6.5}
\end{figure}

We shall therefore cover its operation in more detail than that of the transistor.
Besides the plate (anode) and heated filament (cathode), the bulb of a three-electrode tube has a sealed-in third electrode, called the grid. The electrons freely pass through the grid. Its openings are as much larger than the electron as the earth is larger than a dust speck. \figr{fig-6.5} illustrates how the grid enables the plate current to be controlled. It is evident that a negative voltage impressed on the grid reduces the plate current and a positive voltage increases this current.


\begin{figure}[!ht]
\centering
\includegraphics[width=\textwidth]{figures/fig-06-06.pdf}
\caption{Tube characteristic curve.}
\label{fig-6.6}
\end{figure}


Let us conduct a simple experiment. First we apply 100 volts across the cathode and anode (filament and plate). Then we begin to vary the grid voltage as shown, for instance, in \figr{fig-6.6} in a range from minus eight volts to plus five. Using an ammeter, we measure the current in the plate circuit. From this data we can plot the upper curve shown in \figr{fig-6.6}. This is called the \emph{tube characteristic}. Next we repeat the experiment with a plate voltage of 90 volts. We obtain a similar curve (lower curve in \figr{fig-6.6}).

Take note of the following outstanding feature. As is obvious from the hatched triangle, we can increase the plate current by 5 milliamperes in two different ways: either by increasing the plate voltage by 10 volts or by increasing the grid voltage by 2 volts. The introduction of the grid makes an amplifier out of the vacuum-tube triode. The amplification factor in our example equals 5 (ten divided by two). In other words, the effect of the grid voltage on the plate current is five times that of the plate voltage.

Now, we shall discuss how a triode enables us to generate waves of a definite length.

\begin{figure}[!ht]
\centering
\includegraphics[width=0.4\textwidth]{figures/fig-06-07.pdf}
\caption{Working of a triode.}
\label{fig-6.7}
\end{figure}

This is illustrated by the extremely simplified circuit diagram given in \figr{fig-6.7}. When the plate voltage is applied (switched on), capacitor $C_{\textrm{ckt}}$ of the oscillatory circuit is charged through the tube. The lower plate is positively charged. The capacitor is immediately discharged through the inductor $L_{\textrm{ckt}}$. This produces free (natural) oscillations that would be damped if there were no continuous energy input from the tube. How can we ensure that this energy is supplied at the proper times so that the oscillatory circuit builds up in the same way as the amplitude of a swing when you push it at the right times? This requires what is called \emph{feedback}. The current of the oscillatory circuit induces an emf in coil $L_{\textrm{fb}}$, of the same frequency as that of free oscillations. Thus the grid produces a pulsating current in the plate circuit; this current builds up the circuit with its natural frequency.



The two ingenious principles described above are the basis on which radio engineering and allied fields have been established. The electronic tube has become obsolete, making way for the transistor, but the idea of amplifying and generating electromagnetic oscillations has remained the same.

As in a vacuum-tube triode, low power in the input circuit of a transistor can control high power in the output circuit. There is a difference in the way control is accomplished. The plate current of a tube, as we have seen above, depends upon the grid voltage. The current of a collector in a transistor depends on the emitter current.

But we have not described a transistor yet. It has three electrodes. The emitter corresponds to the cathode, the collector to the plate (anode) and the base to the grid. The lead from the emitter is the input and that from the collector is the output.

\begin{figure}[!ht]
\centering
\includegraphics[width=\textwidth]{figures/fig-06-08.pdf}
\caption{The \emph{n-p-n} and \emph{p-n-p} transistors..}
\label{fig-6.8}
\end{figure}

As shown in \figr{fig-6.8}, the transistor consists of two \emph{p-n} junctions. At the left is \emph{p-n-p} transistor with the \emph{n}-layer in the middle between two \emph{p}-layers. We can also have a \emph{p}-layer in the middle, in which case we have an \emph{n-p-n} transistor (at the right).

We always bias the emitter positively so that it can produce a large number of majority carriers of charges. When the low-resistance emitter circuit changes the current in the high-resistance collector circuit, we obtain amplification.
The ways in which transistors are connected into circuits and employed as amplifiers and generators are similar, in the main, to the principles of the vacuum-tube triode. But we shall not discuss this extremely vital branch of up-to-date physics here.


\section{Radio Transmission}

The kinds of radio transmission can be classified on the basis of the power rating of the broadcasting station. Large stations transmit signals at powers ranging up to a megawatt. A miniature transmitter of the walkie-talkie type, employed by a traffic cop to notify his colleague that a Lada car with the license number 31416P has just jumped a stop light and is to be held up, radiates about one milliwatt. Even less power is sufficient for some purposes.

There are essential differences in the layout and design of stations operating with waves over several metres long and transmitters radiating ultrashort waves with a length of dozens of centimetres or even only fractions of a centimetre. Within each of the wavelength and power ranges, the engineer designing the station can apply any of a huge number of circuits and layouts that may be dictated by the locality, specific aims, economic considerations or, simply, engineering intuition.

The basic unit of a radio transmitter is the radio wave generator, or oscillator. What kind will you employ? You have at least five choices. You can use a vacuum-tube oscillator. Its range is exceptionally wide. Power ratings may vary from fractions of a watt to hundreds of kilowatts, and frequencies from dozens of kilohertzes to several gigahertzes. But if your power requirement is small, of the order of tenths of a watt, only a transistor oscillator will suit you. On the contrary, you will have to reject transistors for the time being (probably not for long) if your power requirement exceeds several hundred watts. If, however, the power rating is such that both types of oscillators can be efficiently applied, the designer will evidently prefer the transistor version. Such an engineering solution is doubtlessly more elegant.

A transistorized transmitter occupies considerably less space and, if necessary, can much more easily be designed as a portable model than a transmitter with a vacuum-tube oscillator.

Magnetron and klystron oscillators have more specialized application. The former can be extremely useful in sending pulses of several megawatts into space. The frequency band for which magnetron oscillators can be used is much narrower: it lies approximately between 300 megahertzes and 300 gigahertzes.

Klystron oscillators are used for the same range of ultrashort waves. But they find application only in low power installations: not exceeding several watts in the centimetre and several milliwatts in the millimetre bands.

The last two types of oscillator, as well as the fifth type, the quantum oscillator, are highly specific and require a special discussion. As to transistor and vacuum-tube transmitters, they resemble each other. There is a clear-cut engineering rule enabling one to replace a vacuum tube with an equivalent transistor.

The choice of the generator of electromagnetic oscillations is by far not all that is required for designing a transmitter. You must decide how to amplify the power produced by the primary (or, as they say, master) oscillator. You must also select a method for modulating the carrier wave by the audio frequencies. There are also many ways for transmitting power to the antenna field. The arrangement of the antenna field itself provides abundant opportunity for exercising engineering ingenuity.
\begin{figure}[!ht]
\centering
\includegraphics[width=\textwidth]{figures/fig-06-09.pdf}
\caption{A block diagram of a radio broadcasting station.}
\label{fig-6.9}
\end{figure}

Radio engineers frequently resort to so-called \emph{block diagrams}. Such a diagram consists of several rectangles with captions. The contents of each rectangle is cleared up as required. A block diagram of a radio broadcasting station is illustrated in \figr{fig-6.9}. The master oscillator generates continuous, almost harmonic, oscillations of the same frequency and wavelength to which you tune
your radio receiver if you wish to listen to the given station. The second unit is the power amplifier. Its name speaks for itself and we shall not describe its design. The task of the unit called the modulator is to convert sound vibrations into electric oscillations and superimpose them on the carrier wave of the broadcasting station.

Modulation can be accomplished in various ways. The simplest kind to explain is frequency modulation. In many designs the microphone is a capacitor whose capacitance is varied by the sound pressure because the capacitance depends upon the distance between the plates. Imagine now that such a capacitor is connected into the oscillatory circuit which generates the wave. Then the frequency of the wave varies with the sound pressure.

Since we have ``invaded'' the oscillatory circuit with our microphone, a band of frequencies is transmitted into space rather than a strictly definite frequency. It is sufficiently evident that ideally this spreading should include the whole audio interval of frequencies which, as we know, equals about 20 kHz.

If the station is broadcasting on long waves, corresponding to a frequency of the order of 100 kHz, the passband is about one-fifth of the carrier frequency. It is clear that long waves cannot provide for a large number of non-overlapping broadcasting stations. Short waves are an entirely different matter. For a frequency of 20 MHz, the band width is only a fraction of one per cent of the carrier frequency.

There is probably not a single home in the USSR that has no plug socket for listening to the radio. You receive these broadcasts from the so-called rebroadcasting network. It is also called wire broadcasting.

The first single-program rebroadcasting network was established in Moscow in 1925. The broadcast was transmitted simultaneously through 50 loudspeakers.


Single-program rebroadcasting is carried out on audio frequencies. From the broadcasting station, the program is transmitted by wire to the central amplifying station. From the central station it is transmitted, again by wire, to the control points, where it is again amplified and transmitted along the trunk feeder mains to the transformer substations. From each substation the wires branch off to substations of the next lower rank. Depending upon the size of the city or region, the number of steps in the network and, consequently, the number of times the voltage is stepped down, may vary. In the subscriber's lines the voltage is equal to 30 V.

Since 1962, three-wire rebroadcasting is being installed in the cities of the USSR. The transmission of two additional programs is accomplished along independent networks by amplitude modulation on carrier frequencies of 78 and 120 kHz. These two broadcasts are demodulated (i.e. the sound is separated out and the high frequency is filtered out) in the home by turning the knob of your Mayak wire-broadcasting plug-in receiver or some other model.

Thus, in three-program rebroadcasting, a single wire carries three programs simultaneously: the main program is on audio frequencies and two are non-demodulated. Therefore, the broadcasts do not interfere with one another. A simple idea, but what excellent results! Economy, reliability and high fidelity of the broadcasts are factors that indicate that wire broadcasting has a great future, including the installation of wire networks for television.

\section{Radio Reception}
Radio receiving sets are available in innumerable designs. The field of radio electronics is being developed at an exceptionally rapid rate, so that radio sets soon become obsolete and new items, better than the previous models, are available each year in the shops.

What do we mean by ``better'' with respect to receivers? Each and every reader knows the answer even if he does not understand the physics involved. A good receiving set must separate out of the chaos of radio waves that reach the antenna only those signals that are required. This property is called \emph{selectivity}. A radio set must also be as sensitive as possible, i.e. it should he able to receive even the weakest signals. Finally, it must have high fidelity, i.e. it must reproduce the music and speech broadcasted from the tuned-in station without any distortion whatsoever.

Thus, sensitivity, selectivity and fidelity. We could, perhaps, add one more requirement: the set should operate well on all wave bands.

\begin{figure}[!ht]
\centering
\includegraphics[width=\textwidth]{figures/fig-06-10.pdf}
\caption{A block diagram of a radio receiver.}
\label{fig-6.10}
\end{figure}

The block diagram of a receiver with straight amplification is sufficiently clear (\figr{fig-6.10}). It is necessary, first of all, to separate out the required wavelength and to amplify the radio-frequency oscillations produced in the antenna by the wave of the broadcasting station. Next, it is necessary to accomplish rectification, or demodulation, which is the name of the process of ``discarding'' the carrier frequency and sifting out the information carried by the sound from the electric current. Finally, it is necessary to provide one more amplifier, but for audio frequencies. The concluding stage is to convert these electric oscillations into sound. This is done by means of a dynamic loudspeaker or by headphones. The latter are used by considerate people who do not wish to be a nuisance to their neighbours.


\begin{figure}[!ht]
\centering
\includegraphics[width=0.4\textwidth]{figures/fig-06-11.pdf}
\caption{Schematic block diagram of antenna linked to inductive oscillatory circuits for tuning.}
\label{fig-6.11}
\end{figure}


The antenna of the receiving set is usually inductively linked to the oscillatory circuits of several frequency bands. When we turn the band switch knob, we perform the operation shown schematically in \figr{fig-6.11}. Within the limits of each band we usually tune the set by varying the capacitance of the oscillatory circuit capacitor. The capacity of the radio set to select the frequency in the optimal way is determined by the resonance curve of the oscillatory circuit.

I am looking at the specifications in the instruction book of an automobile radio set. Its selectivity is within 9 kHz of resonance for the long- and medium-wave frequency bands. This, of course, is not the limit that can be reached.

The sensitivity of a radio receiver is characterized by the minimum emf in its antenna that enables us to hear the broadcast with sufficient clarity (I cannot say that this definition is very exact). In the automobile radio set the sensitivity for long waves is better than \SI{175}{\micro\volt} and for ultrashort waves, better than \SI{5}{\micro\volt}.

The sensitivity depends upon the amplification factor and on set noise. The amplification factors of radio sets vary from \num{d5} to \num{d8}. This means that the broadcasting station I want to listen to must produce an induced emf of at least \SI{d-8}{\micro\volt} in the antenna of the receiver.

\section{Radio-Wave Propagation}

The simplest case is the propagation of radio waves in free space. At a relatively short distance away the radio transmitter can be regarded as a point. If so, then the radio wave front can be assumed spherical. If we imagine several spheres concentrically surrounding the transmitter, it will be clear that, in the absence of absorption, the energy passing through the spheres remains constant. As we know, the surface of a sphere is proportional to the square of the radius. Hence, the wave intensity, i.e. the amount of energy transmitted by the wave in unit time through unit area perpendicular to the direction of wave propagation, decreases as we move away from the source in inverse proportion to the square of the distance.

This important rule is applicable, of course, if no special measures have been taken to obtain a narrow directed beam of radio waves.

Various techniques exist for producing directed radio beams. One way is to use the proper beam, or antenna, array. The antennas should be arranged so that the waves they transmit are sent in the required direction, ``hump to hump''. Reflectors of various shapes are also employed for this purpose.


Radio waves travelling in space may deviate from the straight-line direction of propagation by being reflected, dispersed or refracted if they meet with an obstacle commensurate with the wavelength.

Of greatest interest is the behaviour of waves travelling near the surface of the earth. Each case may prove to be entirely unique, depending upon the wavelength.
Of cardinal importance are the electrical properties of the earth's surface and of the atmosphere. If the surface can conduct a current, it does not ``release'' the radio waves. Electric lines of force of an electromagnetic field enter a metal (or, for that matter, any conductor) at a right angle.

Now just imagine that the radio transmission is near the surface of the sea. Sea water contains dissolved salts and is therefore an electrolyte. Sea water is an excellent current conductor. Consequently, it ``holds'' the radio waves, making them travel along the surface of the sea.

Plains and timbered areas are also good conductors for currents of not especially high frequency. In other words, plains and forests behave like metals with respect to long waves.

For these reasons, long waves are contained by the whole surface of the earth and are capable of circumventing the globe. Incidentally, the velocity of radio waves can be measured in this way. Radio engineers know that it takes a radio wave 0.13 second to go around the earth. How about mountains? As far as long waves are concerned, mountains are not so high, and a wave a kilometre long can go around a mountain with more or less ease.

The feasibility of long-range reception of short waves is based on the presence of the ionosphere. The sun's rays are capable of breaking up the molecules of air in the higher regions of the atmosphere. The molecules are converted into ions and form several charged layers at altitudes from 100 to 300 km. Thus, the space through which waves of short length travel is a layer of dielectric sandwiched between two conducting surfaces.

Since plains and forest lands are poor conductors for short waves, they are incapable of holding them. Short waves depart on their journey into free space but encounter the ionosphere which reflects them like the surface of a metal.

The ionosphere is nonuniformly ionized and this varies from day to night. Therefore, short radio waves may travel along various paths. They may reach your radio set after being reflected time after time by the ionosphere and the earth. The actual fate of any short wave depends upon the angle it makes with the ionosphere layer. If this angle is close to a right angle, no reflection occurs and the wave goes through and off into outer space. More frequently, however, total internal reflection takes place and the wave returns to the earth.

The ionosphere is transparent to ultrashort waves. Radio reception of such waves is therefore possible only along a line of sight (with the receiver antenna within sight of the transmitting antenna) or with the aid of communications satellites. If we direct our wave at the satellite, we can pick up the signals reflected from it at enormous distances.

Satellites ushered in a new era in radio communications techniques by making radio and TV reception feasible on ultrashort waves.

Transmission by centimetre, millimetre and submillimetre waves presents extremely interesting opportunities. Waves of this length can be absorbed by the atmosphere. It was found, however, that there are ``windows'' and, if the wavelength is properly selected, we can make use of waves that are already within the optical band. The advantages of such waves are well known: in a narrow wave band we can ``fit in'' a huge number of non-overlapping broadcasts.

\section{Radar}
The principles of radar (\textbf{ra}dio \textbf{d}etecting \textbf{a}nd \textbf{r}anging) are sufficiently simple. We transmit a signal, it is reflected by the object that interests us and returns to our receiver. If the object is 150 m away, the signal returns in 1 microsecond, if the distance is 150 km, it returns in 1 millisecond. The direction in which the signal returns is along a line from the point where the airplane, rocket or automobile was at the instant it met the radio beam.

Naturally, the radio wave must be a pencil beam; the angle of the beam, within which the greater part of its power is concentrated, should be of the order of one degree of arc.

The principle is really not at all complicated, but the apparatus required is far from simple. To begin with, exceptionally high requirements are made to the oscillator. Vacuum-tube oscillators are used for the metre and decimetre bands (longer waves are obviously inapplicable), and klystrons and magnetrons, for the centimetre band.

The pulsed system of operation is evidently the most natural one. Very short pulses are periodically transmitted into space. The duration, or length (of time), of the pulse ranges from 0.1 to 10 microseconds in up-to-date radar transmitters. The frequency of the pulses must be selected so that the echo signal has time to return and be received in the pause before the next pulse is sent.

The maximum range at which an airplane or rocket can be detected is limited only by the condition that the object must be on a line of sight from the radar set. The reader knows, without doubt, that up-to-date radar installations are capable of picking up signals reflected from any planet of our solar system. The waves they use must, of course, pass unimpeded through the ionosphere. Fortunately, shortening the wavelength also has a direct effect on the range of radar operation because this range is proportional to the frequency of radiation and not only to the energy of the transmitted pulse.

Traces of the transmitted and received pulses are visible on the screen of an oscilloscope (cathode-ray tube). If the airplane is flying toward the observer, the trace of the echo signal moves toward the trace of the transmitted pulse.

Radar installations do not necessarily have to operate with a pulsed system. Assume that the airplane is flying toward the transmitter antenna with the velocity $v$. The radio beam is being continuously reflected from the airplane. Due to the Doppler effect, the frequency of the received waves is related to that of the transmitted waves by the equation
\begin{equation*}%
v_{r} = v_{tr} \left( 1 + \frac{2v}{c} \right)
\end{equation*}
Frequency values are determined with great precision by radio engineering methods. By tuning into resonance, we determine the value of $\nu$ and from it we calculate the airplane's velocity $v$. If, for instance, the frequency of the transmitted signal equals \SI{d9}{\hertz} and the airplane or rocket is approaching the radar antenna at a velocity of \SI{1000}{\kilo\meter\per\hour}, the received frequency exceeds the transmitted frequency by \SI{1850}{\hertz}.

The reflection of a radar beam from an airplane, rocket, steamship or automobile is not the same as from a reflector. The wavelength is commensurate with or substantially smaller than the size of the reflecting object which, moreover, is of a complex shape. As the rays are reflected by various points of the object, they (the rays) interfere with one another and are scattered to the sides. Owing to these two phenomena, the effective reflecting surface of the object differs considerably from its true surface. Calculations here are extremely complicated and only the skill and experience of the radar operator enable him to determine what kind of object is encountered by the radar beam.

You have, of course, seen radar antennas: large spherical reflectors of wire lattice structure, always in motion, surveying space. A great many different kinds of motion can be imparted to the radar reflector. For example, the beam can be moved so that it scans space in lines or circles. With this mode of operation, the path of flight of the airplane can be followed besides determining its range.
This technique is used to home aircraft into an airport under conditions of a total lack of visibility. This job may be done either by a person or even by an automatic device.

A radar set can be ``deceived''. In the first place, the object can be covered by a material that absorbs radio waves. Coal dust or rubber can be used for this purpose. Moreover, to reduce the reflection factor, the coating may be corrugated so that the major part of the radiation is disorderly scattered in all directions. If packages of aluminium foil strips or metallized fibres are thrown overboard from the airplane, the radar set is completely disoriented. This trick was first used by the English flying forces during World War II. A third way is to fill space with false radio signals.

Radar is an interesting field of engineering and has found extensive applications for peaceful purposes. It is impossible to conceive of defense today without the use of radar.

A competitor of radar is the laser (\textbf{l}ight \textbf{a}mplification by \textbf{s}timulated \textbf{e}mission of \textbf{r}adiation). The principles of location and ranging with a laser in no way differ from those described above.

Communication between spacecraft and the earth is based on the principles of radar. Radio telescopes are located so that they keep the spacecraft in view. The antennas of these telescopes are of huge size, some hundreds of metres in diameter. Such large antennas are required so that they can transmit powerful signals and pick up extremely weak signals from a radio transmitter. Of vital importance, naturally, is a narrow radio beam. If the antenna operates with a frequency of 2.2 thousand million oscillations per second (the wavelength being about 1 cm), the beam diverges only to a diameter of 1000 km over the distance to the moon. True, when the beam reaches Mars (300 million km away), its diameter is already equal to 700 000 km.

\section{Television}
Since 99 out of 100 readers daily spend an hour or two watching some TV program, it is only fair to say a few words about this wonderful invention. We shall discuss only the principles involved in television broadcasting and reception.

The idea of sending pictures over some distance consists in the following. The picture to be transmitted is divided into small squares. A physiologist can tell us how small the square must be for the eye not to be able to distinguish variations of brightness within it (the square). The luminous energy of each portion of the picture can be converted by the photoelectric effect into an electric signal. Next a way was devised for reading these signals. This is, of course, done in a strictly definite sequence as in reading a book. The electric signals modulate the electromagnetic carrier wave in exactly the same manner as in radio transmission. What happens further on is also identical to radio communications. The modulated oscillations are amplified and detected. The TV receiver reconverts the electric pulses into a visible image.

The televisor tube in the transmitter is called an image iconoscope, superorthicon or a vidicon. By means of a lens in the TV camera, the image is focussed on a photo-cathode. The most extensively used are cesium-oxide or cesium-antimony cathodes. The photocathode is mounted in an evacuated tube together with the target screen.

It is possible, in principle, to transmit an image by successively projecting the luminous flux from each element of the image. In this case, the photocurrent should flow only during the short time each element is being transmitted. Such operations would be inconvenient, however, and the camera tube contains a large number of photocells, rather than a single one, this number equalling the number of elements into which the image being transmitted is divided. This receiving screen is called the target and is of mosaic design.

The mosaic target is a thin plate of mica with a large number of grains of silver, insulated from one another and coated by cesium oxide, applied on one side. Each grain is a photocell. The other side of the mica plate is coated by a metallic film. A small capacitor is formed between each grain of the mosaic and the metal, and is charged by electrons emitted from the photocathode. It is evident that the charge of each small capacitor is proportional to the brightness of the corresponding spot on the image being transmitted.

Thus, a latent electric image of the object is produced in the metallic plate. How do we take it from the plate? By means of an electron beam which is made to scan the plate in the same way as the eye passes along the lines of print on the page of a book. The electron beam plays the role of a key switch which, for an instant, closes the electric circuit through a microcapacitor. At the instant the circuit is closed, the current is uniquely related to the brightness of the image.

Each signal can and should be amplified many times by the ordinary means applied in radio engineering. In transmitting the image, the eye should not be able to distinguish the fact that the electron beam is successively scanning the various spots of the luminous screen. A complete image, obtained on the screen of the kinescope in the TV receiver, during one cycle of electron beam motion, is called a frame. The frames must change at a frequency sufficiently high so that persistence of vision eliminates flickering of the picture being watched. 

What no-flickering frame frequency should be chosen? It must be a number related to the current frequency in the mains. The fact is that the pulsating voltage applied to the grid of the cathode-ray tube produces dark and bright lines on the screen. These lines are stationary and imperceptible only if the frequency with which the frames change is equal to or a multiple of the power mains frequency. Continuous motion is seen when the frames change with a frequency of at least \SI{20}{\hertz}; in television a frequency of \SI{25}{\hertz} is taken but a small flickering is still noticeable. It is undesirable to change the frames with a frequency of \SI{50}{\hertz} and therefore TV designers resorted to what is called \emph{interlacing} in the scanning process. The frequency is left at \SI{25}{\hertz}, but the electron beam (also called the \emph{exploring spot}) first scans the odd numbered lines and then the even numbered lines. Thus, the frequency with which the semi-frames change equals \SI{50}{\hertz} and any flickering in the brightness of the image becomes unnoticeable.

The frequency with which the frames change and the line-by-line scanning should be strictly synchronized. The technical details of this synchronization are beyond the scope of our book. Consequently, we shall not explain why the number of lines must be odd and consist of several whole factors. In the USSR the frame is divided into 625 lines, i.e, $5^{4}$. Since 25 frames are changed per second, lines are scanned with a frequency of \SI{15625}{\hertz}. From this condition follows the width of the band of TV signal frequencies.

The lowest frequency of \SI{50}{\hertz} is the semiframe frequency. The highest frequency is determined by the time required to transmit a single element.

Simple calculations, which we shall not carry out here, indicate that it is necessary to take \SI{6.5}{\mega\hertz} as the higher frequency. It follows that the carrier frequency must be at least 40 or \SI{50}{\mega\hertz}, because the frequency of the carrier wave must be at least 6 to 7 times the bandwidth of the transmitted frequencies. Now we understand why only ultrashort waves can be used in television transmission and why the range is limited to the line of sight from the transmitting antenna.

This, however, is a slip; I should have written: was limited. A breakthrough, enabling TV reception to be carried out at any distance, is the use of communications satellites. The USSR first employed satellites for this purpose. Today the whole of the Soviet Union is covered by communications accomplished by making use of a number of satellites.

Without going into the design of powerful TV stations, we shall only cite some interesting figures that characterize the huge possibilities possessed by up-to-date radio engineering in the amplification of signals. Before amplification, an ordinary video signal has a power up to \SI{d-3}{\watt}; the power amplifier increases this power by a million times. This pulse of \SI{d3}{\watt} is fed to a parabolic antenna about \SI{30}{\meter} in diameter. This antenna produces a narrow directed beam which is reflected by the satellite. After the electromagnetic wave has travelled about 35 000 km to the satellite, its power equals only \SI{d-11}{\watt}.

The amplifier mounted on the repeater satellite increases the power of this exceptionally weak signal to about \SI{10}{\watt}. When this signal, reflected from the satellite, returns to the earth, its power is \SI{d-17}{\watt}. Amplification returns the power of this video signal to its initial \SI{d-3}{\watt}.

I think that ten years ago not even the most optimistic engineer would have believed these figures.

\section{Microelectronic Circuits}
It is impossible to end this chapter devoted to radio engineering without at least a few words about the new revolution that we are witnessing in this field.
We have in mind the fantastic miniaturization of all electronic devices and apparatus. This became feasible in going over from apparatus built up of separate elements, such as resistors, transistors, etc., connected together by wiring, to electric circuits ``drawn'' by special techniques on a piece of silicon several millimetres in size.

This new technology (at least one of its versions) consists in using various kinds of masks and a series of chemical compounds enabling \emph{p}-type and \emph{n}-type impurities to be introduced at the required places of a crystal of silicon or germanium. Ion-beam treatment is used for this purpose.

An electric circuit consisting of tens of thousands of elements (!) is arranged on an area with linear dimensions of about two millimetres. When we mentioned the 
``drawing'' of the circuit, the reader may have gotten the impression that the problem involved only some surface treatment of a piece of semiconductor material. This is not so, however. The techniques are much more complicated. Each electronic element is of three-dimensional structure. Several layers containing different amounts of impurities must be built up on a tiny area of silicon.

How is this done? First a layer of oxide is applied to the surface of the silicon. This is coated with a light-sensitive material. The obtained layer-cake item is exposed to ultraviolet light through a mask of the designed shape. After developing, recesses are formed in the surface of the piece of silicon only at the places where the light passed through the mask.

The next stage consists in treating the electronic circuit being manufactured with hydrofluoric acid. The acid removes the silicon dioxide but has no effect on either the primary surface (i.e. silicon) or the light-sensitive layer. In the final stage another solvent is used to remove the light-sensitive layer. As a result, a piece of semiconductor material is coated with an insulating layer -- silicon dioxide -- wherever required by the design. The recesses of the required shape are bare silicon. These recesses are treated with an ion beam to dope the silicon with the required amount of impurities.

The manufacture of microelectronic circuits is one of the most rapidly advancing fields of engineering today. 

New ideas and new discoveries in semiconductor physics indicate that the fantastic results that have already been achieved are only the beginning.


% !TEX root = pfe-book1.tex
%!TEX TS-program = pdflatex
%!TEX encoding = UTF-8 Unicode

\chapter{Pressure}
\section{Hydraulic Press}
A hydraulic press is an ancient machine, but it has retained its significance to the present day. Take a look at \figr{fig-7-01} depicting a hydraulic press. Two pistons -- small and large -- can move in a vessel with water. If we press one piston with our hand, the pressure is transmitted to the other piston -- it will rise. Just as much water will rise above the initial position of the second piston as the first piston presses down into the vessel.
\begin{figure}[!ht]
\centering
\includegraphics[width=0.65\textwidth]{figures/fig-7-1.pdf}
\caption{A hydraulic press.}
\label{fig-7-01}
\end{figure}


If the areas of the pistons are $S_{1}$ and $S_{2}$ and their displacements are $l_{1}$ and $l_{2}$ the equality of the volumes
yields 
\begin{equation*}%
S_{1} l_{1} = S_{2} l_{2}, \quad \textrm{or} \quad \dfrac{l_{1}}{l_{2}} =\dfrac{S_{2}}{S_{1}} 
\end{equation*}
We must discover the equilibrium condition for the pistons.

We shall find such a condition without difficulty, starting out from the fact that the work performed by the balancing forces should be equal to zero. Then during the displacement of the pistons the works done by the forces exerted on them should be equal (with opposite signs). Therefore,
\begin{equation*}%
F_{1} l_{1} = F_{2} l_{2}, \quad \textrm{or} \quad \dfrac{F_{2}}{F_{1}} =\dfrac{l_{1}}{l_{2}} 
\end{equation*}
Comparing this with the preceding equality, we see that
\begin{equation*}%
\dfrac{F_{2}}{F_{1}} =\dfrac{S_{2}}{S_{1}} 
\end{equation*}
This equation implies the possibility of an enormous
multiplication of force. The piston transmitting pressure
can have an area which is hundreds or thousands of
times smaller. The force acting on the large piston will
be just as many times greater compared to the muscular
force.

With the aid of a hydraulic press, one can forge and
punch metals, press the juice out of grapes and raise
weights.

Of course, the gain in force will be accompanied by a loss in path. In order to compress a body by \SI{1}{\centi\meter} with a press, one's hand would have to cover a path as many times greater as the forces $F_{2}$ and $F_{1}$ differ.

Physicists call the ratio of the force to the area, $F/S$, the \emph{pressure} (it is denoted by the letter $p$). Instead of saying, ``One kilogram-force acts on an area of  \SI{1}{\centi\meter\squared} we shall say more concisely, ``The pressure $p = \SI{1}{\kgf\per\centi\meter\squared}$.'' This pressure is called the technical atmosphere (\SI{1}{\kgf\per\centi\meter\squared} = \SI{1}{\atmos}).

Instead of the relation 
\begin{equation*}%
\dfrac{F_{2}}{F_{1}} =\dfrac{S_{2}}{S_{1}},
\end{equation*}
one can now write:
\begin{equation*}%
\dfrac{F_{2}}{S_{2}} =\dfrac{F_{1}}{S_{1}}, \quad \textrm{i.e.} \quad p_{1}=p_{2}
\end{equation*}
Thus, the pressures on both the pistons are the same. Our reasoning does not depend on where the pistons are located or whether their surfaces are horizontal, vertical or inclined. And in general, it is not a matter of pistons. One may conceptually choose any two portions of a surface enclosing a liquid, and assert that the pressures on them are identical.

It turns out, therefore, that the pressure within a liquid is the same at all its points and in all the directions. In other words, an identical force is exerted on area elements of a definite size, irrespective of their orientation. This fact is called \emph{Pascal's law}.

\section{Hydrostatic Pressure}
Pascal's law is valid for liquids and gases. However, it fails to take into account an important circumstance -- the existence of weight.

Under terrestrial conditions, this should not be forgotten. Even water has weight. It is therefore obvious that two area elements situated at different depths under water will experience different pressures. But what will this difference be equal to? Let us conceptually single out within a liquid a right cylinder with horizontal bases. The water inside it presses on the surrounding
water. The resultant force of this pressure is equal to the weight $mg$ of the liquid in the cylinder \figr{fig-7-02}. This resultant force is made up of the forces acting on the bases of the cylinder and on its lateral surface. But the forces acting on opposite sides of the lateral surface are equal in magnitude and opposite in direction. Therefore, the sum of all the forces acting on the lateral surface
is equal to zero. Hence, the weight $mg$ will be equal to the difference in force, $F_{2} - F_{1}$. If the height of the
cylinder is $h$, the area of its base is $S$, and the density
of the liquid is $\rho$, we may write $\rho g h S$ instead of $mg$.

\begin{figure}[!ht]
\centering
\includegraphics[width=0.65\textwidth]{figures/fig-7-2.pdf}
\caption{Hydrostatic pressure on a column of water.}
\label{fig-7-02}
\end{figure}


The difference in force is equal to this quantity. In order to obtain the difference in pressure, we must divide the
weight by the area $S$. The difference in pressure turns out equal to $\rho gh$.

In accordance with Pascal's law, the pressure on differently oriented area elements located at the same depth will be identical. Hence, at two points of a liquid situated one above the other at height $h$ the difference in pressure will equal the weight of a column of the liquid whose cross-sectional area is equal to unity and whose height is $h$:
\begin{equation*}%
p_{2} - p_{1} = \rho gh
\end{equation*}
A pressure exerted by water caused by its weight is called \emph{hydrostatic}.

Under terrestrial conditions, air most often presses down on the free surface of a liquid. The pressure exerted by air is called \emph{atmospheric}. The pressure at a depth is composed of atmospheric and hydrostatic pressures.

In order to compute the force due to water pressure, it is only necessary to know the size of the area element on which it is exerted and the height of the column of liquid above it. By virtue of Pascal's law, nothing else plays any role.

\begin{figure}[!ht]
\centering
\includegraphics[width=0.85\textwidth]{figures/fig-7-3.pdf}
\caption{Are the forces at the bottom of the two vessels same?}
\label{fig-7-3}
\end{figure}


This may seem surprising. Is it possible for the forces acting on the identical bottoms of the two vessels depicted \figr{fig-7-3} to be the same? Indeed, there is much more water in the vessel on the left. In spite of this, the forces acting on the bottoms are equal to $\rho g h S$ in both cases. This is greater than the weight of the water in the vessel on the right and less than the weight of the water in the vessel on the left. The sloping walls of the vessel on the left support the weight of the ``extra'' water, but on the right, on the contrary, they add reaction forces to the weight of the water. This interesting phenomenon is sometimes called the hydrostatic paradox.

If two vessels of different form, but with water at the same level, are connected by means of a tube, water will not flow from one vessel to another. Such a flow could take place in case the pressures in the vessels were different. But this is not the case, and so the liquid in communicating vessels will always stand at one and the same level regardless of their form.

On the contrary, if the water levels in communicating vessels are different, water will begin moving and the levels will equalize.

Water pressure is much greater than air pressure. At a depth of \SI{10}{\meter}, water pressure is twice atmospheric pressure, at a depth of \SI{1}{\kilo\meter}, it is equal to \SI{100}{\atmos}.

Oceans have depths greater than \SI{10}{\kilo\meter} at certain places. The forces due to water pressure at such depths are exceptionally great. Pieces of wood which are lowered to a depth of \SI{5}{\kilo\meter} are so compressed by this enormous pressure that, after such a ``baptism'', they sink like bricks in a barrel of water.

This enormous pressure makes great difficulties for investigators of marine life. Deep-sea descents are carried out in steel globes -- the so-called bathyspheres or bathyscaphes -- which have to withstand pressures greater than \SI{1000}{\atmos}.

But submarines can dive to a depth of only \SIrange{100}{200}{\meter}.

\section{Atmospheric Pressure}

We live on the bottom of an ocean of air -- the atmosphere. Each body, every grain of sand, any object situated on the Earth is subject to air pressure.

Atmospheric pressure isn't so small. A force of about \SI{1}{\kgf} acts on each square centimetre of a body's surface. The cause of atmospheric pressure is obvious. Just as water, air possesses weight and, therefore, exerts a pressure equal (just as for water) to the weight of the column of air above the body. The higher we climb up a mountain, the less air there will be above us and, therefore, the lower will atmospheric pressure become. One must know how to measure pressure for scientific
and everyday purposes. There exist special instrument -- \emph{barometers} -- for this.

It isn't difficult to make a barometer. Mercury is poured into a tube with one end sealed off. Closing the open end with a finger, one turns the tube upside-down and submerges its open end in a cup of mercury. When this is done, the mercury in the tube will fall, but will not all pour out. The space above the mercury in the tube is undoubtedly airless. The mercury is supported in the tube by the pressure of the external air \figr{fig-7-4}.


\begin{figure}[!ht]
\centering
\includegraphics[width=0.65\textwidth]{figures/fig-7-4.pdf}
\caption{A barometer.}
\label{fig-7-4}
\end{figure}


Whatever the dimensions of the cup with mercury, and whatever the diameter of the tube, the mercury will stand at about one and the same height -- \SI{76}{\centi\meter}.

If we take a tube shorter than \SI{76}{\centi\meter}, it will be completely filled by mercury and we will not see any vacuum. A \SI{76}{\centi\meter} column of mercury presses down on the support with the same force as the atmosphere. This mercury column with a cross-sectional area of \SI{1}{\centi\meter\squared} presses down with a force of
\SI{1.033}{\kgf}. This number is the volume of the mercury $1 \times \SI{76}{\centi\meter\cubed}$ multiplied by its density and the acceleration of free fall.

As you see, the average atmospheric pressure (usually called standard atmospheric pressure) that is exerted on everything on the Earth is close to the pressure that \SI{1}{\kilo\gram} weight exerts on an area of \SI{1}{\centi\meter\squared}.

Various units are used in measuring pressures. One often simply indicates the height of a mercury column in millimetres. For example, we say that the pressure is above normal today, it is equal to \SI{768}{\milli\meter\mercury} (i.e. of mercury).

A pressure of \SI{760}{\milli\meter\mercury} is sometimes called a \emph{standard atmosphere}. A pressure of \SI{1}{\kgf\per\centi\meter\squared} is called a \emph{technical atmosphere}. Since the difference between a physical
atmosphere and a technical atmosphere is very small, from now on we will not distinguish between them.

Physicists also make frequent use of another unit of pressure, the \emph{bar}; 
\begin{equation*}%
\SI{1}{\bar} =\SI{d6}{\dyne\per\centi\meter\squared}
\end{equation*}
 Since \SI{1}{\gf} = \SI{981}{\dyne}, one bar is approximately equal to one atmosphere. More precisely, standard (normal) atmospheric pressure roughly equals 1013 millibars.
 
The unit of pressure in the SI system is the \emph{pascal} (\si{\pascal}), which is the pressure produced by a force of \SI{1}{\newton} acting on an area of \SI{1}{\meter\squared}. This is very little pressure, as can be seen from the fact that 
\begin{equation*}%
\SI{1}{\pascal} = \SI{1}{\newton\per\metre\squared} = 10\,\,\textrm{dyn}/ \si{\centi\metre\squared} = 10^{-5}\,\, \textrm{bar}.
\end{equation*}
Computing the area of the Earth's surface with the aid of the formula $4 \pi R^{2}$ , we find that the weight of the entire atmosphere is expressed by the enormous figure of \SI{5d18}{\kgf}.

Barometer tubes can have the most varied forms; only one thing is important: one of the ends of the tube must be sealed off in such a way that there be no air above the surface of the mercury. Atmospheric pressure acts on the other level of the mercury.

Atmospheric pressure can be measured by a mercury
barometer with a very great accuracy. Of course, it isn't
necessary to use only mercury; any other liquid is suitable.
But mercury is the heaviest liquid, and so the height
of a mercury column under standard pressure will be
minimum. The mercury barometer is not a particularly
convenient instrument. It is not good to leave a surface
of mercury open (mercury vapour is poisonous); further-
more, this instrument is not portable.

These drawbacks are not shared by aneroid barometers -- aneroids (i.e. airless). Everyone has seen such a barometer. It is a small round metal box with a scale and a pointer. Values of pressure are marked on the scale, usually in centimetres of a mercury column. The air has been pumped out of the metal box. The cover of the box is kept in place by a strong spring, since it would otherwise be crushed by atmospheric pressure.

\begin{figure}[!ht]
\centering
\includegraphics[width=0.65\textwidth]{figures/fig-7-5.pdf}
\caption{A siphon used to transfer liquids.}
\label{fig-7-5}
\end{figure}

With a change in atmospheric pressure, the cover either
bends or straightens. The pointer is connected to the cover
in such a manner that the pointer moves to the right
when the cover is bent.

Such a barometer is graduated by comparing its readings with those of a mercury barometer. If you want to know the pressure, don't forget to knock on the barometer with your finger. The pointer of the dial experiences considerable friction and usually gets stuck at ``yesterday's weather''.

A simple mechanism -- the siphon -- is based on atmospheric pressure,

A driver wants to help his friend, who is out of gas. But how can gasoline be poured from the tank of his car? It can't be inclined like a tea-kettle.

A rubber tube comes to his aid. He lowers one of its ends into his gas tank and orally sucks the air out of the other end. Then a rapid motion-the open end is stopped up with a finger and placed at a height below the gas tank. Now the finger can be removed -- the gasoline will pour out of the hose \figr{fig-7-5}.

A bent rubber tube is just what a siphon is. The liquid moves in this case for the same reason as through a straight inclined tube. In the final analysis, the liquid flows downwards in both cases.

Atmospheric pressure is necessary for the action of
a siphon: it ``props up'' the liquid and doesn't let the
column of liquid in the tube break. If there were no
atmospheric pressure, the column would break at the
transfer point and the liquid would slip into both vessels.

The siphon starts functioning when the liquid in the
right-hand (i.e. the ``pouring'') part of the tube drops
below the level of the liquid being siphoned off, into
which the left end of the tube has been lowered. The
liquid would otherwise flow back.

\section{How Atmospheric Pressure Was Discovered}

Suction pumps were already known to ancient civilizations. Water could be raised to a considerable height with their aid. Water very obediently followed the piston of such a pump.

Ancient philosophers thought about the causes for this and arrived at the following profound conclusion: water follows the piston because nature fears a vacuum and so does not leave any free space between the piston and the water.

It is told that an artisan constructed for the Duke of Tuscany in Florence a suction pump whose piston was supposed to draw water to a height of more than \SI{10}{\meter}. But no matter how they tried to begin sucking up water with this pump, nothing came of it. The water rose \SI{10}{\meter} with the piston, but after that the piston left the water behind, and so the very same vacuum which nature fears was formed. 

When Galileo was asked to explain the cause of this failure, he answered that nature really dislikes a vacuum, but only up to a certain point. A disciple of Galileo, Evangelista Torricelli (1608-1647), evidently used this case as an excuse to perform his famous experiment in 1643 with a tube filled with mercury. We have just described this experiment -- the constructing of a mercury barometer is precisely Torricelli's experiment.

Taking a tube of height more than \SI{76}{\centi\meter}, Torricelli created a vacuum over the mercury (it is often called a \emph{Torricellian vacuum} in his honour) and thus proved the existence of atmospheric pressure.

By means of this experiment, Torricelli cleared up the misunderstanding of the Duke of Tuscany's artisan. In fact, it is easy to see how many metres water will humbly follow the piston of a suction pump. This motion will continue until the column of water with an area of \SI{1}{\centi\meter\squared} acquires a weight of \SI{1}{\kgf}. Such a column of water will have a height of \SI{10}{\meter}. This is why nature fears a vacuum \ldots{}, but only up to \SI{10}{\meter}.

In 1654, 11 years after Torricelli's discovery, the
action of atmospheric pressure was graphically demonstrated by the Burgomaster of Magdeburg, Otto von Guericke (1602-1686). It wasn't so much the physical essence of the experiment as the theatricality of its performance that brought the author renown.

Two copper hemispheres were connected by an annular
washer. The air was pumped out of the sphere so obtained
through a pipe attached to one of the hemispheres, after
which it was impossible to separate the hemispheres.
A detailed description of Guericke's experiment has been
preserved. The atmospheric pressure on the hemispheres
can now be calculated: for a diameter of \SI{37}{\centi\meter}, the force was approximately equal to \SI{1000}{\kgf}. In order to separate
the hemispheres, Guericke ordered that two teams of
eight horses each be harnessed. Ropes passing through
the rings attached to the hemispheres were tied to the
harnesses. The horses proved unable to separate the
Magdeburg hemispheres.

The forces supplied by eight horses (exactly eight and not sixteen, since the second team harnessed for greater effect could have been replaced by a hook nailed to the wall, with no change in the force acting on the hemispheres) were not enough to break the Magdeburg hemispheres.

If there is a cavity between two bodies in contact, these bodies will not come apart because of atmospheric pressure.

\section{Atmospheric Pressure and Weather}

Pressure fluctuations caused by the weather are very
irregular. At one time people thought that pressure
alone determines the weather. Therefore, the following
inscriptions have been placed on barometers up to the
present day: clear, dry, rain, storm. You can even find
the inscription ``earthquake''.

Changes in pressure really do playa big role in changing the weather. But this role is not decisive. Average or
standard pressure at sea level is equal to 1013~millibars. Pressure fluctuations are comparatively small. The pressure rarely falls below 935-940 millibars or rises to 1055-1060. The lowest pressure -- 885 millibars -- was registered on August 18, 1927, in the South China Sea. The highest -- about 1080 millibars -- was registered on January 23, 1900, at the Barnaul station in Siberia (all figures are taken with respect to sea level).



A map used by meteorologists analyzing changes in the weather is depicted on next page \figr{fig-7-6}. The lines drawn
on the map are called isobars. The pressure is the same along each such line (its value is indicated). Note the regions of the lowest and highest pressures -- the pressure ``peaks'' and ``pockets''. The directions and strengths of winds are related to the distribution of atmospheric pressure. 

Pressures are not identical at different places on the Earth's surface, and a higher pressure ``squeezes'' air into places with a lower pressure. It would seem that a wind should blow in a direction perpendicular to the isobars, i.e. where the pressure is falling most rapidly. However, wind maps show otherwise. The Coriolis force interferes with air pressure and contributes corrections which are
very significant.


\begin{figure}[!ht]
\centering
\includegraphics[width=0.75\textwidth]{figures/fig-7-6.pdf}
\caption{A map for weather forecasting showing isobars.}
\label{fig-7-6}
\end{figure}


As we know, a Coriolis force directed to the right of the motion acts on any body moving in the Northern Hemisphere. This also pertains to air particles. ``Squeezed out'' of places of higher pressure and into places where the pressure is lower, the particle should move across the isobar, but the Coriolis force deflects it to the right, and so the direction of the wind forms an angle of about \ang{45} with the direction of the isobar.

A strikingly large effect for such a small force! This is explained by the fact that the obstacles to the action of the Coriolis force -- the friction between layers of air -- are also very insignificant.

The influence of the Coriolis force on the direction of winds at pressure ``peaks'' and ``pockets'' is even more interesting. Owing to the action of the Coriolis force, the air leaving a pressure ``peak'' does not flow in all directions along radii, but moves along curved lines -- spirals. These spiral air streams twist in one and the same direction and create a circular whirlwind displacing air masses clockwise in a high-pressure area. \figr{fig-2.16} clearly shows how a radial motion is converted into a spiral motion under the action of a constant deflecting force (see p.~\pageref{fig-2.16}).

The same thing also happens in a low-pressure area. In the absence of the Coriolis force, the air would flow towards this area uniformly along all radii. However, along the way air masses are deflected to the right. In this case, as is clear from the figure, a circular whirlwind is formed moving the air counterclockwise.

Winds in low-pressure areas are called \emph{cyclones}; winds in high-pressure areas are called \emph{anticyclones}. 

You shouldn't think that every cyclone implies a hurricane or a storm. The passing of cyclones or anticyclones through the city where we live is an ordinary phenomenon related, it is true, more often than not to a change in weather. In many cases, the approach of a cyclone means the coming of bad weather, while the approach of an anticyclone the coming of good weather. 

Incidentally, we shall not embark on the path of a weather forecaster.

\section{Change of Pressure with Altitude}

Pressure falls with an increase in altitude. This was first clarified by the Frenchman Florin P\'erier in 1648 on the instructions of Blaise Pascal. Mt. Puy de Dome, near where P\'erier lived, was \SI{975}{\meter} high. Measurements showed that the mercury in a Torricellian tube falls by \SI{8}{\milli\meter} when this mountain is climbed.


A fall in air pressure with an increase in altitude is
quite natural, for a smaller column of air then presses
down on the instrument.

If you have ever flown in an airplane, you should know
that there is an instrument on the front wall of the cabin
indicating the altitude of the airplane with an accuracy
to within tens of metres. This instrument is called an
altimeter. This is an ordinary barometer, but it has been
calibrated to show heights above sea level.

Pressure falls with an increase in altitude; let us find
a formula of this dependence. We single out a small
layer of air with an area of \SI{1}{\centi\meter\squared}  located between altitudes $h_{1}$ and $h_{2}$. The change of density with altitude is hardly noticeable within a layer which is not too large. Therefore, the weight of the volume of air we have singled
out (it is a small cylinder of height $h_{2} - h_{1}$ and base area of \SI{1}{\centi\meter\squared}) will be
\begin{equation*}%
mg = \rho \, (h_{2} - h_{1}) \, g
\end{equation*}
This weight is just what yields the fall in pressure caused by rising from altitude $h_{1}$ to altitude $h_{2}$, that is
\begin{equation*}%
\dfrac{p_{1} - p_{2}}{\rho} = g \, (h_{2} - h_{1}) 
\end{equation*}
But according to Boyle's law, which should be known to the reader (and if not, he will find it in the second book, p.~33), the density of a gas is proportional to its pressure. Consequently,
\begin{equation*}%
\dfrac{p_{1} - p_{2}}{p} \propto  (h_{2} - h_{1}) 
\end{equation*}

On the left is the fraction by which the pressure grew when the altitude was lowered from $h_{2}$ to $h_{1}$. Hence, a growth in pressure by one and the same percent will correspond to identical drops of $h_{2} - h_{1}$.

Measurements and calculations in complete agreement with each other show that the pressure will fall by 0.1 of its value for each kilometre rise above sea level. The same also holds for descents into deep shafts under sea level -- the pressure will increase by 0.1 of its value when we descend by one kilometre.

We are talking about a change of 0.1 from the value at the previous altitude. This means that during an ascent of \SI{1}{\kilo\meter}, the pressure decreases to 0.9 of the pressure at sea level; during an ascent through the next kilometre, it will become equal to 0.9 of 0.9 of the pressure at sea level; at an altitude of  \SI{3}{\kilo\meter}, the pressure will be equal to 0.9 of 0.9 of 0.9, i.e. $0.9^{3}$ , of the pressure at sea level. It is not difficult to continue this reasoning further.

Denoting the pressure at sea level by $p_{0}$, we can write out the pressure at altitude $h$ (expressed in kilometres):
\begin{equation*}%
p = p_{0} (0.87)^{h} = p_{0} \times 10^{-0.06h} 
\end{equation*}
A more precise number is written in parentheses: 0.9 is the rounded-off value. The formula presupposes the
identical temperature at all altitudes. But as a matter of fact, the temperature of the atmosphere changes with
altitude and does so, moreover, in accordance with a rather complicated law. Nevertheless, the formula
yields fairly good results and may be used for altitudes up to hundreds of kilometres.

It is not hard to determine with the aid of this formula that on the top of the Elbrus -- about \SI{5.6}{\kilo\meter} -- the pressure will fall by a factor of approximately two, while at an altitude of \SI{22}{\kilo\meter} (the record height of a stratospheric balloon's ascent with people), the pressure will fall to \SI{50}{\milli\meter\mercury}.

When we say that a pressure of  \SI{760}{\milli\meter\mercury} is standard, we must not forget to add, ``at sea level''. At an altitude
of  \SI{5.6}{\kilo\meter}, the standard pressure will not be 760, but  \SI{380}{\milli\meter\mercury}.

Along with pressure, air density also falls with an increase in altitude according to the same law. At an altitude of \SI{160}{\kilo\meter}, not much air will remain.

In fact,
\begin{equation*}%
(0.87)^{160} = 10^{-10} 
\end{equation*}

The air density at the Earth's surface is equal to about  \SI{1000}{\gram\per\centi\metre\cubed}, which means that according to our formula there should be \SI{d-7}{\gram} of air in  \SI{1}{\metre\cubed} at an altitude of \SI{160}{\kilo\meter}. But in reality, as measurements performed with the aid of rockets show, the air density at this height is ten times as great.

Our formula gives us an even greater underestimation for heights of several hundreds of kilometres. The change of temperature with altitude and also a particular phenomenon -- the decay of air molecules under the action of solar radiation -- are responsible for the fact that the formula becomes useless at great heights. Here we shall not go into this.

\section{Archimedes' Principle}

Let us hang a weight on a spring balance. The spring will stretch and show how much the weight weighs. Without taking the weight off the spring balance, let us submerge it in water. Will the reading of the spring balance change? Yes, the weight of the body seems to decrease. If the experiment is done with an iron kilogram weight, the ``loss'' in weight will constitute approximately
140~grams.

But what is the matter? For it is clear that neither mass of the weight nor its attraction by the Earth could have changed. There can be only one cause of the loss in weight: an upward force of 140~gf acts on the weight submerged in water. But where does this buoyant force discovered by the great scientist of antiquity, Archimedes, come from? Before considering a solid body in water, let
us consider ``water in water''. We conceptually single out an arbitrary volume of water. This volume possesses
weight, but does not fall to the bottom. Why? The answer is obvious -- the hydrostatic pressure of the surrounding water prevents this. This implies that the resultant of this pressure in the volume under consideration is equal to the weight of the water and directed vertically upwards.

If this volume is now occupied by a solid body, it is clear that the hydrostatic pressure will remain the same.

Thus, as a result of hydrostatic pressure, a force acts on a body immersed in a fluid. The force is directed vertically upwards and is equal in magnitude to the weight of the fluid displaced by the body. This is \emph{Archimedes' principle}.

It is said that Archimedes lay in a bath-tub and thought about how to determine whether or not there is any silver in a gold crown. A person taking a bath distinctly feels a buoyant force. Suddenly the principle came to light, presented itself to Archimedes in its remarkable simplicity. With a cry of  ``Eureka!'' (which means ``I found it!''), Archimedes jumped out of the bath-tub and ran into
the room containing the precious crown in order to immediately determine its loss of weight in water.

The loss of weight of a body in water will be equal to the weight of the water displaced by the body. Knowing the weight of the water, we shall immediately determine its volume, which is equal to the volume of the crown. Knowing the weight of the crown, we can immediately find the density of the material out of which it was made and, knowing the density of gold and silver, find the fraction of silver in the crowd.

\begin{figure}[!ht]
\centering
\includegraphics[width=0.55\textwidth]{figures/fig-7-7.pdf}
\caption{Areometer for density measurement.}
\label{fig-7-07}
\end{figure}
Archimedes' principle is valid, of course, for any fluid. If a body of volume $V$ is immersed in a fluid of density $\rho$,
then the weight of the displaced fluid -- and this is just the buoyant force -- will be equal to $\rho gV$.

The working of simple instruments controlling properties of fluid products is based on Archimedes' principle.
If alcohol or milk is diluted in water, its density will change; but it is possible to judge its composition on the
basis of its density. Such a measurement is simply and easily performed with the aid of an areometer \figr{fig-7-07}.

An areometer lowered into a liquid will be immersed to a greater or smaller depth depending on its density. An areometer will be in a state of equilibrium when the buoyant force becomes equal to the weight of the areometer.

Divisions are marked off on an areometer, and the density of a liquid is read from the marking which appears at its surface. Areometers applied for the control of alcohol are called alcoholometers, and those for the control of milk lactometers.

The average density of a person's body is somewhat greater than unity. Anyone unable to swim will drown in fresh water. Salt water has. a density greater than unity. The salinity of the water in most seas is insignificant, and its density, although greater than unity, is less than the average density of the human body. The density of the water in the Bay of Kara-Bogaz-Gol in the Caspian Sea is 1.18. This is greater than the average density of the human body. It is impossible to drown in this bay. One can lie on the water and read a book.

Ice floats on water. The preposition ``on'', incidentally, is somewhat out of place here. The density of ice is about 10\% less than that of water, so it follows from Archimedes' principle that approximately nine-tenth of a piece of ice is submerged in water. It is precisely this circumstance that makes it so dangerous for ocean liners to come across icebergs.

If a balance scale is in equilibrium in air, this does not imply that it will be in equilibrium in a vacuum. Archimedes' principle refers to air to the same degree as to water. A buoyant force equal to the weight of the displaced air acts on a body in air. A body ``weighs'' less in air than in a vacuum. The greater the volume of a body, the greater will be its loss of weight. A ton of
wood loses more weight than a ton of lead. To the humorous question of which is lighter, there is the same kind
of answer: a ton of lead is heavier than a ton of wood if
they are weighed in air.

The loss of weight in air is slight as long as we are considering small bodies. However, in weighing a piece the size of a room, we would ``lose'' several tens of kilograms. For exact weighing, the correction due to the loss of weight of large bodies in air should be taken into account.

The buoyant force in air permits us to construct balloons, aerostats and dirigibles of various types. For this
one must have a gas lighter than air.

If a balloon of volume \SI{1}{\meter\cubed}  is filled with hydrogen, \SI{1}{\meter\cubed} of which has a weight equal to \SI{0.09}{\kgf}, then the lift -- the difference between the buoyant force and the weight of the gas -- will equal
\begin{equation*}%
\SI{1.29}{\kgf} - \SI{0.09}{\kgf} = \SI{1.20}{\kgf}
\end{equation*}
 \SI{1.29}{\kilo\gram\per\centi\metre\cubed} is the density of air.

Hence, a load of about a kilogram can be attached to such a balloon, and this will not prevent it from flying
above the clouds.

It is clear that with relatively small volumes -- of several hundred cubic metres -- hydrogen balloons are capable of raising considerable loads into the air.

A serious defect of hydrogen aerostats is the inflammability of hydrogen. Together with air, hydrogen forms
an explosive mixture. Tragic accidents have marked the history of the creation of aerostats.

Therefore, when helium was discovered, people started filling balloons with it. Helium is twice as heavy as hydrogen and the lift of a balloon filled with it is smaller. But will this difference be significant? The lift of a \SI{1}{\meter\cubed} balloon filled with helium is found as the difference $\SI{1.29}{\kgf} - \SI{0.18}{\kgf} = \SI{1.11}{\kgf}$.
The lift has decreased by only 8\%. At the same time, the advantages of helium are obvious.

The aerostat was the first apparatus with whose aid people rose in the air. Aerostats with a hermetically sealed car have been used up to the present day for investigating the upper layers of the atmosphere. They are called stratospheric balloons. They rise to a height of more than \SI{20}{\kilo\meter}.
\begin{figure}[!ht]
\centering
\includegraphics[width=0.55\textwidth]{figures/fig-7-8.pdf}
\caption{A stratospheric balloon for studying upper layers of the atmosphere with various measuring instruments.}
\label{fig-7-08}
\end{figure}
Balloons equipped with various measuring devices and transmitting the results of their measurements by radio \figr{fig-7-08} are widely used at the present time. Such radiosondes contain miniature radio transmitters with batteries which report on the humidity, temperature and atmospheric pressure at various heights by means of prearranged signals.

One can send an unguided aerostat on a long journey and determine rather accurately where it will land. For this it is necessary that the aerostat climb to a great height, of the order of \SIrange{20}{30}{\kilo\meter}. Air currents are extremely stable at such heights, and the path of the aerostat can be calculated quite well beforehand. When necessary, one can automatically change the lift of the aerostat by letting out gas or throwing off ballast. 

Aerostats on which a motor with a propeller was installed were previously used for flights. Such airships were streamlined. Airships lost the competition with airplanes even in comparison with planes of 30 years ago, they are clumsy, difficult to control, move slowly and have a ``low ceiling''. It is believed that airships would be advantageous for carrying cargo.


\section{Extremely Low Pressures. Vacuum}
A vessel which is technically empty still contains an enormous number of molecules.

Molecules of gas constitute a considerable hindrance in many physical instruments. Radio tubes, $X$-ray tubes,
accelerators of elementary particles-all these instruments require a vacuum (derived from the Latin \emph{vacuus}
meaning ``empty''), i.e. space free of gas molecules. There should also be a vacuum in an ordinary electric lamp. If air enters a lamp, it will oxidize and immediately burn out.

In the best vacuum instruments, vacuum of the order of \SI{d-8}{\milli\meter} -\si{\mercury} is produced. A completely negligible pressure, it would seem: the level of mercury in a manometer would move by a hundred-millionth of a millimetre if the pressure changed by such an amount. 

However, there are still several hundred million molecules in \SI{1}{\centi\meter\cubed} at this meagre pressure.

It is interesting to compare the void of interstellar space with such a vacuum -- there one finds an average of one elementary particle of matter in several cubic centimetres.

Special pumps are employed in order to obtain vacuum. An ordinary pump removing gas by means of the motion of a piston can create a vacuum of at best \SI{0.01}{\milli\meter\mercury}. A good or, as one says, high vacuum can be obtained with the aid of a so-called diffusion (mercury or oil) pump in which gas molecules are caught up in a stream of mercury or oil vapour.

Mercury pumps, bearing the name of their inventor, Langmuir, start working only after a preliminary exhaustion to a pressure of about \SI{0.1}{\milli\meter\mercury}; such a preliminary rarefaction is called a forevacuum.

This is the way it works. A small glass container is connected to a vessel with mercury, an evacuated space and a forepump, The mercury is heated and the forepump carries away its vapour. The mercury vapour captures molecules of the gas along the way and brings them to the forepump. The mercury vapour condenses (cooling by means of running water is provided for), and the liquid
trickles down into the vessel from which the mercury began its journey.

A vacuum obtained under laboratory conditions, as we have just said, is still far from empty in the absolute sense of the word. A vacuum is greatly rarefied gas. The properties of such a gas may differ essentially from those of an ordinary gas.

The motion of the molecules ``forming a vacuum'' changes its character when the mean free path of a molecule becomes greater than the dimensions of the vessel containing the gas. The molecules then rarely collide with each other and travel in straight zigzags striking against first one and then another wall of the vessel. We shall speak in detail about the motion of molecules
in the second book. It is known to the reader that mean free path of a molecule in air at atmospheric pressure is equal to \SI{5e-6}{\centi\meter}. If we increase it by a factor of \num{d7}, it will be \SI{50}{\centi\meter}, i.e. will be noticeably greater than an average sized vessel. Since the mean free path is inversely proportional to the density, and hence also to the pressure, the pressure must be \num{d-7} of atmospheric pressure, or approximately \SI{d-4}{\milli\meter\mercury}.

Even interplanetary space is not entirely empty. But the density of the matter in it is about \SI{5e-24}{\gram\per\cubic\centi\metre}.
The main component of interplanetary matter is atomic hydrogen. At the present time, it is considered that cosmic space contains several hydrogen atoms per \SI{1}{\gram\per\centi\metre\cubed}. If a hydrogen molecule were enlarged to the size of a pea and placed in Moscow, its nearest  ``cosmic neighbour'' would prove to be in Tula.

\section{Pressures of Millions of Atmospheres}

We daily come across high pressures exerted on small surfaces. Let us estimate, for example, what the pressure will be at the point of a needle. Assume that the tip of a needle or nail has a linear dimension of \SI{0.1}{\milli\meter}. This implies that the area of the point will be about \SI{0.0001}{\centi\meter\squared}. If a rather modest force of \SI{10}{kgf} acts on such a nail, then the tip of the nail will exert a pressure of \SI{100000}{\atmos}. It's no wonder that the pointed objects so easily penetrate deeply into dense bodies.

It follows from this example that to create high pressures on small surfaces is quite a common thing. The situation is completely different if the question is to create high pressures on large surfaces.

\begin{figure}[!ht]
\centering
\includegraphics[width=0.4\textwidth]{figures/fig-7-9.pdf}
\caption{Apparatus for creating high pressures.}
\label{fig-7-9}
\end{figure}

The creation of high pressures under laboratory conditions is accomplished with the aid of powerful presses,
for example, hydraulic ones \figr{fig-7-9}. The force of  the press is transmitted to a piston of small area, and
the piston forces its way into the vessel within which we wish to create a high pressure.

Pressures of several thousand atmospheres can be created in this manner without any particular difficulty. But in order to obtain ultrahigh pressures, we must complicate the experiment, since the material composing the vessel cannot withstand such pressures.

Here nature has met us half-way. It turns out that metals become considerably stronger under pressures of the order of \SI{20000}{\atmos}. Therefore, an apparatus for obtaining ultrahigh pressures is submerged in a liquid which is under a pressure of the order of \SI{30000}{\atmos}. In this case, one is able to create pressures of several hundred thousands of atmospheres (but again with
a piston). The highest pressure -- \SI{400000}{\atmos} -- was obtained by the American physicist Percy Williams Bridgman.

Our interest in obtaining ultrahigh pressures is far from idle. Phenomena which are impossible to induce by other methods can occur at such pressures. Artificial diamonds were obtained in 1955. A pressure of \SI{100000}{\atmos} and, in addition, a temperature of \SI{2000}{\kelvin} were required for this. Ultrahigh pressures of the order of \SI{300000}{\atmos} on large surfaces are formed during explosions of solid or liquid explosive materials -- nitroglycerine, trotyl, etc. Incomparably higher pressures attaining \SI{d13}{\atmos} arise, within an atomic bomb during its explosion. Pressures during an explosion exist for a very short time.

There are constant high pressures deep inside celestial bodies including the Earth, of course. The pressure at the centre of the Earth is equal to approximately 3 million atmospheres.

\cleardoublepage
\includepdf{pfe-book1-bc.pdf}
\end{document}

% Local Variables:
% TeX-engine: pdflatex
% End:
